\chapter{The Clifford Tower and Constituent Duplication}
\label{ch:ch05_tower}

The question of this chapter is what the tower $\Cl(3n,3n)$ establishes about the number of particles. It treats the centralisers
along the tower, the Lorentz actions on several constituents, the duplication of $\Cl(3,3)$ into $\Cl(6,6)$, the second rung (rest
energies, complex structures, the Fock reading, exchange), the processes that take one constituent to two, and the ladder with its spin
parity and dual pair. The algebra is proved; reading its quanta as particles is conditional. Fock quanta are not particles,
$\Lambda^n\C^4=0$ for $n>4$ is a statement about one mode, and graded exchange does not prove a spin--statistics theorem.


% --- source: tex/s06_tower#0 ((head))


% --- source: tex/s06_tower#1 (The tower: centralisers, nesting, and carriers)
\section{The tower: centralisers, nesting, and carriers}
\label{sec:tower}

Climbing the tower adds six directions per rung and leaves the local structure of each factor unchanged. The interior four-plane
stays in the first factor; the question is what the added directions do to the internal symmetry.

\begin{theorem}[Internal centraliser]
\label{thm:centraliser}
In $\Cl(3n,3n)$ the bivectors commuting with the interior Lorentz action form
\[
\gint^{(n)}\cong\so{3n-1,3n-3},\qquad \dim\gint^{(n)}=\tfrac12(6n-4)(6n-5).
\]
\end{theorem}

\begin{proof}
By \eqref{eq:bladecomm} a bivector sharing exactly one index with an interior bivector anticommutes with it, and one sharing none
or two commutes. A bivector commuting with all six interior bivectors therefore has either both indices in $E$, which gives
$\lint$ itself and is excluded by commutation with the boosts, or both indices in the Lorentz complement
\[
W_n=\langle e^{(1)}_2,e^{(1)}_3\rangle\oplus\bigoplus_{k=2}^n\langle e^{(k)}_1,\dots,e^{(k)}_6\rangle ,
\]
which has $3n-1$ positive and $3n-3$ negative directions. All bivectors of $W_n$ commute with $\lint$, and they span $\so{W_n}$.
\end{proof}

\begin{corollary}[The photon at one constituent]
\label{cor:photon1}
At $n=1$, $W_1=U$ and $\gint^{(1)}=\R J$. The electromagnetic $\uu1$ is not a symmetry of the interior structure but the
symmetry, by exhaustion within the Lorentz-commuting grade-two connection algebra. In agreement with \cref{prop:gauge1}, no other charge
direction of that algebra gauges at $n=1$. This does not by itself select a gauging, a Maxwell kinetic term or the electromagnetic
identification, and it says nothing about symmetries of a larger endomorphism algebra.
\end{corollary}

\begin{proposition}[One frame, and nesting]
\label{prop:frame}
The interior four-plane lies in the first tensor factor, so the six interior bivectors have the same index sets at every rung and
identical structure constants. The tower is one Lorentz frame with an increasing family of modules over it. Moreover
$W_1\subset W_2\subset\cdots$, hence
\[
\gint^{(1)}\subset\gint^{(2)}\subset\gint^{(3)}\subset\cdots,\qquad \dim=1,\ 28,\ 91,\ 190,\ 325,\dots,
\]
with direct limit $\so{\infty,\infty}$. A connection valued at one rung is a connection at every higher rung with no matching
condition. The inclusion runs in the wrong direction to justify a truncation to low rungs, which would be a separate postulate.
\end{proposition}

A gauge algebra has to be compact, and every compact subalgebra of a reductive real Lie algebra is conjugate into a maximal compact
one \cite{Knapp,Helgason}, here $\so{3n-1}\oplus\so{3n-3}$. Since $\su2$ needs a faithful real representation of dimension at least
three and $\su3$ one of dimension at least six \cite{GoodmanWallach}, the rungs can be sorted by what they can carry.

\begin{theorem}[Carrier census]
\label{thm:census}
\[
\begin{array}{c|l|l|r|r|c|c|c}
\toprule
n & \gint^{(n)} & \text{maximal compact} & \dim & \text{rank} & \uu1\ \text{alone} & \su2 & \su3\\
\midrule
1 & \so{2,0} & \so2 & 1 & 1 & \text{yes} & \text{no} & \text{no}\\
2 & \so{5,3} & \so5\oplus\so3 & 13 & 3 & \text{no} & \text{yes} & \text{no}\\
3 & \so{8,6} & \so8\oplus\so6 & 43 & 7 & \text{no} & \text{yes} & \text{yes}\\
4 & \so{11,9} & \so{11}\oplus\so9 & 91 & 9 & \text{no} & \text{yes} & \text{yes}\\
\bottomrule
\end{array}
\]
At $n=2$ the projection of a simple $\su3$ to $\so3$ must vanish and $\so5$ has no proper subalgebra of dimension above six; at
$n=3$ colour enters through $\su3\subset\su4\cong\so6$.
\end{theorem}

\begin{lemma}[Rank obstruction]
\label{lem:rank}
$\rank(\su3\oplus\su2\oplus\uu1)=4$, while the maximal compact subalgebras above have ranks $1,3,7,9$. The full Standard-Model
algebra therefore embeds at no rung $n\le2$. At $n=2$ the maximal compact subalgebra has dimension $13\ge12$, so dimension counting
alone would admit it, and rank is the sharp invariant.
\end{lemma}

\begin{figure}[t]
\centering
\begin{tikzpicture}[scale=1.0]
\foreach \n/\alg/\dimc/\cmp/\cr in {1/{\so{2,0}}/{1}/{\so2}/{\uu1},2/{\so{5,3}}/{28}/{\so5\oplus\so3}/{\uu1,\ \su2},3/{\so{8,6}}/{91}/{\so8\oplus\so6}/{\uu1,\ \su2,\ \su3},4/{\so{11,9}}/{190}/{\so{11}\oplus\so9}/{\uu1,\ \su2,\ \su3}}{
  \pgfmathsetmacro{\y}{1.35*(\n-1)}
  \pgfmathsetmacro{\w}{1.6+0.6*\n}
  \fill[tcol!8] (-\w,\y) rectangle (\w,\y+1.0);
  \draw[tcol,thick] (-\w,\y) rectangle (\w,\y+1.0);
  \node at (0,\y+0.68) {\small $\gint^{(\n)}=\alg$, $\dim\dimc$};
  \node at (0,\y+0.28) {\small compact part $\cmp$};
  \node[anchor=west] at (4.3,\y+0.5) {\small $n=\n$: carries $\cr$};
}
\foreach \n in {1,2,3}{\pgfmathsetmacro{\y}{1.35*\n-0.35}\draw[->,thick,gray] (0,\y-0.0) -- (0,\y+0.35);}
\node[anchor=north] at (0,-0.15) {\small nested: $\gint^{(1)}\subset\gint^{(2)}\subset\gint^{(3)}\subset\cdots$};
\end{tikzpicture}
\caption{The tower of internal centralisers (\cref{thm:centraliser,thm:census}). Each rung contains the previous one; the compact part
decides which gauge factors can be carried. The electromagnetic $\uu1$ is the whole centraliser only at the first rung, the one rung
with no room for a nonabelian factor.}
\label{fig:tower}
\end{figure}

\begin{corollary}[The tower's trade-off]
\label{cor:tradeoff}
The $\uu1$ is the whole internal centraliser only at $n=1$, the one rung with no room for either nonabelian carrier. Inside
$\gint^{(n)}$ the centraliser of the embedded one-rung charge $e^{(1)}_2e^{(1)}_3$ has dimensions $1,16,67,154,277$ for $n=1,\dots,5$,
since by \eqref{eq:bladecomm} exactly $12(n-1)$ internal bivectors fail to commute with it. From $n=2$ onwards the photon is one
direction in a space that grows without bound; its uniqueness is lost to no dynamical process, only to the arrival of room.
\end{corollary}

At $n\ge2$ the embedded one-rung charge $e^{(1)}_2e^{(1)}_3$ must not be confused with the complex structure $J=\wI$ of
\cref{thm:i}, which commutes with every internal bivector and has all of $\gint^{(n)}$ as centraliser; they agree only at $n=1$.

\begin{remark}[The gaugeable fraction does not improve]
Of the $\binom{6n}2=3n(6n-1)$ bivectors of $\Cl(3n,3n)$, $2\binom{3n}2=3n(3n-1)$ pair directions of the same sign: the fractions
$6/15$, $30/66$, $72/153$ rise only to $1/2$. Climbing adds boosts at essentially the rate at which it adds rotations. The census
permits the chain $\uu1\subset\uu1\oplus\su2\subset\uu1\oplus\su2\oplus\su3$ at $n=1,2,3$; it permits it, and nothing selects those
subalgebras among others of the same rank.
\end{remark}

% --- source: tex/s10_actions#0 ((head))


% --- source: tex/s10_actions#1 (Other Lorentz actions, and the vector bosons of the tower)
\section{Other Lorentz actions, and the vector bosons of the tower}
\label{sec:actions}

Everything so far takes the Lorentz action from one constituent and treats the others as passengers. \Cref{thm:centraliser} is a
statement about that choice, which we call the \emph{spectator} action. The choice is adopted, not forced, and the alternatives are
different mathematics with different verdicts.

% --- source: tex/s10_actions#2 (The diagonal action)
\subsection{The diagonal action}
\label{sec:diagonal}

At three constituents let $\gamma^{(k)}_\mu$ be the interior four-plane of constituent $k$, of metric $\eta$. We call the two
remaining generators $s^{(k)}_1,s^{(k)}_2$ of that constituent its surplus pair; both square to $+\id$. Write $\mathcal S$ for the six surplus
directions and put
\[
\Sigma^{(k)}_{\mu\nu}=\tfrac12\gamma^{(k)}_\mu\gamma^{(k)}_\nu,\qquad \Delta_{\mu\nu}=\sum_{k=1}^3\Sigma^{(k)}_{\mu\nu}.
\]
Bivectors of distinct constituents commute, so $\Delta$ satisfies the $\so{1,3}$ brackets, and as a quadratic space
\begin{equation}
\label{eq:tensorpic}
\R^{9,9}=\bigl(\R^3\otimes\R^{1,3}\bigr)\oplus\R^{6,0},
\end{equation}
the first factor carrying a copy index with $\delta_{kl}$ and a Lorentz index with $\eta_{\mu\nu}$. The actions $\Delta$ and $\Sigma^{(1)}$ are
not conjugate: their centralisers have different dimensions and their Casimirs different spectra.

\begin{definition}
$M_{kl}=\tfrac12\sum_\mu\eta^{\mu\mu}\gamma^{(k)}_\mu\gamma^{(l)}_\mu$ for $k<l$ span the copy algebra $\mathfrak k_{\rm cop}$, and
$\mathfrak k_{\mathcal S}$ is the span of the fifteen surplus bivectors.
\end{definition}

\begin{theorem}[The diagonal centraliser is compact]
\label{thm:diagcent}
$\mathfrak z(\Delta)=\mathfrak k_{\mathcal S}\oplus\mathfrak k_{\rm cop}\cong\so6\oplus\so3$, of dimension $18$, and every direction of it is compact.
\end{theorem}

\begin{proof}
Under $\Delta$ the vector representation is $3V\oplus6\cdot\mathbf1$ by \eqref{eq:tensorpic}, with $V$ irreducible and nontrivial. By Schur
an equivariant endomorphism acts as $A\otimes\id_V$ on the first summand and arbitrarily on the second, with no mixing since
$\Hom_{\so{1,3}}(V,\mathbf1)=0$; antisymmetry for $\delta\otimes\eta$ makes $A\in\so3$, and the surplus block gives $\so6$. Both are
realised: the surplus bivectors commute with $\Delta$ by \eqref{eq:bladecomm}, and $M_{kl}$ is the spin lift of the copy rotation
$E_{kl}-E_{lk}$ tensored with the identity on $\R^{1,3}$. Every surplus direction squares to $+\id$, and each term of $M_{kl}$ pairs
$\gamma^{(k)}_\mu$ with $\gamma^{(l)}_\mu$, two generators of the same square, so every generator is compact by \cref{thm:sig}.
\end{proof}

For comparison, the centraliser of $\Sigma^{(1)}$ is $\so{8,6}$, of dimension $91$ with $43$ compact and $48$ boost directions; that of all
three $\Sigma^{(k)}$ together is $\mathfrak k_{\mathcal S}$ alone, of dimension $15$; and that of $\Sigma^{(1)}$ and $\Delta$ together is
$\mathfrak k_{\mathcal S}\oplus\R M_{23}$, of dimension $16$. The diagonal choice exchanges $91$ internal directions, only $43$ of them
compact, for $18$ of which all are compact. By \cref{thm:ghost} below the noncompact directions are exactly the ghost-signed ones, so the
diagonal internal algebra is free of ghosts as a whole while the spectator one is not.

% --- source: tex/s10_actions#3 (A reductive dual pair)
\subsection{A reductive dual pair}
\label{sec:dualpair}

Put $C_\Delta=-\sum_{a<b}\eta^{aa}\eta^{bb}\Delta_{ab}^2$, $C_{\rm cop}=-\sum_{k<l}M_{kl}^2$ and $C_{\mathcal S}=-\sum_{s<s'}(\tfrac12ss')^2$. For a
single constituent $-\sum_{a<b}\eta^{aa}\eta^{bb}(\Sigma^{(k)}_{ab})^2=\tfrac32$ identically, the Dirac value, and $C_{\mathcal S}=\tfrac{15}4$
identically, so the surplus $\so6$ acts isotypically on $S_3$ and its Casimir separates nothing.

\begin{theorem}[Spectrum of the diagonal Casimir]
\label{thm:diagCspec}
$\spec C_\Delta=\{\tfrac32\,(160),\ \tfrac{11}2\,(288),\ \tfrac{15}2\,(64)\}$, each eigenprojector being Lorentz invariant. The
$C_\Delta=\tfrac32$ block, the \emph{Dirac block} $\mathcal D$, has rank $160$.
\end{theorem}

\begin{proof}
A quadratic Casimir of $\slr{2,\C}\oplus\slr{2,\C}$ is $\alpha j_L(j_L+1)+\beta j_R(j_R+1)$, and the Dirac value $\tfrac32$ together with the
vector value $3$ gives $\alpha=\beta=2$. Since the constituent Casimir is $\tfrac32$ identically, each $\R^8$ is $M\oplus M$ with $M$ the
real Majorana module, and $S_3\cong M^{\otimes3}\otimes\R^8$. Writing $A=(\tfrac12,0)$ and $B=(0,\tfrac12)$,
\[
\begin{aligned}
AAA&=(\tfrac32,0)\oplus2(\tfrac12,0), & AAB&=(1,\tfrac12)\oplus(0,\tfrac12),\\
ABB&=(\tfrac12,1)\oplus(\tfrac12,0), & BBB&=(0,\tfrac32)\oplus2(0,\tfrac12),
\end{aligned}
\]
with three words of each mixed type. This gives multiplicities $1,3,5$ for $(\tfrac32,0)\oplus(0,\tfrac32)$, $(1,\tfrac12)\oplus(\tfrac12,1)$ and
$(\tfrac12,0)\oplus(0,\tfrac12)$, of complex dimensions $8,12,4$; the totals $8,36,20$ add up to $64$, and the multiplicity factor $8$ turns them
into $64,288,160$.
\end{proof}

\begin{theorem}[A Casimir identity]
\label{thm:casid}
As operators on $S_3$,
\begin{equation}
\label{eq:casid}
C_{\rm cop}+C_\Delta=\tfrac{15}2\,\id ,
\end{equation}
and for $n$ constituents with a $d$-dimensional interior, $C_{\rm cop}+C_\Delta=\bigl[nd(d-1)/8+\binom n2d/4\bigr]\id$; for $d=4$ this is
$\tfrac12n(n+2)\,\id$.
\end{theorem}

\begin{proof}
Write $B^{kl}_{ab}=\gamma^{(k)}_a\gamma^{(k)}_b\gamma^{(l)}_a\gamma^{(l)}_b$, symmetric in $k\leftrightarrow l$ and in $a\leftrightarrow b$. Expanding,
\[
\Delta_{ab}^2=\tfrac14\Bigl[-n\,\eta_{aa}\eta_{bb}+2\sum_{k<l}B^{kl}_{ab}\Bigr],\qquad
M_{kl}^2=-\tfrac d4-\tfrac12\sum_{\mu<\nu}\eta^{\mu\mu}\eta^{\nu\nu}B^{kl}_{\mu\nu},
\]
the first from $\gamma_a\gamma_b\gamma_a\gamma_b=-\eta_{aa}\eta_{bb}$, the second from
$\gamma^{(k)}_\mu\gamma^{(l)}_\mu\gamma^{(k)}_\nu\gamma^{(l)}_\nu=-B^{kl}_{\mu\nu}$ for $\mu\neq\nu$. Hence $C_\Delta=nd(d-1)/8-\tfrac12\mathsf S$ and
$C_{\rm cop}=\binom n2d/4+\tfrac12\mathsf S$ with $\mathsf S=\sum_{k<l}\sum_{a<b}\eta^{aa}\eta^{bb}B^{kl}_{ab}$, the only term in which two
constituents interact. It cancels in the sum.
\end{proof}

At $n=3$, $d=4$ the constant is $3\cdot4\cdot3/8+3\cdot4/4=\tfrac92+3=\tfrac{15}2$. The identity is inherited: its general form $nd(d+n-2)/8$ is the
case $d_{\rm N}=n$, $k_{\rm N}=d/2$ of the relation $C_{\mathfrak o(d_{\rm N})}+C_{\mathfrak o(2k_{\rm N})}=\tfrac14k_{\rm N}d_{\rm N}(d_{\rm N}+2k_{\rm N}-2)$ between the
Casimirs of the $\mathfrak o(d_{\rm N})$--$\mathfrak o(2k_{\rm N})$ dual pair on a fermionic Fock space of $k_{\rm N}$ species \cite[eq.~(73)]{Neergard19}, an
integer $k_{\rm N}$ ($2$ for $d=4$, $3$ for $d=6$) at every rung, odd ones included (see also \cite{Neergard20,Howe89}). The proof given here is an independent derivation in the
real Clifford realisation.

\begin{corollary}[A dual pair]
\label{cor:howe}
$\so3_{\rm cop}$ and $\so{1,3}_\Delta$ are mutual centralisers inside $\so{3,9}$, and the joint spectrum of $(C_{\rm cop},C_\Delta)$ has exactly
three points:
\begin{center}
\begin{tabular}{@{}cccccc@{}}
\toprule
$j_{\rm cop}$ & $\dim$ & $C_{\rm cop}$ & $C_\Delta$ & Lorentz content & rank in $S_3$\\
\midrule
$0$ & $1$ & $0$ & $\tfrac{15}2$ & $(\tfrac32,0)\oplus(0,\tfrac32)$ & $8\cdot1\cdot8=64$\\
$1$ & $3$ & $2$ & $\tfrac{11}2$ & $(1,\tfrac12)\oplus(\tfrac12,1)$ & $12\cdot3\cdot8=288$\\
$2$ & $5$ & $6$ & $\tfrac32$ & $(\tfrac12,0)\oplus(0,\tfrac12)$ & $4\cdot5\cdot8=160$\\
\bottomrule
\end{tabular}
\end{center}
The multiplicities $1,3,5$ of \cref{thm:diagCspec} are the copy representations, so $\mathcal D\cong L\otimes U$ with
$L=(\tfrac12,0)\oplus(0,\tfrac12)$ and $U=[j_{\rm cop}=2]\otimes\R^8$ carrying the whole gauge action.
\end{corollary}

\begin{figure}[t]
\centering
\begin{tikzpicture}
\begin{axis}[width=8.6cm,height=6.2cm,xmin=-0.6,xmax=7.6,ymin=0,ymax=8.6,
  xlabel={$C_{\rm cop}$},ylabel={$C_\Delta$},xtick={0,2,6},ytick={1.5,5.5,7.5},
  yticklabels={$\tfrac32$,$\tfrac{11}2$,$\tfrac{15}2$},axis lines=left,clip=false]
\addplot[gray,dashed,domain=-0.3:7.5] {7.5-x};
\addplot[only marks,mark=*,mark size=3pt,tcol] coordinates {(0,7.5) (2,5.5) (6,1.5)};
\node[anchor=west] at (axis cs:0.25,7.75) {\small $j_{\rm cop}=0$: $(\tfrac32,0)\oplus(0,\tfrac32)$, rank $64$};
\node[anchor=west] at (axis cs:2.25,5.75) {\small $j_{\rm cop}=1$: $(1,\tfrac12)\oplus(\tfrac12,1)$, rank $288$};
\node[anchor=north east] at (axis cs:5.8,1.25) {\small $j_{\rm cop}=2$: Dirac block, rank $160$};
\node[gray,rotate=-34,anchor=south] at (axis cs:4.0,3.6) {\small $C_{\rm cop}+C_\Delta=\tfrac{15}2$};
\end{axis}
\end{tikzpicture}
\caption{The diagonal sector at three constituents (\cref{thm:casid,cor:howe}). The copy Casimir and the Lorentz Casimir add up to a constant
on the whole module, so copy spin is locked to Lorentz type: the Dirac block is exactly the copy quintet.}
\label{fig:dualpair}
\end{figure}

A projector selected by a rotation condition alone need not be Lorentz invariant. The rotation Casimir $-\sum_{i<j}\Delta_{ij}^2$ has spectrum
$\{\tfrac34\,(256),\tfrac{15}4\,(256)\}$, and its rank-$256$ spin-$\tfrac12$ projector commutes with the rotations and with no boost: inside the
irreducible $(1,\tfrac12)$ a boost carries the $j=\tfrac12$ piece into the $j=\tfrac32$ piece. The Lorentz-invariant answer is rank $160$.

% --- source: tex/s10_actions#4 (Gamma matrices on the Dirac block)
\subsection{Gamma matrices on the Dirac block}

The averaged gammas $V_\mu=3^{-1/2}\sum_k\gamma^{(k)}_\mu$ satisfy $\{V_\mu,V_\nu\}=2\eta_{\mu\nu}$ and transform as a Lorentz vector, but they are a
copy vector, not a copy singlet: they map $\mathcal D$ into itself and into the rank-$288$ block with equal norms, and their compression to
$\mathcal D$ degenerates on a rank-$32$ locus. The block does carry gamma matrices; they come from a copy-singlet cubic.

\begin{theorem}[A gauge-invariant $\Cl(1,3)$ on the Dirac block]
\label{thm:blockgamma}
Set $\Psi_{abc}=\varepsilon_{klm}\gamma^{(k)}_a\gamma^{(l)}_b\gamma^{(m)}_c$,
\[
Q_\mu=\tfrac12\eta^{\nu\rho}\Psi_{\mu\nu\rho}
=\sum_\nu\eta_{\nu\nu}\bigl(\gamma^{(1)}_\mu\gamma^{(2)}_\nu\gamma^{(3)}_\nu+\gamma^{(1)}_\nu\gamma^{(2)}_\mu\gamma^{(3)}_\nu+\gamma^{(1)}_\nu\gamma^{(2)}_\nu\gamma^{(3)}_\mu\bigr),
\qquad
G_\mu=\tfrac16\,\omega_{\mathcal S}Q_\mu\big|_{\mathcal D},
\]
with $\omega_{\mathcal S}$ the product of the six surplus generators. Then $\Psi$ is a copy singlet and totally symmetric in $(a,b,c)$; each $Q_\mu$
is a sum of ten grade-three blades, commutes with all eighteen gauge generators and with $C_\Delta$, is a Lorentz vector, and satisfies
$\{Q_\mu,Q_\nu\}|_{\mathcal D}=-72\,\eta_{\mu\nu}$. Consequently, on all $160$ dimensions,
\[
\{G_\mu,G_\nu\}=2\eta_{\mu\nu}\id_{\mathcal D},\qquad [\Delta_{ab},G_c]=\eta_{bc}G_a-\eta_{ac}G_b,\qquad [G_\mu,\mathfrak z(\Delta)]=0,\qquad
\tfrac14[G_a,G_b]=\Delta_{ab}|_{\mathcal D}.
\]
\end{theorem}

\begin{proof}
Reordering $\gamma^{(\sigma1)}\gamma^{(\sigma2)}\gamma^{(\sigma3)}$ into copy order costs $\operatorname{sign}\sigma$, because generators of distinct
constituents anticommute; against the factor $\varepsilon$ this gives $(\operatorname{sign}\sigma)^2=1$, so antisymmetry in the copies is symmetry
in the Lorentz labels. Block preservation follows because a copy singlet commutes with $C_{\rm cop}$ and hence, by \eqref{eq:casid}, with
$C_\Delta$. Finally $\omega_{\mathcal S}$ commutes with $Q_\mu$ and $\omega_{\mathcal S}^2=-\id$, so $\{G_\mu,G_\nu\}=\tfrac1{36}(-\id)(-72\eta_{\mu\nu})$.
The two restricted identities $\{Q_\mu,Q_\nu\}|_{\mathcal D}=-72\,\eta_{\mu\nu}$ and $\tfrac14[G_a,G_b]=\Delta_{ab}|_{\mathcal D}$ are not derived by hand here;
they are certified by exact integer computation on the $160$-dimensional block (a9-1, a9-2).
\end{proof}

The sign flip is not cosmetic. $\Cl(1,3)\cong M_2(\HH)$ has minimal real module of dimension eight, so the four-dimensional Majorana module
carries no real $\Cl(1,3)$; the $Q_\mu$ alone generate $\Cl(3,1)$, and $\omega_{\mathcal S}$, with square $-\id$, converts them. The operator
$\slashed D=G^\mu\nabla_\mu$ is Lorentz covariant and gauge invariant on $\mathcal D$; $G_5=G_0G_1G_2G_3$ has $G_5^2=-\id$, a complex structure
rather than a chirality, as $M_2(\HH)$ requires.

% --- source: tex/s10_actions#5 (The diagonal sector is vector-like)
\subsection{The diagonal sector is vector-like}

\begin{theorem}[Vector-likeness]
\label{thm:diagvectorlike}
Let $W_\mu=\gamma^{(1)}_\mu\gamma^{(2)}_\mu\gamma^{(3)}_\mu$. Then $W_\mu^2=-\eta_{\mu\mu}$, so $W_\mu$ is invertible; $[W_\mu,\mathfrak z(\Delta)]=0$;
$W_\mu$ anticommutes with $\Delta_{ab}$ when $\mu\in\{a,b\}$ and commutes otherwise, hence commutes with $C_\Delta$ and preserves every block;
and $W_\mu$ anticommutes with the gauge-invariant chirality $\omega=E_1\cdots E_{18}$. Therefore $W_0$ is an invertible
$\mathfrak z(\Delta)$-equivariant map exchanging the two chiral halves of every $C_\Delta$ block, and all of $S_3$ is vector-like for $\mathfrak z(\Delta)$.
\end{theorem}

\begin{proof}
Each property is \eqref{eq:bladecomm} with $|I|=3$: against a surplus bivector the sign is $(-1)^{6}$; against $M_{kl}$, the term with $\nu=\mu$
meets $W_\mu$ in two indices and the others in none; against $\Delta_{ab}$ it is $(-1)^{6-1}$ when $\mu\in\{a,b\}$; against $\omega$ it is
$(-1)^{54-3}$. The intertwining conclusion is Schur's lemma.
\end{proof}

By \cref{cor:howe} the reason is structural: $\mathcal D\cong L\otimes U$ with the gauge algebra acting on $U$ only, so chirality and gauge charge
are carried by different tensor factors. Vector-like means that the gauge representation permits a mass for every state, not that any mass is
large; nothing here fixes a scale. A Weyl projection of the block is possible, and it keeps the kinetic operators and the gauge algebra, but it
is then a choice of field content, and at three constituents it carries a colour anomaly (\cref{prop:diagtower}).

% --- source: tex/s10_actions#6 (The candidate gauge algebra of the constituent-wise Lorentz )
\subsection{The candidate gauge algebra of the constituent-wise Lorentz actions}

A Lorentz action on the tower can be taken from each constituent in one of three ways: by the interior algebra $\lint$, by the algebra $\lG$
that centralises the binding operator (\cref{sec:emfield}), or not at all. These constituent-wise actions are the class treated here. Embeddings
of the Lorentz algebra that mix constituents, for example by non-diagonal pairings of the Lorentz bases of different constituents, are not
covered, and the no-go statements below that rely on this classification hold within the tested class (\cref{sec:assessment}).

\begin{theorem}[Candidate algebras]
\label{thm:algebra}
Let the Lorentz algebra act on rung $n$ by $\lint$ on $k_I$ constituents, by $\lG$ on $k_B$ constituents and trivially on $k_S$ constituents,
$k_I+k_B\ge1$. The grade-two elements of $\Cl(3n,3n)$ commuting with the action form
\[
\mathfrak g(k_I,k_B,k_S)\cong\so{k_I}\oplus\so{2k_I+3k_S,\,3k_S}\oplus\so{k_B,\C},
\]
a direct sum of ideals. Its maximal compact subalgebra is $\so{k_I}\oplus\so{2k_I+3k_S}\oplus\so{3k_S}\oplus\so{k_B}$, and the algebra is
compact exactly when $k_S=0$ and $k_B\le1$.
\end{theorem}

\begin{proof}
Grade two is $\so{V_n}$, and $V_n=U\oplus W\oplus F$ with $U=\R^{k_I}\otimes\R^{1,3}$ the interior four-planes, $W$ the Lorentz-trivial directions
(the surplus pairs of the interior constituents and all directions of the spectators, of signature $(2k_I+3k_S,3k_S)$), and $F=\R^{k_B}\otimes V_f$ with
$V_f=(1,0)\oplus(0,1)$ the binding frames. The modules $(\tfrac12,\tfrac12)$, $\mathbf1$ and $V_f$ are pairwise inequivalent and irreducible, so
the mixed parts of $\Lambda^2V_n$ contain no invariant. $\Lambda^2W$ is invariant and gives $\so W$. In
$\Lambda^2U=(\Lambda^2\R^{k_I}\otimes S^2\R^4)\oplus(S^2\R^{k_I}\otimes\Lambda^2\R^4)$ the invariants of $S^2\R^4$ are the multiples of $\eta$ and those of
$\Lambda^2\R^4$ vanish, leaving $\so{k_I}$. On $V_f$ the invariant symmetric forms are spanned by $\eta$ and $\eta(\cdot,G\cdot)$ and there is no invariant
antisymmetric form, leaving $\Lambda^2\R^{k_B}\otimes\operatorname{span}\{\id,G\}\cong\so{k_B,\C}$.
\end{proof}

Here ``compact'' means contained in the chosen compact orthogonal subgroup, the one compatible with a positive kinetic form, in the stated
embedding; it is not a statement about abstract Lie algebras. A compact candidate centraliser is not by itself a gauge theory.

\begin{proposition}[Orthogonal against spin actions]
\label{prop:orthospin}
A copy rotation of interior four-planes is a sum of four commuting terms of square $-\tfrac14$, and one of binding frames a sum of six; their eigenvalues
on $S_n$ are $i$ times integers, and a rotation through $2\pi$ is the identity. A rotation $\tfrac12e_se_t$ of two Lorentz-trivial directions of equal
sign squares to $-\tfrac14$ and its rotation through $2\pi$ is $-\id$. The copy algebras act through orthogonal groups, the algebra of the Lorentz-trivial
directions through spin groups. Consequently the spin content of $S_n$ is half-integer if $k_I$ is odd and integer if $k_I$ is even: fermions require an
odd number of interior constituents.
\end{proposition}

% --- source: tex/s10_actions#7 (The diagonal tower)
\subsection{The diagonal tower}

Under the diagonal interior action the $4n$ four-plane generators and the $2n$ surplus generators $e^{(k)}_2,e^{(k)}_3$ anticommute, so
$S_n\otimes\C\cong S_{4n}\otimes S_{2n}$, with the Lorentz and copy algebras acting on the first factor and the surplus $\so{2n}$ on the second. The
connection algebra is the compact $\so n\oplus\so{2n}$, of dimension $n(5n-3)/2$.

\begin{proposition}[The diagonal tower]
\label{prop:diagtower}
\begin{enumerate}[label=(\roman*),leftmargin=2.2em]
\item For $n\le7$ the factor $S_{4n}$ decomposes under $\lint\oplus\so n$ without multiplicity. Spin is half-integer exactly at odd $n$, and the Dirac
components occur only at odd $n$, each once, tensored with the copy representation of highest weight $(2,\dots,2)$, of dimension $1,5,35,294$ at
$n=1,3,5,7$.
\item At odd $n$ the product $\mathcal J$ of the $2n$ surplus generators is a complex structure commuting with the whole connection algebra, and
$\Gamma=\mathcal J\Omega$ is a chirality. With $\mathcal J$ as $i$ the Dirac block is $[(\tfrac12,0)\oplus(0,\tfrac12)]\otimes V_{(2,\dots,2)}\otimes S^+$,
with $S^+$ a half-spinor of the surplus $\so{2n}$, and the charge $q=-\mathcal J\tilde J/2$ of $\tilde J=\sum_kJ^{(k)}$ takes on $S^+$ the values
\begin{center}
\begin{tabular}{@{}lll@{}}
\toprule
$n$ & $\dim S^+$ & values of $q$ (multiplicity)\\
\midrule
$1$ & $1$ & $\tfrac12$\\
$3$ & $4$ & $-\tfrac32\,(1)$, $\tfrac12\,(3)$\\
$5$ & $16$ & $-\tfrac32\,(5)$, $\tfrac12\,(10)$, $\tfrac52\,(1)$\\
$7$ & $64$ & $-\tfrac72\,(1)$, $-\tfrac32\,(21)$, $\tfrac12\,(35)$, $\tfrac52\,(7)$\\
\bottomrule
\end{tabular}
\end{center}
\item At $n=3$ the centraliser of $\tilde J$ in the surplus $\so6$ is $\uu3$, whose derived algebra is a compact $\su3$; on a Weyl half it acts as
a triplet with $q=\tfrac12$ and a singlet with $q=-\tfrac32$, so that with $\BmL=\tfrac23q$ the triplet has $\BmL=\tfrac13$ and the singlet $-1$:
lepton number as the fourth colour \cite{PatiSalam}. The eight gluons are neutral, the six remaining generators of $\so6$ carry $q=\pm2$ and turn
the triplet into the singlet, as leptoquark bosons do. The copy $\so3$ acts on the Dirac block with Casimir $6$, that is isospin two: it is not weak
isospin.
\item At $n=5$, with colour on three surplus planes and weak isospin on two, $S^+$ is one Standard-Model family with a right-handed neutrino,
$(\mathbf3,\mathbf2)_{1/6}\oplus(\bar{\mathbf3},\mathbf1)_{-2/3}\oplus(\bar{\mathbf3},\mathbf1)_{1/3}\oplus(\mathbf1,\mathbf2)_{-1/2}\oplus(\mathbf1,\mathbf1)_1\oplus(\mathbf1,\mathbf1)_0$,
the spinor $\mathbf{16}$ of $\so{10}$ \cite{FritzschMinkowski,Slansky}, occurring in $35$ copies. At $n=7$, $S^+$ under $\so{10}\oplus\so4$ is two
families and two mirror families, the pattern found for the spinor of $\so{18}$ \cite{WilczekZee}.
\item The Dirac block is vector-like at every odd $n$. Its Weyl projection admits no gauge-invariant mass, since no weight of $S^+$ has its negative in
$S^+$; at $n=3$ it carries a colour anomaly, $\Tr(P_L\mathcal JH^3)=960$ for the colour Cartan element $H=J^{(1)}+J^{(2)}-2J^{(3)}$, and at $n=5$ every
Standard-Model anomaly sum vanishes and $\so{10}$ has no cubic anomaly \cite{GeorgiGlashow72}.
\end{enumerate}
\end{proposition}

\begin{lemma}[No grade-two generator is chiral by itself]
\label{lem:nochiral}
Let $\Gamma$ be, up to sign, a product of $k\ge5$ distinct generators with $\Gamma^2=\id$, and $P_{L,R}=\tfrac12(\id\mp\Gamma)$. A grade-two element $X$
with $XP_R=0$ is zero.
\end{lemma}

\begin{proof}
$XP_R=0$ means $X\Gamma=-X$. For $a\neq b$, $e_ae_b\Gamma$ is a blade of grade $k-2$, $k$ or $k+2$, so $X\Gamma$ has no grade-two component, and
$X\Gamma=-X$ forces $X=0$.
\end{proof}

Under the diagonal and the spectator actions the chirality is, up to sign, the volume element of $\Cl(3n,3n)$, so a gauge interaction built from
connection components can be chiral only through the representation the Weyl block carries.

% --- source: tex/s10_actions#8 (The compact part of the spectator algebra)
\subsection{The compact part of the spectator algebra}
\label{sec:spectatorcompact}

Under the spectator action the connection components form $\so{3n-1,3n-3}$, noncompact for $n\ge2$, and the gauge algebra is its compact part
$\so{3n-1}\oplus\so{3n-3}$, acting through spinors. Let $w$ be the product of the $6n-4$ Lorentz-trivial generators; it commutes with the whole
connection algebra and serves as its complex structure, and $\Gamma=w\,\Omega^{(1)}$ is the chirality.

\begin{proposition}[What the spectator gauge algebra acts on]
\label{prop:spectator}
At $n=2$ the Weyl block is $(\tfrac12,0)\otimes(\mathbf4,\mathbf2)$ under $\spp2\oplus\su2\cong\so5\oplus\so3$: isospin doublets without colour, in a
real representation, so the weak interaction it carries is vector-like. At $n=3$ the Weyl block is
\[
(\tfrac12,0)\otimes\bigl[(\mathbf8_s,\mathbf4)\oplus(\mathbf8_c,\bar{\mathbf4})\bigr]\quad\text{under }\so8\oplus\so6 .
\]
With colour on three negative planes, weak isospin on two positive planes and a family $\so4$ on the remaining two, its $64$ weights are two
Standard-Model families with $\nu^c$ and two mirror families. The cubic $\su4$ sum and the Standard-Model anomaly sums vanish, there are sixteen weak
doublets, the weak roots carry $Q=\pm1$, the gluons are neutral and the leptoquarks carry $\pm\tfrac23$. The block is chiral and anomaly free under
$\so8\oplus\so6$ and vector-like under the Standard-Model algebra.
\end{proposition}

\noindent Executable check: e5-3a, e5-3b, e5-3c for $n=2$ and a7-6 for $n=3$ (exact symbolic or integer arithmetic; \cref{app:methods}).

At $n=3$ the compact spectator algebra therefore already contains the Pati--Salam algebra $\so4\oplus\so6\cong\su2_L\oplus\su2_R\oplus\su4$. The
projector $P_*$ of \cref{sec:pstar} is a way of cutting one chiral family out of this content.

% --- source: tex/s10_actions#9 (Vector bosons rung by rung)
\subsection{Vector bosons rung by rung}
\label{sec:rungs}

\begin{center}
\small
\begin{tabular}{@{}cllcl@{}}
\toprule
$n$ & diagonal connection algebra & bosons & fermions & action on the Dirac block\\
\midrule
$1$ & $\so2=\uu1_J$ & $1$ & yes & one charge, also the complex structure\\
$2$ & $\so2\oplus\so4$ & $7$ & no & integer spin only\\
$3$ & $\so3\oplus\so6\cong\su2\oplus\su4$ & $18$ & yes & $5\times(\mathbf3\oplus\mathbf1)$; the $\su2$ is the copy algebra\\
$4$ & $\so4\oplus\so8$ & $34$ & no & ---\\
$5$ & $\so5\oplus\so{10}$ & $55$ & yes & $35\times\mathbf{16}$, one family each\\
$7$ & $\so7\oplus\so{14}$ & $112$ & yes & $294\times(2\cdot\mathbf{16}\oplus2\cdot\overline{\mathbf{16}})$\\
\midrule
$n$ & spectator gauge algebra & bosons & & action on the Weyl block\\
\midrule
$1$ & $\so2$ & $1$ & & one charge\\
$2$ & $\so5\oplus\so3$ & $13$ & & $(\mathbf4,\mathbf2)$: doublets, no colour, real\\
$3$ & $\so8\oplus\so6$ & $43$ & & two families and two mirror families\\
\bottomrule
\end{tabular}
\end{center}
The tempting assignment ``$n=1$: photon; $n=2$: $W$ and $Z$; $n=3$: gluons'' can now be tested. The photon at $n=1$ is supported: the single
connection component $J$ gauges the Dirac spinor and is also its complex structure. Chiral $W$ and $Z$ at $n=2$ are supported by no action: the
diagonal action has no fermions there, the action that is interior on one constituent and binding on the other has the abelian algebra $\uu1$ only, and
the spectator Weyl block is a real representation. Gluons at $n=3$ are supported, without weak bosons under the diagonal action, where the $\su2$ is
the copy algebra, and with them under the spectator action. Rung three is the first at which some gauge algebra contains
$\su3\oplus\su2\oplus\uu1$ acting on fermions in Standard-Model representations, there together with mirror families; among the actions with a compact
connection algebra the first such rung is $n=5$, through $\so{10}$, with one family per copy and thirty-five copies. Nothing computed here reduces the
copy multiplicity to three, and the rung index counts tensor factors, so none of this identifies particles with constituents.

% --- source: rest_energy/d5_duplication#0 ((head))
% Track d5 -- Duplication: the mathematics of Cl(3,3) -> Cl(6,6).

% --- source: rest_energy/d5_duplication#1 (Duplication: from Cl(3,3) to Cl(6,6))
\section{\texorpdfstring{Duplication: from $\Cl(3,3)$ to $\Cl(6,6)$}{Duplication: from Cl(3,3) to Cl(6,6)}}
\label{d5:sec:main}

The author's second remark asks which physical and mathematical process makes a unitary energy ``result in a duplication'', from
$\Cl(3,3)$ to $\Cl(6,6)$. A duplication acts on four levels at once: on the vector space, $V\to V_2:=V\oplus V$; on the algebra,
$\Cl(V)\to\Cl(V_2)=\Cl(6,6)$; on the spinor module, $S\to S_2=S\otimes S$; and on the energy, $E_6\to E_6\oplus E_6$. This section takes
the four steps in turn, asks which are forced and which are chosen, and records what the duplicate can do that one copy cannot. The
answers, with their classes, are these. The step on the algebra is forced: it is the graded tensor product, its Klein factor is the
Jordan--Wigner sign, and $S_2$ is the Fock space of six fermionic modes (inherited; \cref{d5:sec:algebra}). The step on the energy is a
choice within a three-sphere of energy-preserving, Lorentz-covariant duplications, which is the algebra of a lossless beam splitter. Only a
circle in it, the copy rotation, respects $\eta$, and the two notions of rest energy that coincide on one copy separate on two (fixed;
\cref{d5:sec:energy}). The duplicate carries structures one copy cannot carry: an $\eta$-isometric complex structure, a compact,
Lorentz-invariant copy charge, the copy algebra $\so{2,\C}$, a closure that runs at constant rate with non-null copies ($E_0(\pi^{(k)})\neq0$), and a non-rigid,
Lorentz-invariant vacuum (fixed, partly restated from earlier chapters; \cref{d5:sec:closure,d5:sec:actions,d5:sec:spin,d5:sec:vacuum}). Finally, the
copy rotation is the only admissible $b$-field between two binding frames that commutes with the binding root and the rotations. When
the two rates differ it cannot be removed, and it confines a frame that has no rate of its own; it creates no rate where both frames
are free (fixed, with the mixing formulas inherited; \cref{d5:sec:bfield}). Every finite-dimensional statement was checked by
\texttt{d5\_duplication\_math.py}; identifiers d5-$n$ refer to its \texttt{PASS} lines (\cref{d5:tab:claims}).

\paragraph{Notation.} $V_2=V\oplus V$ has the basis $e_1,\dots,e_6$ of copy~$1$ followed by that of copy~$2$, the metric
$\eta_2=\id_2\otimes\eta$, and the binding root $G_2=\id_2\otimes G$ of both copies (as $G_3$ in \cref{thm:so9C}).
$\varepsilon=\bigl(\begin{smallmatrix}0&-1\\1&0\end{smallmatrix}\bigr)$, $\sigma_x$, $\sigma_z$ act on the copy index, and
$C_v:=\varepsilon\otimes\id_6$ is the copy rotation on vectors. A pair of field strengths is $(z_1,z_2)\in\C^3\oplus\C^3$, with
$q_2:=q(z_1)+q(z_2)$ and duplicated energy $E_6(z_1)+E_6(z_2)$. On $S_2=\R^8\otimes\R^8$ the tower of equation~\cref{eq:tower} gives
$e^{(1)}_a=e_a\otimes\id$, $e^{(2)}_a=\omega\otimes e_a$, $\omega=e_1\cdots e_6$, and $E_{6(k-1)+a}=e^{(k)}_a$. Statements about $V_2$ are
pointwise linear algebra and apply to whichever field is placed there. Where the field matters we name it: the frame field $X$, or the
momentum field $\pi=\eta^{-1}p$ of the closure, to which the Casimirs belong (\cref{sec:closure}, \cref{sec:restclosure}). The
rates $\kappa_k$ of the binding, the spectrum of the self-adjoint part of the shift, are kept apart from the lines of the motion, the
spectrum of the shift itself: a $b$-field can move the lines but never the rates (\cref{d7:thm:schr}(iii)). Reading either as a rest energy through
$mc^2=\hbar\kappa$ is an added identification.

% --- source: rest_energy/d5_duplication#2 (The duplicated algebra: graded tensor product and Klein fact)
\subsection{The duplicated algebra: graded tensor product and Klein factor}
\label{d5:sec:algebra}

\begin{proposition}[The duplicated algebra]
\label{d5:prop:graded}
\begin{enumerate}[label=(\roman*),leftmargin=2.2em]
\item $\omega^2=\id$, $\omega^{\!\top}=\omega$, $\Tr\omega=0$, and $\omega$ anticommutes with every $e_a$. The twelve matrices $e^{(k)}_a$ satisfy
the Clifford relations of $\R^{6,6}$. Their $4096$ blades are signed permutation matrices, and all except $\id$ are traceless, so
$\Cl(6,6)\cong M_{64}(\R)$.
\item For blades $e_I,e_J$ of $\Cl(3,3)$, $e^{(1)}_Ie^{(2)}_J=(e_I\,\omega^{|J|})\otimes e_J$ and $e^{(2)}_Je^{(1)}_I=(-1)^{|I||J|}e^{(1)}_Ie^{(2)}_J$.
Hence $a\,\hat\otimes\,b\mapsto a\,\omega^{|b|}\otimes b$ is an algebra isomorphism $\Cl(3,3)\hat\otimes\Cl(3,3)\to\Cl(6,6)$, where $\hat\otimes$
is the graded tensor product, $(a\hat\otimes b)(c\hat\otimes d)=(-1)^{|b||c|}ac\,\hat\otimes\,bd$.
\item The Klein factor is forced. In $M_{64}(\R)$ the elements commuting with every $e_a\otimes\id$ form $\id\otimes M_8(\R)$, and those
anticommuting with every $e_a\otimes\id$ form $\omega\otimes M_8(\R)$. Any realisation of $\Cl(V_2)$ on $\R^{64}$ that extends
$e_a\mapsto e_a\otimes\id$ is therefore conjugate to the tower by an element of $\id\otimes GL(8,\R)$.
\end{enumerate}
\end{proposition}

\begin{proof}
(i) $\omega^2=(-1)^{15}(-1)^3\id$ by the transpose rule \cref{eq:transpose}, and $\omega$ has even grade in an even dimension. The relations
follow from $(\omega\otimes e_a)(\omega\otimes e_b)=\id\otimes e_ae_b$ and $(e_a\otimes\id)(\omega\otimes e_b)=-(\omega\otimes e_b)(e_a\otimes\id)$. A product
of signed permutation matrices is one, and if $I\neq J$ then $e_I^{\!\top}e_J=\pm e_{I\triangle J}$ is traceless, so the blades are orthogonal
for the trace form; $4096=64^2$. (ii) $e^{(2)}_J=\omega^{|J|}\otimes e_J$, and moving the $|J|$ copy-2 generators past the $|I|$ copy-1
generators costs $(-1)^{|I||J|}$. (iii) By \cref{eq:bladecomm} a blade $e^{(1)}_Ie^{(2)}_J$ commutes with $e^{(1)}_a$ exactly when
$|I|+|J|-[a\in I]$ is even. For all six $a$ this forces either $I=\emptyset$ and $|J|$ even, or $I=\{1,\dots,6\}$ and $|J|$ odd: $64$ blades,
each $\pm\id\otimes e_J$. The anticommuting case is the same with the parities reversed. If $e'_a$ anticommutes with all $e_b\otimes\id$, then
$e'_a=\omega\otimes f_a$, and $(\omega\otimes f_a)(\omega\otimes f_b)=\id\otimes f_af_b$, so the $f_a$ satisfy the relations of $\Cl(3,3)\cong M_8(\R)$. By the
Skolem--Noether theorem $f_a=ue_au^{-1}$, and conjugation by $\id\otimes u$ fixes $e_a\otimes\id$ and $\omega\otimes\id$.
\end{proof}

Everything here is inherited: it is the exponential law $\Cl(V\oplus W)\cong\Cl(V)\hat\otimes\Cl(W)$ \cite{Chevalley54,LawsonMichelsohn89}, and
(i) restates equation~\cref{eq:tower}. The arena contributes only the fixed realisation (d5-1a to d5-1d). The duplication of the
algebra is not a choice; given copy~$1$, copy~$2$ is determined up to conjugation.

\begin{proposition}[Fock space and Jordan--Wigner strings]
\label{d5:prop:fock}
Put $n^\pm_i=\tfrac12(e_i\pm e_{i+3})$, $i=1,2,3$. Then $\{n^\pm_i,n^\pm_j\}=0$ and $\{n^+_i,n^-_j\}=\delta_{ij}$, so $\Cl(3,3)$ is the real
algebra of canonical anticommutation relations of three fermionic modes and $S=\R^8$ is their Fock space. The occupation numbers are
$N_i=n^+_in^-_i=\tfrac12(\id-\omega_i)$, with $\omega_i=e_ie_{i+3}$ the pole operators. The eight eigenlines of \cref{prop:grading}(i) are therefore the eight Fock states, and $(-1)^{N_1+N_2+N_3}=-\omega$. At $n=2$ the tower supplies six modes. The modes of copy~$2$ are
$\omega\otimes n^\pm_i$, each carrying the string $-(-1)^{N^{(1)}}$ of copy~$1$. The total number has spectrum $0,\dots,6$ with multiplicities
$1,6,15,20,15,6,1$, and $(-1)^N$ is the volume element of $\Cl(6,6)$.
\end{proposition}

Inherited. In split signature the vectors $e_i\pm e_{i+3}$ span two maximal isotropic subspaces, and the spinor module is the exterior
algebra of one of them \cite{Chevalley54}. The string is the Jordan--Wigner transformation \cite{JordanWigner28}, and the Klein factor is the
device that makes the fields of two subsystems anticommute \cite{Klein38}. On the level of states, the mathematical process behind
``$\Cl(3,3)\to\Cl(6,6)$'' is the second-quantisation rule $\mathrm{Fock}(\mathcal H\oplus\mathcal H)=\mathrm{Fock}(\mathcal H)\otimes\mathrm{Fock}(\mathcal H)$
(d5-1e).

% --- source: rest_energy/d5_duplication#3 (Centre-of-mass and relative halves)
\subsection{Centre-of-mass and relative halves}
\label{d5:sec:cm}

\begin{proposition}[Two halves, and the exchange]
\label{d5:prop:cmrel}
Let $f_a=(e^{(1)}_a+e^{(2)}_a)/\sqrt2$ and $g_a=(e^{(1)}_a-e^{(2)}_a)/\sqrt2$.
\begin{enumerate}[label=(\roman*),leftmargin=2.2em]
\item $\{f_a,f_b\}=\{g_a,g_b\}=2\eta_{ab}$ and $\{f_a,g_b\}=0$ for all $a,b$, so $\Cl(6,6)=\Cl(f)\hat\otimes\Cl(g)$: two mutually anticommuting
copies of $\Cl(3,3)$.
\item The copy rotation $C_b=\tfrac12\sum_a\eta_{aa}e^{(1)}_ae^{(2)}_a$ acts on vectors by $e^{(1)}_a\mapsto-e^{(2)}_a$, $e^{(2)}_a\mapsto e^{(1)}_a$, that is by $-C_v$, so it lifts the copy rotation with the opposite orientation; it is a sum of
six commuting terms of square $-\tfrac14$, so $\exp(2\pi C_b)=\id$. Conjugation by $\exp(\tfrac\pi4C_b)$ carries $e^{(1)}_a$ to $g_a$ and $e^{(2)}_a$ to
$f_a$, and $C_b=-\tfrac12\sum_a\eta_{aa}f_ag_a$. The split into halves is the split into constituents turned by an eighth of a turn of the
copy rotation, which rotates one half into the other.
\item Let $\omega=(-1)^{\mathsf p}$ on the basis of equation~\cref{eq:realisation}, in which $\omega$ is diagonal, let $\tau(x\otimes y)=y\otimes x$, and let
$D=(-1)^{\mathsf p\otimes\mathsf p}$. The Koszul swap $\tau_K=D\tau$ satisfies $\tau_Ke^{(1)}_a\tau_K=e^{(2)}_a$, $\tau_K^2=\id$, $\Tr\tau_K=0$, and
$\tau_K=g_1g_2\cdots g_6$. So the exchange commutes with every $f_a$ and anticommutes with every $g_a$. The naive swap $\tau$ does not implement
the exchange.
\item The summed clock is $e^{(1)}_1+e^{(2)}_1=\sqrt2f_1$. Every diagonal bivector satisfies $\sum_ke^{(k)}_ae^{(k)}_b=f_af_b+g_ag_b$; this covers the
diagonal interior and binding Lorentz generators and the summed charge $J^{(1)}+J^{(2)}=f_2f_3+g_2g_3$, which act on both halves alike
with no cross term. $G_2$ preserves both halves. The copy rotations $C_b$ and $M_{12}$ (the copy rotation of the interior four-planes, \cref{sec:diagonal}) are pure
cross terms and anticommute with $\tau_K$.
\end{enumerate}
\end{proposition}

\begin{proof}
(i) and (iv) are expansions, using $e^{(1)}_ae^{(2)}_a=-f_ag_a$. (ii) By \cref{eq:bladecomm}, $[C_b,e^{(1)}_b]=-e^{(2)}_b$. The six terms
$\tfrac12\eta_{aa}e^{(1)}_ae^{(2)}_a$ share no index, so they commute, and each squares to $-\tfrac14$; hence
$\exp(\theta C_b)=\prod_a(\cos\tfrac\theta2+\sin\tfrac\theta2\,\eta_{aa}e^{(1)}_ae^{(2)}_a)$, which is $(-1)^6\id$ at $\theta=2\pi$. On vectors $\exp(\theta\ad C_b)$ turns
$e^{(1)}_a$ into $\cos\theta\,e^{(1)}_a-\sin\theta\,e^{(2)}_a$, which is $g_a$ at $\theta=\tfrac\pi4$. (iii) The swap of the copies fixes every $f_a$ and negates every $g_a$: it is the reflection of $V_2$
in the span of the $f_a$ along the span of the $g_a$, and the product $g_1\cdots g_6$ of an even number of unit vectors implements it by conjugation. The Koszul swap implements
it too, by the sign rule. The commutant of $\Cl(6,6)=M_{64}(\R)$ is $\R$, so the two agree up to a scalar, which is $+1$. $\Tr\tau_K$ counts
$\pm1$ over the diagonal states $x_i\otimes x_i$ with sign $(-1)^{\mathsf p_i}$, and $\omega$ is traceless.
\end{proof}

Inherited: centre-of-mass and relative coordinates, and the graded swap, which is the exchange operator of two fermionic subsystems. The
identity $\tau_K=g_1\cdots g_6$, with its sign, and the list in (iv) are facts of the realisation (d5-1f, d5-2a to d5-2c). Because $\tau_K$
anticommutes with the copy rotation, exchanging the constituents reverses the copy charge, as exchanging two real vector fields $W^1,W^2$
exchanges $W^\pm$ (inherited). The halves are not invariant under the copy rotation. Only the exchange selects them, as its two eigenspaces
on vectors.

% --- source: rest_energy/d5_duplication#4 (Duplicating the undivided energy)
\subsection{Duplicating the undivided energy}
\label{d5:sec:energy}

\begin{proposition}[Symmetry of the duplicated energy]
\label{d5:prop:energysym}
\begin{enumerate}[label=(\roman*),leftmargin=2.2em]
\item The linear maps of $V_2$ that commute with the diagonal rotations and with $G_2$ form $\mathfrak{gl}(2)\otimes\operatorname{span}\{\id,G\}\cong\mathfrak{gl}(2,\C)$.
The infinitesimal symmetries of $E_6\oplus E_6$ among them form $\uu2=\operatorname{span}\{C_v,\ \id\otimes G,\ \sigma_z\otimes G,\ \sigma_x\otimes G\}$. Those
in $\so{6,6}$ form $\so{2,\C}=\operatorname{span}\{C_v,\ \varepsilon\otimes G\}$, and those in both form $\so2=\R\,C_v$. The overall duality $\id\otimes G=G_2$, the
relative duality $\sigma_z\otimes G$ (the binding rotation in opposite senses on the two copies) and $\sigma_x\otimes G$ are $\eta_2$-self-adjoint,
like $G$.
\item Inside $\uu6=\so{12}\cap\Cent(G_2)$, of dimension $36$, the diagonal $\uu3$ and $\uu2\otimes\id$ are each other's commutants: $U(3)\times U(2)$ is
a dual pair in $U(6)$.
\end{enumerate}
\end{proposition}

Part (i) restates rows three to five of \cref{prop:commutantsn} (fixed there), with the generators named. Row three is stated there for the diagonal $\lG$, whose commutant in $\so{6,6}$ is the same $\so{2,\C}$ (\cref{d5:prop:so6C}(iii)). The four generators are
those of a lossless beam splitter on two modes \cite{CamposSalehTeich89}: a real splitter $C_v$, an overall phase, a relative phase and a
symmetric splitter. Their group $U(2)$ is the energy-preserving duality group of two free Maxwell fields \cite{GaillardZumino81}, as
\cref{sec:spintower} notes. Part (ii) is inherited from Howe duality for $\C^3\otimes\C^2$ \cite{Howe89} (d5-3a, d5-3b).

\begin{theorem}[Duplication maps]
\label{d5:thm:dup}
\begin{enumerate}[label=(\roman*),leftmargin=2.2em]
\item The $G$-linear, rotation-equivariant maps $V\to V_2$ are $\iota_{\alpha,\beta}(z)=(\alpha z,\beta z)$ with $\alpha,\beta\in\C$. Such a map preserves
the energy, $E_6(\alpha z)+E_6(\beta z)=E_6(z)$, if and only if $|\alpha|^2+|\beta|^2=1$, which gives a three-sphere $U(2)/U(1)$: the first column of
a beam splitter. It is an isometry $(V,\eta)\to(V_2,\eta_2)$ if and only if $\alpha^2+\beta^2=1$ over $\C$, a non-compact complex conic, the orbit
of $\iota_{1,0}$ under the group of the copy algebra $\so{2,\C}$. An energy-preserving $\iota_{\alpha,\beta}$ is an isometry if and only if
$\alpha,\beta$ are real with $\alpha^2+\beta^2=1$, which gives a circle, the orbit of the copy rotation.
\item On the image, $q_2=(\alpha^2+\beta^2)\,q$, $E_0^{(1)}+E_0^{(2)}=(|\alpha|^2+|\beta|^2)E_0$, and the summed Maxwell stress tensor is
$T(\alpha z)+T(\beta z)=(|\alpha|^2+|\beta|^2)\,T(z)$. Every energy-preserving duplication therefore preserves the summed stress tensor, the sum
of the copies' rest energies, and the pair invariant $E_{\rm pair}^2:=\tfrac14T_{\rm tot}{}^\mu{}_\nu T_{\rm tot}{}^\nu{}_\mu$, which equals $E_0^2$ by
Rainich's identity $(T\eta)^2=E_0^2\,\id$ \cite{Rainich1925}. By contrast $|q_2|=|\alpha^2+\beta^2|\,|q|=\bigl(1-4(\Im\alpha\bar\beta)^2\bigr)^{1/2}|q|$ sweeps $[0,|q|]$. The complex invariant $q_2$ is preserved exactly on
the circle. Its modulus $|q_2|$ is preserved exactly where $\Im\alpha\bar\beta=0$, that is on the torus $(\alpha,\beta)=e^{i\psi}(\cos t,\sin t)$, the orbit of the copy
rotation and the overall duality rotation $e^{\psi G_2}$; for example $\iota_{i,0}$ gives $q_2=-q$. At the equal split with relative duality angle $\theta$ ($\beta=e^{i\theta}\alpha$), $|q_2|=|\cos\theta|\,|q|$, which vanishes at
$\theta=\pi/2$.
\item Without equivariance, consider $X\mapsto(\cos\phi\,X,\ \sin\phi\,RX)$ with $R\in U(3)$ and $0<\phi<\pi/2$. It preserves the summed stress tensor
for all $X$ iff $R=e^{\vartheta G}$ is a duality rotation, preserves $q_2$ iff $R\in O(3)$, and preserves both iff $R=\pm\id$. More generally,
a $G$-linear embedding $W=U+iV:\C^3\to\C^6$ that preserves both $E_6\oplus E_6$ and $q_2$ is real, since
$W^*W-W^{\!\top}W=2V^{\!\top}V-2iV^{\!\top}U$.
\end{enumerate}
\end{theorem}

\begin{proof}
(i) The rotations act on $V=\R^2\otimes\R^3$ through the absolutely irreducible second factor, so by Schur's lemma \cite{Schur} the equivariant maps
are $(A,B)\otimes\id_3$ with $A,B\in M_2(\R)$, and $G$-linearity puts $A,B$ in $\operatorname{span}\{\id,G\}\cong\C$. With $G^{\!\top}G=\id$,
$G^{\!\top}\eta G=-\eta$ and $G^{\!\top}\eta=\eta G$, the map $L=\iota_{\alpha,\beta}$ has $L^{\!\top}L=(|\alpha|^2+|\beta|^2)\id$ and
$L^{\!\top}\eta_2L=\Re(\alpha^2+\beta^2)\,\eta+\Im(\alpha^2+\beta^2)\,\eta G$. Since $|\alpha|^2+|\beta|^2-\Re(\alpha^2+\beta^2)=2[(\Im\alpha)^2+(\Im\beta)^2]$, the two
conditions together force real coefficients. (ii) $q$ is holomorphic and quadratic. The components of $T$ are $\tfrac12|z|^2$,
the textbook Poynting vector $-\mathbf S=-\tfrac i2\bar z\times z$, and $\tfrac12\delta_{ij}|z|^2-\Re(\bar z_iz_j)$, all sesquilinear, so $T(\alpha z)=|\alpha|^2T(z)$. Rainich's identity is the
tensor form of the identity $E_6^2=E_0^2+|\mathbf S|^2$ of \cref{sec:energy-identity}. (iii) $T$ determines the Hermitian matrix $zz^*$, the real part from $T^{ij}$ and
the imaginary part from $T^{0i}$, and conversely. So $T(Rz)=T(z)$ for all $z$ iff $R$ commutes with every rank-one Hermitian matrix, iff $R$ is a
phase. $q(Rz)=q(z)$ for all $z$ iff $R^{\!\top}R=\id$, which for unitary $R$ means $\bar R=R$. If $W^*W=W^{\!\top}W=\id$ the identity gives
$V^{\!\top}V=0$.
\end{proof}

All three parts are inherited: they are the linear algebra of a beam splitter acting on $\C^3\otimes\C^2$ (d5-3c to d5-3e). What is not
inherited is the reading.

\begin{remark}[Two rest energies of a duplicate]
\label{d5:rem:twoE0}
On one copy $E_0=\tfrac12|q|$ is at once the holomorphic invariant and the stress-energy invariant; this is Rainich's identity, the field-point
form of $E^2=(pc)^2+(mc^2)^2$ (\cref{sec:energy-identity}). On two copies the two separate. The stress-energy content ($T_{\rm tot}$,
$\sum_kE_0^{(k)}$, $E_{\rm pair}$) is preserved by the whole three-sphere, the holomorphic invariant $q_2$ only by the circle (fixed). Reading this as the
point where duplication changes the meaning of ``rest energy'' is interpretive. Which one deserves the name depends on the group. The invariants of a pair
under the diagonal $\so{3,\C}$ are the entries of the complex Gram matrix $Q_{kl}=z_k\cdot z_l$; the copy algebra $\so{2,\C}$ preserves
$\Tr Q=q_2$ and $\det Q$ but not $Q_{11}$ (d5-3f); and the closure sees $|q_2|$ of the momentum field (\cref{d5:prop:balance}). Two
counter-propagating null fields, $z_1=(1,i,0)$ and $z_2=(1,-i,0)$, have $E_0^{(k)}=0$ and $q_2=0$, yet $E_{\rm pair}^2=2$ (d5-3i). Their rest energy as a pair is the rest-energy density $M_{\rm pair}=2$ of the summed stress tensor (\cref{d6:thm:pair}), not
$E_{\rm pair}=\sqrt2$: $E_{\rm pair}$ is positive exactly when $M_{\rm pair}$ is, and equals it only when $\det Q=0$, as on every duplicate of one field
(\cref{d6:thm:pair}). A nonzero $M_{\rm pair}$ is a kinematic analogue, at one field point, of the invariant mass of two photons, which permits $\gamma\gamma\to e^+e^-$
\cite{BreitWheeler34}; it is a local rest-energy density, not an invariant mass, and the analogy describes no production process. A single null field
duplicated by any $\iota_{\alpha,\beta}$ stays null in each copy, since $q(\alpha z)=\alpha^2q(z)$; that is the analogue of the absence of
$\gamma\to e^+e^-$ in vacuum (the photon kinematics inherited; both analogies interpretive). No isometric linear map copies the energy, $z\mapsto(z,z)$ has $E_6\oplus E_6=2E_6$, and the
genuine duplication of a quantum, $z\mapsto z\otimes z$, is quadratic; that is the no-cloning theorem \cite{WoottersZurek82} (inherited; treating the field point $z$ as a quantum state is interpretive). A
linear duplication divides the undivided energy; it does not multiply it.
\end{remark}

\begin{proposition}[An isometric complex structure exists only on the duplicate]
\label{d5:prop:cx}
\begin{enumerate}[label=(\roman*),leftmargin=2.2em]
\item No $J\in O(3,3)$ satisfies $J^2=-\id$.
\item On $\R^{6,6}$, $C_v$ satisfies $C_v^2=-\id$, $C_v^{\!\top}C_v=\id$ and $C_v^{\!\top}\eta_2C_v=\eta_2$, and commutes with $G_2$ and with the diagonal
$\lG$. It is a compact charge, invariant under the Lorentz group that acts on both frames alike.
\item $\uu6\cap\so{6,6}=\so6$ has dimension $15$; it consists of the paired rotations and is the compact part of $\Cent(G_2)$. At $n=1$ the
corresponding algebra is $\uu3\cap\so{3,3}=\so3$. Inside $\so6$ the centraliser of $C_v$ is
$\uu3'=\{\id\otimes\operatorname{diag}(A,A)+\varepsilon\otimes\operatorname{diag}(B,B):A^{\!\top}=-A,\ B^{\!\top}=B\}$, of dimension $9$, and its part
commuting with the diagonal boosts is $\R\,C_v$.
\end{enumerate}
\end{proposition}

\begin{proof}
(i) If $J^{\!\top}\eta J=\eta$ and $J^2=-\id$, then $h(x,y)=\eta(x,y)+i\,\eta(x,Jy)$ is a Hermitian form on $(\R^6,J)\cong\C^3$ with $\Re h=\eta$. A
Hermitian form of signature $(p,q)$ has real signature $(2p,2q)$, never $(3,3)$. (ii) and (iii) are computed (d5-3g, d5-3h).
\end{proof}

Part (i) is inherited, the signature parity of Hermitian forms. It bites because the arena's triples are odd, and that much is fixed. Read
physically, $w=X^{(1)}+iX^{(2)}$ combines two neutral copies into one field charged under $C_v$, as two real fields form a charged one
(inherited). The pointwise unitary symmetry $U(3)$ of $E_6$ is cut to $O(3)$ by Lorentz compatibility on one copy
(\cref{sec:energy}). It reappears inside $SO(6,6)$ as $U(3)'$, with the copy index as imaginary unit and not the binding root
(interpretive), and $U(3)'$ is not Lorentz invariant.

% --- source: rest_energy/d5_duplication#5 (The duplicated closure symmetry)
\subsection{The duplicated closure symmetry}
\label{d5:sec:closure}

\begin{proposition}[$\Cent(G_2)=\so{6,\C}$ and the copy algebra]
\label{d5:prop:so6C}
\begin{enumerate}[label=(\roman*),leftmargin=2.2em]
\item The centraliser of $G_n=G\oplus\cdots\oplus G$ in $\so{3n,3n}$ has dimension $3n(3n-1)$: $6$, $30$ and $72$ at $n=1,2,3$. It is $\so{3n,\C}$, with
compact part $\so{3n}$ of dimension $3$, $15$ and $36$. At $n=2$ every element is $\bigl(\begin{smallmatrix}R&S\\-S&R\end{smallmatrix}\bigr)$ in the
basis of the six temporal directions followed by their partners, with $R,S$ antisymmetric.
\item $\so{6,\C}=\lG^{(1)}\oplus\lG^{(2)}\oplus(\C^3\otimes\C^3)$, of dimensions $6+6+18$. Under the diagonal rotations the cross part is
$\mathbf1\oplus\mathbf3\oplus\mathbf5$ (complex), with Casimir $0$, $2$, $6$ on $2$, $6$, $10$ real dimensions.
\item The commutant of the diagonal $\lG$ in $\so{6,6}$ is the singlet of (ii), the copy algebra $\so{2,\C}=\operatorname{span}\{C_v,\varepsilon\otimes G\}$. Here
$C_v^2=-\id$ is compact, while $(\varepsilon\otimes G)^2=+\id$ is a copy boost, acting on $(z_1,z_2)$ by $\bigl(\begin{smallmatrix}0&-i\\ i&0\end{smallmatrix}\bigr)$; it
preserves $q_2$ and not $E_6\oplus E_6$.
\item $\Cent(G_2)$ meets $\lint^{(1)}$ in $0$, and the elements of $\so{6,6}$ commuting with both $G_2$ and $\lint^{(1)}$ form $\lG^{(2)}\cong\so{3,\C}$.
\end{enumerate}
\end{proposition}

Part (i) is the $n=2$ member of \cref{thm:so9C}, with the same proof, and the dimension $72$ at $n=3$ re-verifies that theorem.
Part (iii) restates row three of \cref{prop:commutantsn} and \cref{thm:algebra} with $k_B=2$ (fixed there). Parts (ii) and (iv) are fixed,
and (iv) is the $n=2$ analogue of \cref{thm:nocommon} (d5-4a to d5-4c).

\begin{proposition}[Duplicated momentum fields in the closure]
\label{d5:prop:balance}
Let one rotor act on two constituents with the same inertia $N=(\nu_1\eta^{-1}+\nu_2G\eta^{-1})/m$, whose momentum fields are a duplication of one
momentum field: $\pi^{(1)}=\alpha\pi$ and $\pi^{(2)}=\beta\pi$, with $\alpha,\beta$ acting through $\operatorname{span}\{\id,G\}$. Then
$w_1+w_2=(\alpha^2+\beta^2)\,w(\pi)$ with $w=q(z(\pi))$ (equation~\cref{eq:wdef}). The effective strength of \cref{prop:several}
is therefore $R_{\rm eff}=|\nu_1+i\nu_2|\,|q_2(\pi)|/m$: the closure sees the modulus of the holomorphic invariant of the duplicated momentum field. For the equal
split with a relative quarter turn of duality, $\pi^{(2)}=G\pi^{(1)}$, equivalently $p^{(2)}=-Gp^{(1)}$, one has $w_2=-w_1$ and $R_{\rm eff}=0$. The
rate is then exactly constant, and nonzero when $h(0)\neq0$, while each copy keeps $E_0(\pi^{(k)})=\tfrac12E_0(\pi)\neq0$. On one constituent with
the same inertia and $(\nu_1,\nu_2)\neq(0,0)$ a constant nonzero rate requires $E_0(\pi_0)=0$ (\cref{cor:constrate}); for
$(\nu_1,\nu_2)=(0,0)$ the rate is constant for all data.
\end{proposition}

This restates \cref{prop:several} and \cref{cor:constrate} (fixed there). The expression of $R_{\rm eff}$ through $q_2$ for duplicates is a
one-line consequence (fixed; d5-9f). So, for $(\nu_1,\nu_2)\neq(0,0)$, a duplicate can do what one copy cannot: run at an exactly constant rate while its copies carry nonzero
rest energy in the momentum field, in the sense that $E_0(\pi^{(k)})\neq0$. Reading $E_0(\pi^{(k)})$ as a rest energy is the interpretive dictionary stated after \cref{en:closure}. The identification $mc^2=\hbar\kappa$ concerns the common rate $\kappa=h/I_0$ of both copies, which does not depend on $E_0(\pi^{(k)})$, so it does not make $E_0(\pi^{(k)})$ a mass (interpretive).

% --- source: rest_energy/d5_duplication#6 (The five Lorentz actions at the second rung)
\subsection{The five Lorentz actions at the second rung}
\label{d5:sec:actions}

What the duplication does depends on how the Lorentz algebra acts on the two copies. \cref{thm:algebra} lists the choices, and at $n=2$
there are five (\cref{d5:tab:actions}).

\begin{table}[t]
\centering
\small
\begin{tabular}{@{}lccccll@{}}
\toprule
action & $(k_I,k_B,k_S)$ & algebra & dim & compact & content of $\R^{6,6}$ & spin on $S_2$\\
\midrule
spectator & $(1,0,1)$ & $\so{5,3}$ & $28$ & $13$ & $V\oplus8\cdot\mathbf1$ & half-integer\\
binding $+$ spectator & $(0,1,1)$ & $\so{3,3}$ & $15$ & $6$ & $\Lambda^2V\oplus6\cdot\mathbf1$ & integer\\
diagonal interior & $(2,0,0)$ & $\so2\oplus\so4$ & $7$ & $7$ & $2V\oplus4\cdot\mathbf1$ & integer\\
mixed & $(1,1,0)$ & $\so2$ & $1$ & $1$ & $V\oplus\Lambda^2V\oplus2\cdot\mathbf1$ & half-integer\\
binding on both & $(0,2,0)$ & $\so{2,\C}$ & $2$ & $1$ & $2\,\Lambda^2V$ & integer\\
\bottomrule
\end{tabular}
\caption{The Lorentz actions at $n=2$ and the grade-two elements commuting with them (\cref{thm:algebra}), with their compact parts,
the content of the vector space (Lorentz Casimir $3$ on a four-vector $V$, $4$ on a field strength $\Lambda^2V$), and the effect of a $2\pi$
rotation on $S_2$. Every entry was computed in the realisation (d5-5a, d5-5b, d5-6b).}
\label{d5:tab:actions}
\end{table}

Read row by row: under the spectator action copy~$2$ is a passenger whose six directions join a non-compact internal algebra. With
$\lG$ on copy~$1$, copy~$2$'s whole $\so{3,3}$ becomes internal symmetry and no spin one-half survives. The diagonal interior action gives a
compact algebra containing the copy rotation $M_{12}$, but only integer spin. The mixed action keeps only $J_1=e^{(1)}_2e^{(1)}_3$. With $\lG$ on both
copies the frames are two field strengths mixed by $\so{2,\C}$. The table restates \cref{thm:algebra} and \cref{prop:orthospin}, \cref{prop:mixed} and \cref{prop:mixedcost} (fixed there), re-verified. For the mixed rung we also re-verified the Casimir identity
$(C_{\rm mix}-\tfrac32)(C_{\rm mix}-\tfrac{11}2)=0$ with ranks $16$ and $48$ (d5-5b). We checked the grade-three element $\Theta_F$ of \cref{sec:mixed}, which identifies the frame field of copy~$2$ with $dA$ of copy~$1$. It and $\Theta_F$ followed by $G$ are invariant and span the
invariants in $\Lambda^2(\text{four-plane}_1)\otimes V^{(2)}$, of dimension two as stated there. Here $\text{four-plane}_1=\langle e^{(1)}_1,e^{(1)}_4,e^{(1)}_5,e^{(1)}_6\rangle$ is the interior four-plane of copy~$1$ and $V^{(2)}$ is the span of the $e^{(2)}_a$, the space of copy~$2$ alone, which \cref{sec:mixed} writes $V_2$. The full space of invariant grade-three elements has
dimension four: the other two lie in $\Lambda^3V^{(2)}$, among the trivectors of copy~$2$ alone, the two $\lG$-invariant trivectors of
$(\Lambda^3\R^6)_\C=2\,(0,0)\oplus2\,(1,1)$ (\cref{sec:readingsexclude}; d5-5c). This is a precision, not a correction. The mixed rung is the only
action at $n=2$ that houses both readings, at the price of the whole internal algebra.

% --- source: rest_energy/d5_duplication#7 (Spin and statistics of the duplicate)
\subsection{Spin and statistics of the duplicate}
\label{d5:sec:spin}

\begin{proposition}[Spin of the duplicate]
\label{d5:prop:spin}
\begin{enumerate}[label=(\roman*),leftmargin=2.2em]
\item The copy rotations $M_{12}$ (four commuting terms of square $-\tfrac14$) and $C_b$ (six) have eigenvalues $i\{-2,\dots,2\}$ and $i\{-3,\dots,3\}$, and
$\exp(2\pi M_{12})=\exp(2\pi C_b)=\id$. A surplus rotation $\tfrac12e_se_t$ has $\exp(\pi e_se_t)=-\id$.
\item A rotation through $2\pi$ is $+\id$ on $S_2$ under the diagonal interior action and $-\id$ under the spectator action. Under the mixed action it
is $-\id$, and under $\lG$ on one or both copies it is $+\id$. The spin Casimir is $0$ on $16$ and $2$ on $48$ dimensions for the diagonal action
($j=0,1$), and $\tfrac34$ on all of $S_2$ for the spectator action.
\item Under the diagonal interior action the Lorentz Casimir is $0,3,4$ on $8,32,24$ dimensions and $C_{\rm cop}=-M_{12}^2$ is $0,1,4$ on $24,32,8$, with
$C_{\rm cop}+C_\Delta=4$. The exchange $\tau_K$ commutes with the diagonal Lorentz action and anticommutes with $M_{12}$, and its $\pm1$ eigenspaces
carry identical Lorentz content, $0,3,4$ on $4,16,12$ dimensions each.
\end{enumerate}
\end{proposition}

This restates \cref{prop:orthospin}, \cref{thm:casid} at $n=2$ ($\tfrac12n(n+2)=4$) and the last paragraph of \cref{sec:mixed} (fixed there),
with the $2\pi$ eigenvalues computed (d5-6a, d5-6b). That two spin-$\tfrac12$ constituents make an integer-spin composite is inherited
\cite{EhrenfestOppenheimer31}: under the diagonal action the duplicate of a fermion is a boson, and fermions need an odd number of interior
constituents. The Lorentz isomorphism of the two exchange sectors in (iii) is fixed (d5-6c). It holds because exchange reverses the copy
charge, pairing charge $\pm m$, and splits the neutral sector evenly; there is no ``singlet antisymmetric, triplet symmetric'' rule once the
surplus directions and the Koszul signs are included.

% --- source: rest_energy/d5_duplication#8 (What the second rung adds: the vacuum and the imaginary unit)
\subsection{What the second rung adds: the vacuum and the imaginary unit}
\label{d5:sec:vacuum}

\begin{proposition}[The photon vacuum and the imaginary unit, re-verified]
\label{d5:prop:vacuum}
\begin{enumerate}[label=(\roman*),leftmargin=2.2em]
\item At $n=1$ the commutant of $\lint$ and $Q=\tfrac12e_2e_3$ in $M_8(\R)$ has dimension four, and its symmetric part is
spanned by $\id$ and $\omega=e_1\cdots e_6$. The rank-four projectors in it are $\tfrac12(\id\pm\omega)$, and every bivector commutes with
$\omega$, so they are rigid (\cref{prop:rigid}).
\item (\cref{thm:towervac} and~\cref{thm:S4}) With $a,b=e^{(1)}_2,e^{(1)}_3$ and $c,d,e=e^{(2)}_1,e^{(2)}_2,e^{(2)}_3$: $H=abcd$ has $H^2=\id$, $H^{\!\top}=H$
and $\Tr H=0$, and $\Pi_0=\tfrac12(\id+H)$ commutes with $Q=\tfrac12ab$, with the four interior gammas and with the internal volume. Moreover
$\pi=[\tfrac12be,\Pi_0]=-\tfrac12acde$ is symmetric with $\Tr\pi^2=16$ and $[Q,\pi]=\tfrac12bcde$. The $Z_i=I_5v_i$ are symmetric of grade four with
$\{Z_i,Z_j\}=2\delta_{ij}$. $\Pi(n)$ is a rank-$32$ projector commuting with the gammas (checked at $n=(2,3,6,0,0)/7$). The tangents
$T_i=\tfrac12I_5v_i$ have $\Tr T_iT_j=16\delta_{ij}$, $[Q,T_a]=-T_b$, $[Q,T_b]=T_a$ and $[Q,T_c]=[Q,T_d]=0$.
\item The grade-four blades in the five positive internal generators that commute with $Q$ are exactly $abcd$, $abce$ and $abde$. Every
such vacuum straddles the two copies.
\item (\cref{thm:i} at $n=2$) The product of the eight internal generators is $J=(e_2e_3)\otimes\omega$, with $J^2=-\id$. The joint commutant
of $\lint^{(1)}$ and the $28$ internal bivectors in $M_{64}(\R)$ is spanned by the blades $\id,\Omega,J,\Omega J$. The solutions of $x^2=-\id$ in it are
exactly $\pm\Omega$ and $\pm J$, and only $\pm J$ commute with the gammas.
\end{enumerate}
\end{proposition}

Parts (i), (ii) and (iv) restate earlier results (fixed there), re-verified in the realisation of $\Cl(6,6)$ (d5-7a to d5-7c, d5-8).
Part (iii) is fixed but elementary (d5-7d). The vacuum of \cref{thm:towervac} is built on four positive internal directions, two of them the surplus pair $a,b$ of copy~$1$. A
grade-four blade in the five positive internal generators commutes with $Q$ exactly when it contains both $a$ and $b$, so it straddles the copies.
Part (iii) covers only these positive blades, not blades that also use the negative internal directions of copy~$2$. At one copy the first rung is
rigid by the commutant argument of \cref{prop:rigid}, without reference to blades, so the vacuum exists because there are two copies. Part (iv) adds one
identity: the imaginary unit of the duplicate is the imaginary unit of copy~$1$ times the volume element of copy~$2$, which is that copy's
fermion parity up to sign (\cref{d5:prop:fock}). It is neither the embedded one-rung charge $J_1\otimes\id$ nor $J^{(1)}+J^{(2)}$.

% --- source: rest_energy/d5_duplication#9 (The copy rotation as a b-field between two manifolds)
\subsection{\texorpdfstring{The copy rotation as a $b$-field between two manifolds}{The copy rotation as a b-field between two manifolds}}
\label{d5:sec:bfield}

On one copy a $b$-field commuting with $G$ is removable along a worldline (\cref{bf:larmor}), and a purely anticommuting one tears
the binding (\cref{bf:pG}). An admissible one with both parts keeps every orbit bounded when the two triples rotate at equal rates,
outside one exceptional ratio, and it is then not removable (\cref{bf:screen,bf:remove}). On two copies the binding term is $K\otimes G$ with $K=\operatorname{diag}(\kappa_1,\kappa_2)$. That is a
multiple of $G_2$ only when the rates agree, and the following theorem measures what changes. The interior Lagrangian of each copy is that
of \cref{sec:bfield}, with $S=K\otimes G+b$; the frame field obeys $\dot X=-SX+\pi/m$ and the momentum field $\dot\pi=S^\dagger\pi$.

\begin{theorem}[The inter-copy $b$-field]
\label{d5:thm:bfield2}
Let $\kappa_1,\kappa_2\ge0$, $\beta\in\R$, $b=\beta\,C_v$, and $M=\bigl(\begin{smallmatrix}-S&\id/m\\0&S^\dagger\end{smallmatrix}\bigr)$ on $(X,\pi)$.
\begin{enumerate}[label=(\roman*),leftmargin=2.2em]
\item Up to scale, $b$ is the only element of $\so{6,6}\cap\so{12}$ commuting with $G_2$ and the diagonal rotations. It is compact, and
$S^\dagger=K\otimes G-b$.
\item $\det(\lambda-S)=\bigl(\lambda^4+(\kappa_1^2+\kappa_2^2+2\beta^2)\lambda^2+(\kappa_1\kappa_2-\beta^2)^2\bigr)^3$, so $\spec S=\{\pm i\omega_\pm\}$ with
$\omega_\pm=\tfrac12(\kappa_1+\kappa_2)\pm\sqrt{\tfrac14(\kappa_1-\kappa_2)^2+\beta^2}$; equivalently, $\omega_\pm$ are the eigenvalues of
$\bigl(\begin{smallmatrix}\kappa_1&i\beta\\-i\beta&\kappa_2\end{smallmatrix}\bigr)$.
\item $[K\otimes G,b]=\beta(\kappa_2-\kappa_1)\,\sigma_x\otimes G$, so for $\beta\neq0$, $b$ is removable along the worldline (Larmor \cite{Larmor1897}) iff $\kappa_1=\kappa_2$. The
centraliser of $K\otimes G$ in $\so{6,6}$ has dimension $30$ if $\kappa_1=\kappa_2\neq0$, $12$ if $\kappa_1\neq\kappa_2$ are both nonzero, and $21$ if
exactly one vanishes.
\item $(M^2+\omega_+^2)(M^2+\omega_-^2)=0$ identically in $\kappa_1,\kappa_2,\beta,m$.
\item Every orbit is bounded iff either $\kappa_1=\kappa_2>0$, or $\kappa_1\neq\kappa_2$ and $\kappa_1\kappa_2\neq\beta^2$. In particular, for $\kappa_1=0<\kappa_2$
the orbits are bounded iff $\beta\neq0$: the inter-copy field confines a frame with no rate of its own. For $\kappa_1=\kappa_2=0$ they are unbounded
for every $\beta$.
\item For $\kappa_1=0<\kappa_2$, $\omega_+\omega_-=-\beta^2$ and $\omega_++\omega_-=\kappa_2$. For $\beta\ll\kappa_2$ the slow line is
$|\omega_-|=\beta^2/\kappa_2-\beta^4/\kappa_2^3+\cdots$, the seesaw. For $\beta\gg\kappa_2$ both lines approach $\pm\beta$, and the orbit size grows like
$2/(m\kappa_2)$.
\end{enumerate}
\end{theorem}

\begin{proof}
(i) is \cref{d5:prop:energysym}(i), and $C_v$ is $\eta_2$-antisymmetric while $K\otimes G$ is self-adjoint. (ii) On the sector $G_2=i$,
$S=(iK+\beta\varepsilon)\otimes\id_3$, whose characteristic polynomial is $\lambda^2-i(\kappa_1+\kappa_2)\lambda-\kappa_1\kappa_2+\beta^2$; the sector $G_2=-i$ is
its conjugate. (iii) $K\varepsilon-\varepsilon K=(\kappa_2-\kappa_1)\sigma_x$. If $\kappa_1=\kappa_2=\kappa$, then $X=e^{-b\tau}Y$ and $\pi=e^{-b\tau}\varpi$ turn
the system into $b=0$ with rate $\kappa$; if $\beta\neq0$ and $\kappa_1\neq\kappa_2$, this substitution makes the binding $e^{b\tau}(K\otimes G)e^{-b\tau}$ time-dependent. The
centraliser dimensions are computed. (iv) is a symbolic identity (d5-9c). (v) If $\kappa_1\neq\kappa_2$ and $\kappa_1\kappa_2\neq\beta^2$, then $\omega_+\neq\omega_-$,
$\omega_-\neq0$ and $\omega_++\omega_-=\kappa_1+\kappa_2>0$, so the polynomial of (iv) has four simple imaginary roots and $M$ is semisimple with imaginary
spectrum. If $\kappa_1=\kappa_2=\kappa>0$, use (iii) and the boundedness of the $b=0$ system (\cref{thm:neutral}); if $\kappa=0$ the substituted
system is free drift. If $\kappa_1\neq\kappa_2$ and $\kappa_1\kappa_2=\beta^2$, then $\omega_-=0$, and $M$ is semisimple at $0$ iff $\ker S^\dagger\subset\operatorname{im}S$.
Because $S^{\!\top}=-S$ and $S^\dagger=-\eta_2S\eta_2$, this says that $\ker S$ is $\eta_2$-totally isotropic. Complexifying, $\ker S$ is spanned by
the real parts of $v\otimes c$ in the $G_2=i$ sector, with $c\in\C^3$ in the basis $e_j-ie_{j+3}$ and $v=(\beta,i\kappa_1)$ (or $v=(1,0)$ when $\beta=\kappa_1=0$), and
$\eta_2(\Re(v\otimes c),\Re(v\otimes c'))=\Re\bigl[(v\cdot v)(c\cdot c')\bigr]$. This vanishes for all $c,c'$ iff $v\cdot v=\beta^2-\kappa_1^2=0$, which together
with $\kappa_1\kappa_2=\beta^2$ means $\kappa_1=\kappa_2$ or $\beta=\kappa_1=0$, and in the second case $v=(1,0)$ is not isotropic; so there is a Jordan block.
(vi) Vieta and expansion. The orbit sizes are numerical (d5-9e, d5-9g).
\end{proof}

Parts (ii) and the seesaw in (vi) are inherited: two-level mixing with detuning $\tfrac12(\kappa_1-\kappa_2)$ and coupling $\beta$, the generalised Rabi
splitting, and the seesaw formula \cite{Minkowski77}. Part (i) restates \cref{prop:commutantsn} (fixed there). Parts (iii)--(v) are fixed
(d5-9a to d5-9d). The exact checks are: bounded at $(\kappa_1,\kappa_2,\beta)=(1,3,\tfrac25),(0,5,\tfrac12),(2,2,\tfrac13),(1,1,1),(2,2,0)$ and
$(0,\tfrac1{100},1)$; unbounded at $(0,0,\tfrac12),(1,4,2),(1,9,3),(0,3,0)$. At $(0,5,\tfrac12)$, $\omega_-=-0.04951$ against $\beta^2/\kappa_2=0.05$, and
$\sup_{t\le600}\|e^{Mt}\|=19.86$ against $1/|\omega_-|=20.2$. At $(0,\tfrac1{100},1)$ the lines are $1.00501$ and $-0.99501$, and the supremum over a beat
period is $200.0=2/(m\kappa_2)$. So the seesaw orbit is bounded, with amplitude of order $1/(m|\omega_-|)$ for $\beta\ll\kappa_2$ and $2/(m\kappa_2)$ for
$\beta\gg\kappa_2$.

\begin{remark}[The author's first remark at two manifolds]
\label{d5:rem:dealbreaker}
On one manifold a $b$ commuting with $G$ is removable along a worldline, and a constant connection alone gives free drift
(\cref{bf:remove,bf:alone}). An admissible $b$ that does not commute with $G$ can already there be non-removable, change the spectrum
and keep every orbit bounded (\cref{bf:screen}), so these two features are not new at $n=2$. What two manifolds with different rates add is
this. The compact inter-copy field commutes with $G_2$, so on one manifold a field of that kind would be removable, but it does not commute with
the binding term $K\otimes G$. It is therefore not removable, it shifts the lines, and it confines a frame with $\kappa_1=0$ by coupling it to a
frame with $\kappa_2>0$. The no-go of \cref{bf:alone} still holds at $n=2$: for $\kappa_1=\kappa_2=0$ the motion is free drift for every
$\beta$. In that sense, and only in it, the $b$-field is a ``deal breaker'' at $n\ge2$ (fixed). Two qualifications remain. The field can set the lines, which approach $\pm\beta$ when $\beta\gg\kappa_2$, but not the rates, and the size of
the orbit is set by $\kappa_2$ and diverges as $\kappa_2\to0$. And reading a line or a rate as a rest energy still needs $mc^2=\hbar\kappa$ (an added
identification).
\end{remark}

% --- source: rest_energy/d5_duplication#10 (Assessment)
\subsection{Assessment}
\label{d5:sec:assess}

\paragraph{What is forced, what is chosen, what is added.} \emph{Forced} (inherited): the algebra step, $\Cl(V\oplus V)=\Cl(V)\hat\otimes\Cl(V)$, with the
Klein factor determined up to conjugation; $S_2=S\otimes S$ as the Fock space of six modes; and the exchange as the Koszul swap.
\emph{Chosen} (constructed): the Lorentz action, one of five; the energy duplication, a point of the three-sphere, or of the circle if $\eta$ is
to be respected, with the angle $\phi$ free; and the vacuum point on $S^4$. \emph{Added} at $n=2$ (fixed; the copy algebra, the copy charge, the constant-rate closure, the vacuum and the mixed rung restate earlier chapters, while the complex-structure statement of \cref{d5:prop:cx}, the comparison $\so6\supset\uu3'$, the form of $R_{\rm eff}$ through $q_2$ and the inter-copy $b$-field are fixed here): the copy algebra $\so{2,\C}$ on two
binding frames (dimension $0$ at $n=1$); an $\eta$-isometric complex structure and a compact, Lorentz-invariant copy charge (none at $n=1$); a
Lorentz-compatible energy symmetry $\so6\supset\uu3'$, against $\so3$; a closure with constant rate and non-null copies, $E_0(\pi^{(k)})\neq0$; a non-rigid,
Lorentz-invariant vacuum; the mixed rung, which houses both readings; and an inter-copy $b$-field that confines (d5-10 collects the
dimensions). \emph{Lost}: spin one-half under the diagonal interior action, and the uniqueness of the photon direction (\cref{cor:tradeoff}).

\paragraph{Processes.} \Cref{d5:tab:processes} lists the physical processes that take one copy to two, and what each preserves.

\begin{table}[t]
\centering
\footnotesize
\begin{tabular}{@{}>{\raggedright\arraybackslash}p{3.3cm}>{\raggedright\arraybackslash}p{4.7cm}>{\raggedright\arraybackslash}p{4.3cm}>{\raggedright\arraybackslash}p{2.4cm}@{}}
\toprule
process & mathematics & preserves & class\\
\midrule
splitting one field into two (beam splitter) & $\iota_{\alpha,\beta}$, $S^3=U(2)/U(1)$; circle if also $\eta$-isometric & $E_6$, $T$, $\sum E_0^{(k)}$; $q_2$ only on the circle & inherited; reading fixed\\
second quantisation & $\Cl(V\oplus V)=\Cl(V)\hat\otimes\Cl(V)$; Fock space; Jordan--Wigner strings & Clifford relations & inherited (forced)\\
two neutral copies into one charged field & $C_v$ as $i$, $w=X^{(1)}+iX^{(2)}$ & $\eta_2$, $E_6\oplus E_6$ & inherited; exists only at $n\ge2$: fixed\\
two-body composite & diagonal interior action; exchange $\tau_K$ & integer spin & inherited; restated\\
pair creation from null energy & linear covariant duplication keeps null copies; two non-collinear null fields have $E_{\rm pair}>0$ & $T$ & inherited; pair-creation reading interpretive\\
cloning, $z\mapsto z\otimes z$ & quadratic, not isometric and linear & --- & obstructed (inherited)\\
closure balance & $\pi^{(2)}=G\pi^{(1)}$, $R_{\rm eff}\propto|q_2(\pi)|=0$ & constant rate, $E_0^{(k)}\neq0$; bounded frames if $h(0)\neq0$ & restated (\cref{prop:several})\\
inter-copy $b$-field & $S=K\otimes G+\beta C_v$ & bounded iff \cref{d5:thm:bfield2}(v) & fixed; mixing inherited\\
\bottomrule
\end{tabular}
\caption{Processes that take one copy of the arena to two.}
\label{d5:tab:processes}
\end{table}

\paragraph{The precise sense of ``a unitary energy results in a duplication''.} Three statements, in increasing order of what the arena adds.
First, the undivided energy $(\C^3,E_6,U(3))$ embeds isometrically in $(\C^3\otimes\C^2,E_6\oplus E_6,U(6))$. The Lorentz-covariant embeddings
form the three-sphere of a beam splitter, and the diagonal $U(3)$ and the copy $U(2)$ are a dual pair (inherited). Second, compatibility with
$\eta$ cuts the copy $U(2)$ to $O(2)$ and $U(6)$ to $O(6)$, since $U(n)\cap O(n,\C)=O(n)$. The copy rotation $SO(2)$, the identity component of
$O(2)$, is the one connected family that preserves the energy, the metric, the binding root and the diagonal Lorentz group at once; the other
component contains the swap of the copies and $\operatorname{diag}(1,-1)$. Neither exists on one copy (fixed). Third, the duplicate's rest
energy splits into a stress-energy part, preserved by every duplication, and a holomorphic part. The complex invariant $q_2$ is preserved only on the circle of the copy rotation. Its modulus $|q_2|$, which is what
the closure sees, is preserved on that circle composed with the overall duality rotation $G_2$, which multiplies $q_2$ by a phase (fixed). Read literally as a multiplication of the energy, the phrase fails either way: the classical map $z\mapsto(z,z)$ is linear and doubles
the quadratic energy, $(E_6\oplus E_6)(z,z)=2E_6(z)$ (\cref{d5:rem:twoE0}), and its normalised version $z\mapsto(z,z)/\sqrt2$ preserves the total
energy and divides it equally between the copies. Copying classical field data is well defined; the no-cloning theorem, which concerns the
universal copying of an unknown quantum state by a physical quantum operation, is a separate statement (inherited).

\begin{table}[t]
\centering
\footnotesize
\begin{tabular}{@{}>{\raggedright\arraybackslash}p{8.1cm}>{\raggedright\arraybackslash}p{5.0cm}>{\raggedright\arraybackslash}p{2.0cm}@{}}
\toprule
statement & class & checks\\
\midrule
\cref{d5:prop:graded}: graded tensor product, $\Cl(6,6)=M_{64}(\R)$, Klein factor forced & inherited (restated: \eqref{eq:tower}) & d5-1a to d5-1d\\
\cref{d5:prop:fock}: Fock space, Jordan--Wigner strings, pole operators $=\id-2N_i$ & inherited (reading of \cref{prop:grading}) & d5-1e\\
\cref{d5:prop:cmrel}: halves, $\pi/4$ copy rotation, $\tau_K=g_1\cdots g_6$, what each half carries & inherited; identities fixed & d5-1f, d5-2a to d5-2c\\
\cref{d5:prop:energysym}: $\uu2$, $\so{2,\C}$, $\so2$; dual pair & fixed (restated: \cref{prop:commutantsn}); inherited & d5-3a, d5-3b\\
\cref{d5:thm:dup}: duplication maps, $S^3$ and circle, $q_2$ versus $T$ & inherited; reading fixed & d5-3c to d5-3e\\
\cref{d5:rem:twoE0}: two rest energies; Gram invariants; pair invariant; no cloning & fixed; inherited; analogies interpretive & d5-3f, d5-3i\\
\cref{d5:prop:cx}: isometric complex structure only at $n\ge2$; $\so6\supset\uu3'$ & inherited (parity); fixed & d5-3g, d5-3h\\
\cref{d5:prop:so6C}: $\Cent(G_n)=\so{3n,\C}$; $\mathbf1\oplus\mathbf3\oplus\mathbf5$; copy algebra; analogue of \cref{thm:nocommon} & fixed (restated: \cref{thm:so9C}, \cref{prop:commutantsn}) & d5-4a to d5-4c\\
\cref{d5:prop:balance}: closure sees $|q_2(\pi)|$; balance with non-null copies & fixed (restated: \cref{prop:several}, \cref{cor:constrate}) & d5-9f\\
\cref{d5:tab:actions}: five Lorentz actions; mixed rung; $\Theta_F$ (dimension $2$ within $\Lambda^2(\text{four-plane}_1)\otimes V^{(2)}$, $4$ in all) & fixed (restated: \cref{thm:algebra}, \crefrange{prop:mixed}{prop:mixedcost}) & d5-5a to d5-5c\\
\cref{d5:prop:spin}: spin and statistics, Casimirs, exchange sectors & fixed (restated: \cref{prop:orthospin}, \cref{thm:casid}); inherited & d5-6a to d5-6c\\
\cref{d5:prop:vacuum}: rigid first rung, $S^4$ vacuum, straddling, $J=J_1\otimes\omega$ & fixed (restated: \cref{prop:rigid}, \cref{thm:towervac}, \cref{thm:S4}, \cref{thm:i}) & d5-7a to d5-7d, d5-8\\
\cref{d5:thm:bfield2}: inter-copy $b$-field: uniqueness, spectrum, removability, boundedness criterion, seesaw & fixed; mixing and seesaw inherited & d5-9a to d5-9e, d5-9g\\
\cref{d5:rem:dealbreaker}: the deal-breaker reading & fixed; obstructed for $\kappa_1=\kappa_2=0$; $mc^2=\hbar\kappa$ an identification & ---\\
\cref{d5:sec:assess}: comparison $n=1$ against $n=2$ & fixed; mostly restated (\cref{prop:commutantsn}) & d5-10\\
\bottomrule
\end{tabular}
\caption{Claims of \cref{d5:sec:main}, their classes, and the checks printed by \texttt{d5\_duplication\_math.py}. ``Restated'': established in an earlier chapter, restated and re-verified here.}
\label{d5:tab:claims}
\end{table}

% --- source: emergence/e1_two_from_one#0 ((head))
% Track e1 -- One particle or two: the second rung Cl(6,6).
% New references: sections/bib_e1.tex.

% --- source: emergence/e1_two_from_one#1 (One particle or two: the second rung)
\section{One particle or two: the second rung}
\label{e1:sec:main}

The conventions of \cref{sec:conventions} say that a constituent is a tensor factor $\Cl(3,3)$ of the tower, that the word never refers to a
particle, and that three constituents do not describe three particles (\cref{sec:conventions}). The duplication
$\Cl(3,3)\to\Cl(6,6)$ of \cref{d5:sec:main} is therefore not yet the emergence of two particles from one. This section
asks in what precise sense, and under which added structures, the second rung $\Cl(6,6)$ carries two particles rather than one
particle with a larger internal space. Four tests separate the two readings: whether the copies have separate mass shells, whether
there is one complex structure under which both copies are complex Dirac spinors, whether some operator counts particles, and whether
exchange of the copies sorts states by statistics. The answer has three parts. In the graded tower the two mass
terms anticommute, they add in quadrature, and no state has a sharp rest energy in each copy (fixed); read as particle content, the graded
rung is one particle (interpretive). A two-particle reading exists, and
it is fixed as an algebra: the commutant of the first copy is a second, commuting $\Cl(3,3)$, whose generators are the graded ones
dressed by the chirality of the first copy (fixed). But it needs three further structures, and none of them is selected. They are
a common complex structure, which exists only on a $32$-dimensional subspace, an additive Hamiltonian that is not a vector of
$\Cl(6,6)$, and an ungraded exchange that is not the Clifford exchange (constructed; the selection of a complex structure is obstructed
under the diagonal action). The Fock number of the tower counts quanta, not particles. No gauge-admissible field that commutes with the
binding root changes it (fixed). It is not invariant under the interior Lorentz algebra, under the boosts of $\lG$, or under the
charge $J$ (fixed), so as a particle number it is an internal reading (interpretive). Identifiers e1-$n$ refer to the \texttt{PASS} lines of
\texttt{e1\_two\_from\_one.py} (\cref{e1:tab:claims}).

\paragraph{Notation.} $S=\R^8$ with the realisation of equation~\cref{eq:realisation}, $\omega=e_1\cdots e_6$, $J=e_2e_3$, $T=e_1e_2e_3$,
$\Omega=e_1e_4e_5e_6$, and $\gamma^\mu=(e_1,e_4,e_5,e_6)$. \cref{app:realisation} records $\omega=J\Omega=\Gamma$, the chirality. On
$S_2=S\otimes S=\R^{64}$ the graded generators are $E_{(1,a)}=e_a\otimes\id$ and $E_{(2,a)}=\omega\otimes e_a$ (equation~\cref{eq:tower}),
and the \emph{ungraded} copy-2 generators are $u_a=\id\otimes e_a$. We write $\Gamma^{(1)}=\omega\otimes\id$ and $\Gamma^{(2)}=\id\otimes\omega$
for the volume elements of the two constituents, $J^{(1)}=J\otimes\id$, $J^{(2)}=\id\otimes J$, $T^{(1)}=T\otimes\id$, $T^{(2)}=\id\otimes T$,
and $\omega_S=E_{(1,2)}E_{(1,3)}E_{(2,2)}E_{(2,3)}$ for the volume element of the four surplus directions. On vectors, $K_n$ is the
para-complex structure $Ke_i=e_{i+3}$, $Ke_{i+3}=e_i$ on each constituent; it is $-\cK$ of \cref{app:realisation}. $G_n$ is the binding
root, and $\lG$ and $\lint$ are the field-strength and interior Lorentz algebras of one constituent. The Fock modes are
$a^\pm_{k,i}=\half(E_{(k,i)}\pm E_{(k,i+3)})$, with $N=\sum_{k,i}a^+_{k,i}a^-_{k,i}$ and $N^{(k)}$ the part of constituent $k$
(\cref{d5:prop:fock}). The rest masses are $m_1,m_2\ge0$ and $M=(m_1^2+m_2^2)^{1/2}$. The diagonal interior Lorentz generators are
$\Delta_{\mu\nu}=\half(E_{(1,\mu)}E_{(1,\nu)}+E_{(2,\mu)}E_{(2,\nu)})$, and $M_{12}=\half\sum_\mu\eta_{\mu\mu}E_{(1,\mu)}E_{(2,\mu)}$ is the interior copy rotation.
The copy rotation is $C_b=\half\sum_a\eta_{aa}E_{(1,a)}E_{(2,a)}$, with the sign of \cref{d5:prop:cmrel}(ii); this sign is used in every
section of this chapter. $k_I\in\{1,2\}$ is the number of constituents on which $\lint$ acts (the spectator and the diagonal action;
\cref{thm:algebra}).

% --- source: emergence/e1_two_from_one#2 (Two copies of (3,3)Cl(3,3) inside (6,6)Cl(6,6))
\subsection{Two copies of \texorpdfstring{$\Cl(3,3)$}{Cl(3,3)} inside \texorpdfstring{$\Cl(6,6)$}{Cl(6,6)}}
\label{e1:sec:copies}

\begin{proposition}[The commuting copy]
\label{e1:prop:copies}
\begin{enumerate}[label=(\roman*),leftmargin=2.2em]
\item The twelve $E_{(k,a)}$ satisfy the Clifford relations of $\R^{6,6}$. The six $u_a$ satisfy those of $\R^{3,3}$ and commute with every
$E_{(1,b)}$. Moreover $u_a=\Gamma^{(1)}E_{(2,a)}=E_{(1,1)}\cdots E_{(1,6)}E_{(2,a)}$, a blade of grade $7$.
\item The commutant of $\Cl(3,3)\otimes\id$ in $M_{64}(\R)$ is spanned by $64$ blades of $\Cl(6,6)$, each equal to $\pm\id\otimes e_J$; it is
$\id\otimes M_8(\R)$, the algebra generated by the $u_a$. A second copy of $\Cl(3,3)$ commuting with the first is therefore unique up to
conjugation by $\id\otimes GL(8,\R)$.
\item The dressing is the chirality and the Fock parity of the first copy: $\Gamma^{(1)}=(J\Omega)\otimes\id=-(-1)^{N^{(1)}}$.
\item Both the ungraded product $\Cl(3,3)\otimes\Cl(3,3)$ and the graded product $\Cl(3,3)\hat\otimes\Cl(3,3)=\Cl(6,6)$ are $M_{64}(\R)$, with
the same $4096$ blades as basis. The twelve elements $e_a\otimes\id$, $u_a$ span no quadratic space: $(e_1\otimes\id+u_1)^2=2\,\id+2\,e_1\otimes e_1$.
\item At rung $n$ the ungraded generator of constituent $k$, $\id^{\otimes(k-1)}\otimes e_a\otimes\id^{\otimes(n-k)}$, equals
$\Gamma^{(1)}\cdots\Gamma^{(k-1)}E_{(k,a)}$, a blade of grade $6(k-1)+1$. Ungraded generators of distinct constituents commute and graded
ones anticommute (checked at $n=3,4$).
\end{enumerate}
\end{proposition}

\begin{proof}
(i) $(\omega\otimes\id)(\omega\otimes e_a)=\omega^2\otimes e_a=\id\otimes e_a$, and $\Gamma^{(1)}=E_{(1,1)}\cdots E_{(1,6)}$. (ii) Conjugation by an
invertible blade multiplies every blade by $\pm1$ (equation~\cref{eq:bladecomm}), so the commutant of a set of blades is spanned by the
blades it contains. A blade commutes with all $E_{(1,b)}$ exactly when it is $\Gamma^{(1)\varepsilon}$ times a blade of copy~2 with grade
parity $\varepsilon$, that is $\pm\id\otimes e_J$ (\cref{d5:prop:graded}(iii)). The uniqueness is the Skolem--Noether theorem in
$\id\otimes M_8(\R)\cong M_8(\R)$. (iii) $\omega=J\Omega$ by moving $e_1$ past $e_2e_3$, and $(-1)^{N_1+N_2+N_3}=-\omega$ is \cref{d5:prop:fock}. (iv) Distinct blades $B,B'$ have $\Tr(B^{\!\top}B')=0$ because $B^{\!\top}B'$ is $\pm$ a blade other than $\id$. The square is
a one-line expansion. (v) $\Gamma^{(j)}=\id^{\otimes(j-1)}\otimes\omega\otimes\id^{\otimes(n-j)}$, and the $\omega$'s cancel in pairs.
\end{proof}

Parts (i)--(iii) restate \cref{d5:prop:graded} and~\cref{d5:prop:fock} (inherited there) and \cref{app:realisation}; naming
$u_a=\Gamma^{(1)}E_{(2,a)}$ as the commuting copy is a restatement in this notation (e1-1a to e1-1d). The
Klein--Jordan--Wigner reading is inherited \cite{JordanWigner28,Klein38}, and so is the exponential law \cite{Chevalley54,LawsonMichelsohn89}. Two
remarks organise the rest of the section. First, bivectors do not see the grading: $E_{(2,a)}E_{(2,b)}=\id\otimes e_ae_b$. Hence each
constituent's $\lint^{(k)}$ and $\lG^{(k)}$, its charge $J^{(k)}=E_{(k,2)}E_{(k,3)}$ and its number $N^{(k)}$ are the same operators in both
readings, and the two Lorentz algebras commute in both. Second, the readings differ only in their odd elements, the gamma matrices and
the mass terms. The algebra $M_{64}(\R)$ is the same; what changes is which of its elements are taken as the vectors.

% --- source: emergence/e1_two_from_one#3 (Rest energies: quadrature against addition)
\subsection{Rest energies: quadrature against addition}
\label{e1:sec:rest}

Put $H_g=m_1E_{(1,1)}+m_2E_{(2,1)}$, the rest Hamiltonian built from the tower's vectors, and $H_u=m_1\,e_1\otimes\id+m_2\,\id\otimes e_1$,
built from the commuting copy. On one constituent the rest flow of the real Dirac form is $\dot\Psi=-JH(0)\Psi$ with $JH(0)=mT$, of angular
frequency $m$ (\cref{sec:firstorder}; \cref{sec:resttworot}).

\begin{proposition}[Graded and ungraded rest energies]
\label{e1:prop:rest}
\begin{enumerate}[label=(\roman*),leftmargin=2.2em]
\item $H_g^2=M^2\,\id$ and $\Tr H_g=0$, so $\spec H_g=\{\pm M\}$, each of multiplicity $32$ for $M>0$ (coincident eigenvalues add their
multiplicities: at $M=0$, $H_g=0$ with the single eigenvalue $0$ on $64$ dimensions).
\item With $B=E_{(2,1)}E_{(1,1)}$, $B^2=-\id$, and $U_\phi=\cos\phi+\sin\phi\,B$, $\tan2\phi=m_2/m_1$: $U_\phi(ME_{(1,1)})U_\phi^{-1}=H_g$. The graded
rest Hamiltonian is one Dirac mass along a unit vector of $\R^{6,6}$, turned by a compact element of $\mathrm{Spin}(6,6)$.
\item $H_u$ has eigenvalue $s_1m_1+s_2m_2$ on the $16$-dimensional joint eigenspace of each sign pair $(s_1,s_2)$, so
$\spec H_u=\{\pm m_1\pm m_2\}$; for generic masses each value has multiplicity $16$, and at $m_1=m_2$ the value $0$ has multiplicity $32$.
$H_u=m_1E_{(1,1)}+m_2\,\Gamma^{(1)}E_{(2,1)}$ has components of grade $1$ and $7$. For $m_1,m_2\neq0$ it is diagonalisable with at least three
distinct eigenvalues, so it is not conjugate to any vector of $\R^{6,6}$: a vector satisfies $v^2=\eta_2(v,v)\id$ and is either
diagonalisable with at most two eigenvalues or nilpotent.
\item Joint mass shells. No nonzero $\psi$ satisfies $E_{(1,1)}\psi=s_1\psi$ and $E_{(2,1)}\psi=s_2\psi$ with $s_1,s_2=\pm1$. Each joint sign
$(s_1,s_2)$ of $e_1\otimes\id$ and $\id\otimes e_1$ occurs on $16$ dimensions. More generally, the graded Dirac operators
$D_k=\sum_\mu p^{(k)}_\mu E_{(k,\mu)}$ anticommute, with $(D_1+D_2)^2=(p^{(1)}\cdot p^{(1)}+p^{(2)}\cdot p^{(2)})\id$, while their ungraded
counterparts commute.
\item Rest flows. $J^{(1)}$ and $J^{(2)}$ commute with $H_g$ and $H_u$. The ungraded generator $Z_u=m_1T^{(1)}+m_2T^{(2)}$ is a sum of two commuting
complex structures, with eigenvalues $i(\pm m_1\pm m_2)$: $e^{-\lambda Z_u}$ winds a two-torus with angular frequencies $m_1$ and $m_2$. The
graded generator $Z_g=J^{(1)}H_g$ has $Z_g^2=-M^2\id$, so $e^{-\lambda Z_g}=\cos(M\lambda)-\sin(M\lambda)Z_g/M$ has the single frequency $M$ (at $M=0$,
$Z_g=0$ and the flow is the identity, the limit of the formula).
\item At rung $n$ the graded $E_{(k,1)}$ anticommute and the ungraded ones commute, with traceless products over every nonempty subset. Hence
$\spec\bigl(\sum_km_kE_{(k,1)}\bigr)=\{\pm(\sum_km_k^2)^{1/2}\}$, while the ungraded sum has spectrum $\{\sum_k\pm m_k\}$, each sign
pattern on $4^n$ dimensions (coinciding values add their multiplicities).
\end{enumerate}
\end{proposition}

\begin{proof}
(i) $E_{(1,1)}$ and $E_{(2,1)}$ are anticommuting traceless involutions. (ii) $U_\phi XU_\phi^{-1}=\cos^2\!\phi\,X+\sin\phi\cos\phi\,[B,X]-\sin^2\!\phi\,BXB$, with
$[B,E_{(1,1)}]=2E_{(2,1)}$ and $BE_{(1,1)}B=E_{(1,1)}$. (iii) Commuting involutions with traceless products have all four joint signs with equal
multiplicity. (iv) If $E_{(1,1)}\psi=s_1\psi$ and $E_{(2,1)}\psi=s_2\psi$, then $s_1s_2\psi=E_{(1,1)}E_{(2,1)}\psi=-E_{(2,1)}E_{(1,1)}\psi=-s_1s_2\psi$. Generators of
distinct constituents anticommute, so the cross terms of $(D_1+D_2)^2$ cancel. (v) $T^{(k)}=J^{(k)}\times(\text{the mass term of copy }k)$, both
$T^{(k)}$ square to $-\id$, and they commute. $J^{(1)}$ commutes with $E_{(1,1)}$ and $E_{(2,1)}$ and squares to $-\id$, so
$(J^{(1)}H_g)^2=-H_g^2$. (vi) As (i) and (iii), with $2^n$ joint signs.
\end{proof}

Parts (i) and (iii) at $m_1=m_2$ restate \cref{d6:prop:dirac}(iv). The general masses, the rotation (ii), the joint mass shells
(iv), the flows (v) and the ladder (vi) are fixed (e1-2a to e1-2g). The free-particle facts behind (iv) are inherited. A system of two
particles has two mass-shell conditions, one per particle, imposed simultaneously, as in the two-body Dirac (Breit) equation
\cite{Breit29}. A single particle in a doubled momentum space has one, $p^{(1)}\cdot p^{(1)}+p^{(2)}\cdot p^{(2)}=M^2$. By (ii) and (iv)
the graded tower is of the second kind. Its two copies admit no simultaneous nonzero-mass Dirac constraints of the graded form, its rest energy is one value $M$, and the second copy
enters the one particle as further momentum and mass directions. The ungraded operators are of the first kind: the energies add, and
the pair threshold is $m_1+m_2$.

% --- source: emergence/e1_two_from_one#4 (Complex structures: the two-particle space and the selected )
\subsection{Complex structures: the two-particle space and the selected unit}
\label{e1:sec:complex}

The paper reads $J=e_2e_3$ as the imaginary unit of one constituent (\cref{thm:i}). A complex two-particle space needs one
imaginary unit acting as $J$ on both factors.

\begin{proposition}[The complex two-particle space]
\label{e1:prop:W}
\begin{enumerate}[label=(\roman*),leftmargin=2.2em]
\item $J^{(1)}$ and $J^{(2)}=E_{(2,2)}E_{(2,3)}$ commute and square to $-\id$, and $J^{(1)}J^{(2)}=\omega_S$ is a symmetric traceless involution. The
subspace $W=\{\psi:J^{(1)}\psi=J^{(2)}\psi\}=\ker(\omega_S+\id)$ has dimension $32$, and its orthogonal complement is
$\overline W=\ker(J^{(1)}+J^{(2)})=\ker(\omega_S-\id)=\operatorname{im}(J^{(1)}-J^{(2)})$. With $J:=J^{(1)}|_W=J^{(2)}|_W$ as imaginary unit,
$W\cong S\otimes_\C S$, and $\overline W\cong S\otimes_\C\overline S$.
\item A blade of $\Cl(6,6)$ with an even number of surplus indices preserves $W$ and $\overline W$. A blade with an odd number exchanges them,
for instance $\id\otimes e_2$ (charge conjugation $C=e_2$ of copy~2 alone, \cref{thm:C}) and $E_{(1,1)}E_{(2,2)}$.
\item $J^{(1)}$ and $J^{(2)}$ commute with $\lint^{(1)}$, $\lint^{(2)}$, the eight graded and the eight ungraded interior gammas, $H_g$, $H_u$ and the
interior copy rotation $M_{12}=\half\sum_\mu\eta_{\mu\mu}E_{(1,\mu)}E_{(2,\mu)}$. All of these are $\C$-linear on $W$ and on $\overline W$.
\item On $W$ the surplus part of the copy rotation vanishes, since $E_{(1,3)}E_{(2,3)}=E_{(1,2)}E_{(2,2)}\omega_S$, so the copy rotation
$C_b=\half\sum_a\eta_{aa}E_{(1,a)}E_{(2,a)}$ of \cref{d5:prop:cmrel}(ii) acts on $W$ as $M_{12}$, $\C$-linearly. The four cross surplus bivectors $E_{(1,s)}E_{(2,t)}$
($s,t\in\{2,3\}$) preserve $W$ but anticommute with $J^{(1)}$ and $J^{(2)}$, so they act on $W$ antilinearly.
\end{enumerate}
\end{proposition}

\begin{proof}
(i) Bivectors of distinct constituents commute; $(J^{(1)}J^{(2)})^2=(J^{(1)})^2(J^{(2)})^2=\id$ and $\Tr(J\otimes J)=(\Tr J)^2=0$. $J^{(1)}\psi=J^{(2)}\psi$ is
equivalent to $J^{(1)}J^{(2)}\psi=-\psi$. $J^{(1)}-J^{(2)}$ is antisymmetric, so its kernel is the orthogonal complement of its image, and
$\omega_S(J^{(1)}-J^{(2)})=J^{(1)}-J^{(2)}$. The canonical map $S\otimes_\R S\to S\otimes_\C S$ has kernel $\operatorname{im}(J^{(1)}-J^{(2)})$, spanned by
$Jx\otimes y-x\otimes Jy$, so it restricts to an isomorphism on $W$; the case of $\overline W$ is the same with $-J$ on the second factor. (ii)
By equation~\cref{eq:bladecomm} a blade commutes with the surplus four-blade $\omega_S$ iff it shares an even number of indices with it, and
otherwise anticommutes. (iii) By the blade rule: the graded gammas, the bivectors of $\lint^{(k)}$ and the terms of $M_{12}$ contain no
surplus index; the ungraded gammas $u_\mu=\Gamma^{(1)}E_{(2,\mu)}$ contain both surplus indices of copy~1, hence meet $J^{(1)}$ evenly and
$J^{(2)}$ not at all. Directly, $\id\otimes\gamma$ commutes with $J\otimes\id$ and $\id\otimes J$. (iv) $E_{(1,2)}E_{(2,2)}\omega_S=-E_{(2,2)}E_{(1,3)}E_{(2,2)}E_{(2,3)}=E_{(1,3)}E_{(2,3)}$, by the anticommutation of distinct
generators. On $W$, $\omega_S=-\id$, so the two surplus terms of $C_b$ cancel. A cross surplus bivector shares one index
with each $J^{(k)}$.
\end{proof}

The tensor-product linear algebra of (i) is inherited. Parts (ii)--(iv) are fixed (e1-3a, e1-3b, e1-3f). With the external $i$ of
\cref{d6:prop:dirac}(i), $J^{(1)}+J^{(2)}$ is $2J$ on $W$ and $0$ on $\overline W$. So $W$ carries total charge $\pm2$, two particles
of equal charge, and $\overline W$ the neutral pairs, among them the particle--antiparticle state $\psi\otimes C\phi$ of that proposition
(restated; the charge reading interpretive). The question is now whether anything in the algebra selects $W$ and the complex
structure on it.

\begin{proposition}[Which imaginary unit the second rung selects]
\label{e1:prop:select}
\begin{enumerate}[label=(\roman*),leftmargin=2.2em]
\item Per-particle selection. The commutant of the eight ungraded gammas $\gamma^\mu\otimes\id$, $\id\otimes\gamma^\mu$ is $\HH\otimes\HH\cong M_4(\R)$,
spanned by $16$ blades. Adding the two temporal volume elements $T^{(1)},T^{(2)}$ (the datum of \cref{rem:ambient} for each copy)
cuts it to $\operatorname{span}\{\id,J^{(1)},J^{(2)},\omega_S\}\cong\C\otimes_\R\C\cong\C\oplus\C$. Its solutions of $x^2=-\id$ are exactly
$\pm J^{(1)}$ and $\pm J^{(2)}$, four isolated points, none of which acts as $J$ on both factors of $S_2$.
\item The spectator action, $k_I=1$. The selected unit of \cref{thm:i} at $n=2$, the product of the eight internal generators, is
$J_{\rm sel}=J\otimes\omega=J^{(1)}\Gamma^{(2)}$ (\cref{d5:prop:vacuum}(iv)). It commutes with $J^{(1)}$ and $J^{(2)}$ and so preserves
$W$. It equals $J^{(1)}$ on $\{\Gamma^{(2)}=1\}$ and $-J^{(1)}$ on $\{\Gamma^{(2)}=-1\}$, and $J^{(2)}$ on $\{\Gamma^{(2)}\omega_S=-1\}$ and
$-J^{(2)}$ on $\{\Gamma^{(2)}\omega_S=1\}$, $32$ dimensions each; on $W$ it is $J$ on $16$ and $-J$ on $16$ dimensions. It anticommutes with every generator of copy~2, graded or ungraded. So no gamma matrix of copy~2 is $J_{\rm sel}$-linear,
and neither $H_g$ nor $H_u$ is for $m_2\neq0$; $m_1E_{(1,1)}$ is.
\item The diagonal action, $k_I=2$. The commutant of the eight graded interior gammas $E_{(k,\mu)}$ is spanned by $16$ blades ($M_2(\HH)$). Adding
the six surplus bivectors cuts it to $\operatorname{span}\{\id,\omega_S\}\cong\R\oplus\R$, with $\omega_S^2=+\id$. No complex structure commutes with
the gammas of both copies and the surplus bivectors. On $W$ alone the restrictions of the same operators have commutant $\R$.
\item Up the ladder, under the diagonal action of rung $n$ the same commutant is $\operatorname{span}\{\id,\omega_S\}$ with
$\omega_S=J^{(1)}\cdots J^{(n)}$ and $\omega_S^2=(-1)^n$ (checked at $n=2,3$ for the commutant and $n=2,3,4$ for the square). An imaginary unit
is selected exactly at odd $n$. The common subspace $J^{(1)}=\cdots=J^{(n)}$, the complex $n$-fold product $S\otimes_\C\cdots\otimes_\C S$, has
dimension $2\cdot4^n$.
\end{enumerate}
\end{proposition}

\begin{proof}
(i) By the blade argument of \cref{e1:prop:copies}(ii), since $\id\otimes\gamma^\mu=\Gamma^{(1)}E_{(2,\mu)}$ and
$T^{(2)}=\Gamma^{(1)}E_{(2,1)}E_{(2,2)}E_{(2,3)}$ are blades. Equivalently, the commutant of the gammas on $S$ is the quaternion algebra of \cref{rem:ambient}, and that of $\{\gamma^\mu,T\}$ is $\R\oplus\R J$. For $x=a+bJ^{(1)}+cJ^{(2)}+d\,\omega_S$,
$x^2=(a^2-b^2-c^2+d^2)+2(ab-cd)J^{(1)}+2(ac-bd)J^{(2)}+2(ad+bc)\omega_S$, and $x^2=-\id$ has exactly the four solutions listed. (ii) $J_{\rm sel}$ is a blade of
grade $8$ containing every copy-2 index, so it anticommutes with each $E_{(2,a)}$; it commutes with $\Gamma^{(1)}$, so it also
anticommutes with each $u_a=\Gamma^{(1)}E_{(2,a)}$. $\Gamma^{(2)}$ commutes with
$J^{(1)}$ and $J^{(2)}$, and $\Tr\Gamma^{(2)}=0$, also on $W$. Finally $J_{\rm sel}\mp J^{(2)}=J^{(1)}(\Gamma^{(2)}\pm\omega_S)$ with $J^{(1)}$
invertible, and $\Gamma^{(2)}\omega_S$ is a traceless involution. (iii) A blade commuting with all eight interior generators meets them in all or none
of them. Meeting all of them requires odd grade, and meeting none requires even grade. Commuting with every surplus bivector forces it to meet
the four surplus directions in all or none. Both conditions together leave $\id$ and $\omega_S$. The algebra generated contains the
central idempotent $\half(\id-\omega_S)$, so its commutant on $W$ is the compression of $\operatorname{span}\{\id,\omega_S\}$, which is $\R$.
(iv) The same parity count with $4n$ interior and $2n$ surplus generators: a blade containing all $4n$ interior generators has grade $4n$ or
$6n$, both even, so it fails. The square follows from equation~\cref{eq:transpose}, and the dimension is $8^n/2^{n-1}$, because the $n-1$
involutions $J^{(k)}J^{(k+1)}$ commute and have traceless products.
\end{proof}

Part (i) is fixed, with the quaternion commutant of one copy restated from \cref{rem:ambient}. Part (ii) restates \cref{thm:i} and \cref{d5:prop:vacuum}(iv) (fixed there), and the comparison with $W$ is fixed. Part (iii) is obstructed under the
stated hypotheses (e1-3c to e1-3e, e1-3g). Part (iv) restates the selection at odd $n$ of \cref{prop:diagtower}(ii) (fixed there)
and adds that nothing is selected at even $n$ (fixed; e1-3h).

\begin{corollary}[No selected two-particle complex structure]
\label{e1:cor:nocx}
At the second rung none of the selection principles of \cref{ch:ch02_lorentz} gives a complex structure under which both copies are complex Dirac spinors
with $\C$-linear gamma matrices. Per copy they give $\pm J^{(1)}$ and $\pm J^{(2)}$ separately. The spectator action gives
$J_{\rm sel}$, under which copy~2 carries no linear Dirac operator, so that copy~2 is internal. The diagonal action gives none. What the
diagonal action does select is the splitting $S_2=W\oplus\overline W$, the eigenspaces of the one nontrivial element $\omega_S$ of its commutant.
The complex structure $J|_W$ that makes $W$ the two-particle space commutes with every interior gamma matrix, with both Hamiltonians and with the copy
rotation, but not with the surplus $\so4$. It is an added datum.
\end{corollary}

The corollary is fixed as a statement about these selection principles (\cref{e1:prop:select}). That the choice of $W$, rather than
$\overline W$, and of $J|_W$ is what turns the second rung into a two-particle space is interpretive. This also bears on the phrase
``two particles from one''. The paper's own second quantisation of one constituent (\cref{thm:spectra}, \cref{thm:subfock}) already has two-excitation states of one Dirac momentum (four fermionic modes), a level of dimension $6$ with charges $q=\pm2$ once each and $q=0$ four times.
In ordinary physics that is where several particles live (inherited \cite{Fock32}). It does not need the second rung.

% --- source: emergence/e1_two_from_one#5 (The Fock reading: quanta, not particles)
\subsection{The Fock reading: quanta, not particles}
\label{e1:sec:fock}

The tower has its own number operator. $S_n$ is the Fock space of $3n$ fermionic modes, and $N$ counts them (\cref{d5:prop:fock}). We ask which elements of $\so{3n,3n}$, the bivectors, conserve $N$.

\begin{proposition}[Bivectors and the number of quanta]
\label{e1:prop:Nbiv}
\begin{enumerate}[label=(\roman*),leftmargin=2.2em]
\item $N=\tfrac{3n}2-\half\sum_{k,i}E_{(k,i)}E_{(k,i+3)}$ is diagonal, with spectrum $0,\dots,3n$; at $n=2$ the multiplicities are
$1,6,15,20,15,6,1$. $N=N^{(1)}+N^{(2)}$, the same operator in the graded and ungraded readings. $(-1)^N=\Gamma^{(1)}\Gamma^{(2)}$ is the volume element of
$\Cl(6,6)$.
\item On vectors $[N,v]=K_nv$. A bivector $X$ commutes with $N$ iff $\ad X$ commutes with $K_n$. These form
$\mathfrak{gl}(L_+)\cong\mathfrak{gl}(3n,\R)$, of dimension $9n^2$ ($36$ at $n=2$), the stabiliser of the transverse maximal isotropic subspaces
$L_\pm=\ker(K_n\mp\id)$. The remaining $3n(3n-1)$ directions ($30$ at $n=2$) are $\Lambda^2L_+\oplus\Lambda^2L_-$, on which $[N,X]=\pm2X$: they create
or annihilate a pair of quanta.
\item $(-1)^N$ commutes with every bivector and anticommutes with every generator. No element of $\so{3n,3n}$ changes $N$ by an odd number.
\end{enumerate}
\end{proposition}

\begin{proof}
(i) $N_i=\half(\id-e_ie_{i+3})$, and the pole operators are diagonal (\cref{app:realisation}). (ii) $[N,a^\pm_{k,i}]=\pm a^\pm_{k,i}$, and
$a^\pm_{k,i}$ is proportional to $e_i\pm e_{i+3}$ in constituent $k$, the $\pm1$ eigenvectors of $K$. $N-\tfrac{3n}2$ is a bivector, so $[N,X]$
is one, and the spin representation of $\so{V_n}$ is faithful. Hence $[N,X]=0$ iff $[\ad N,\ad X]=[K_n,\ad X]=0$. Under $\ad K_n$,
$\so{V_n}\cong\Lambda^2V_n=\Lambda^2L_+\oplus(L_+\otimes L_-)\oplus\Lambda^2L_-$ has eigenvalues $2,0,-2$, and $L_\pm$ are isotropic because
$\eta(e_i\pm e_{i+3},e_j\pm e_{j+3})=\delta_{ij}-\delta_{ij}$. (iii) Grade parity.
\end{proof}

Part (i) restates \cref{d5:prop:fock}. The rest is the standard structure of the spin representation on the exterior algebra of
a maximal isotropic subspace \cite{Chevalley54}, inherited, with its realisation in the tower fixed (e1-4a, e1-4b, e1-4g). The question
is which of these directions are physical.

\begin{theorem}[Compact fields conserve the quanta]
\label{e1:thm:Nrule}
Let $X\in\Cent(G_n)\cap\so{3n,3n}$, the algebra $\so{3n,\C}$ of binding-compatible bivectors (the case $n=3$ is \cref{thm:so9C}, the cases $n\le3$
\cref{d5:prop:so6C}(i); the block form of the proof gives every $n$). Then $K_nXK_n=-X^{\!\top}$, the Cartan involution. Consequently:
\begin{enumerate}[label=(\roman*),leftmargin=2.2em]
\item $X$ commutes with $N$ iff $X$ is compact (antisymmetric), that is iff it is a paired rotation $\bigl(\begin{smallmatrix}R&0\\0&R\end{smallmatrix}\bigr)$.
These form $\so{3n}$, of dimension $\tfrac32n(3n-1)$.
\item The non-compact part, the paired boosts $\bigl(\begin{smallmatrix}0&S\\-S&0\end{smallmatrix}\bigr)$, anticommutes with $K_n$. A nonzero paired
boost has nonzero components in both $\Lambda^2L_+$ and $\Lambda^2L_-$, so it changes $N$ by $+2$ and by $-2$: it has no component commuting
with $N$, and $[N,X]\neq0$ for $X\neq0$. It may still conserve partial occupation numbers; for instance $(e_3e_5-e_2e_6)/2$ conserves $N_1$ and
$N_2-N_3$.
\item By \cref{cor:noboost} a boost generates no compact $U(1)$, so it is not gauge-admissible in the sense fixed after that corollary.
Hence no gauge-admissible, binding-compatible field creates or annihilates a pair of quanta, at any rung.
\end{enumerate}
\end{theorem}

\begin{proof}
On vectors $\eta G_n=-K_n$ and $G_n\eta=-\eta G_n$. For $X$ commuting with $G_n$,
$K_nXK_n=\eta G_nX\eta G_n=-\eta G_nXG_n\eta=-\eta XG_n^2\eta=\eta X\eta$, and $\eta X\eta=-X^{\!\top}$ for $X\in\so{\eta}$. (i) follows from \cref{e1:prop:Nbiv}(ii). The block
forms are those of \cref{d5:prop:so6C}(i). (ii) $X$ maps $L_\pm$ to $L_\mp$, so its components lie in $\Lambda^2L_+\oplus\Lambda^2L_-$.
$G_n$ commutes with $X$, maps $L_+$ onto $L_-$ (because $G_nK_n=-K_nG_n$) and normalises $\so{\eta}$, so conjugation by $G_n$ carries the $\Lambda^2L_+$
component of $X$ to its $\Lambda^2L_-$ component. They vanish together. (iii) Gauge admissibility is compactness, the criterion of
\cref{sec:sigrule}.
\end{proof}

The theorem is fixed. It was checked exactly on vectors at $n=1,\dots,4$, where $\dim\Cent(K_n)\cap\so{3n,3n}=9,36,81,144$,
$\dim\Cent(G_n)\cap\so{3n,3n}=6,30,72,132$ and the compact part has half of it. It was checked entry by entry on $S_n$ at $n=1,\dots,4$ for every
paired rotation $e_ie_j-e_{i+3}e_{j+3}$ and every paired boost $e_ie_{j+3}-e_je_{i+3}$, within one constituent and across two (e1-4c,
e1-4d). At $n=4$, $S_4$ has dimension $4096$, and the generators were handled as signed permutations.

\begin{proposition}[The quanta and the Lorentz algebras]
\label{e1:prop:Nlorentz}
The number $N$ behaves as follows under the copy algebra and under the two Lorentz algebras of one constituent.
\begin{enumerate}[label=(\roman*),leftmargin=2.2em]
\item The copy rotation $C_b$ commutes with the total number $N$, but it does not commute with the number $N^{(1)}$ of the first
constituent. The quarter turn
\[
U=e^{\pi C_b/2}=\tfrac18\textstyle\prod_a(\id+\eta_{aa}E_{(1,a)}E_{(2,a)})
\]
satisfies $UN^{(1)}U^{-1}=N^{(2)}$. The spin lift of the copy boost $\varepsilon\otimes G$ changes $N$ by $\pm2$ only, with both signs.
\item In $\lint$ the elements commuting with $N$ form $\R\,e_1e_4$, a boost (the pole operator $\id-2N_1$). The rotations $e_4e_5,e_4e_6,e_5e_6$, the
boosts $e_1e_5,e_1e_6$ and the charge $J=e_2e_3$ do not commute with $N$.
\item In $\lG$ the elements commuting with $N$ form its compact $\so3$, the paired rotations. No nonzero boost of $\lG$ commutes with $N$.
\end{enumerate}
\end{proposition}

\begin{proof}
(i) $C_b$ is a paired rotation (\cref{e1:thm:Nrule}(i)). Its six terms commute and square to $-\tfrac14$, so
$e^{\theta C_b}=\prod_a(\cos\tfrac\theta2+\sin\tfrac\theta2\,\eta_{aa}E_{(1,a)}E_{(2,a)})$, which at $\theta=\tfrac\pi2$ is $U$; it carries $E_{(1,a)}$ to
$\pm E_{(2,a)}$ (\cref{d5:prop:cmrel}(ii)). The copy
boost is a paired boost (\cref{d5:prop:so6C}(iii)); apply \cref{e1:thm:Nrule}(ii). (ii) $\lint\cap\mathfrak{gl}(3,\R)$ is computed: of
the six planes in $\langle e_1,e_4,e_5,e_6\rangle$ only $\langle e_1,e_4\rangle$ is $K$-invariant, and $J$ involves the modes $a^+_2a^+_3$. (iii) is
\cref{e1:thm:Nrule} at $n=1$, with $\lG=\Cent(G)\cap\so{3,3}$ (\cref{thm:GLorentz}).
\end{proof}

All three parts are fixed (e1-4e, e1-4f), and (i) gives a precise sense to the copy rotation as a hop of quanta between the constituents
at fixed total number. Part (ii) corrects an easy expectation. The Fock number is not invariant under the spatial rotations of the
interior four-plane, only under one of its boosts, and it does not commute with the charge $J$. The quanta counted by $N$ are therefore
not Lorentz-invariant particles. Within $\lint$ they are invariant only under the one boost $e_1e_4$, the pole operator of the first mode;
within $\lG$ only under its compact $\so3$; they are not invariant under the charge $J$ or under any boost of $\lG$. The bivectors that preserve
$N$ form $\mathfrak{gl}(3n,\R)$ (\cref{e1:prop:Nbiv}(ii)), which contains boosts (at $n=1$ six of its nine dimensions), so $N$ is tied to the
split structure $K$ and not to a Lorentz algebra. Reading $N$ as a particle number is therefore an internal reading (interpretive).
What \cref{e1:thm:Nrule} does establish is the following. If quanta are counted by $N$, the only elements of the binding-compatible algebra
that create them in pairs are the paired boosts. These are excluded as gauge fields, so within the binding-compatible algebra pair creation
of quanta would need a non-compact, non-gauge interaction (fixed, with the identification of quanta interpretive). Outside it, compact
gauge-admissible fields that do not commute with $G_n$, such as the charge $J$ and the time--time rotation $e_1e_2$, do create and annihilate
quanta in pairs: both have components with $\Delta N=+2$, $0$ and $-2$. This agrees with the
statement of \cref{d6:sec:dirac} that duplication in the tower's sense is not pair creation.

% --- source: emergence/e1_two_from_one#6 (Exchange and statistics)
\subsection{Exchange and statistics}
\label{e1:sec:exchange}

Two swaps act on $S_2$: the ungraded swap $P(x\otimes y)=y\otimes x$, and the Koszul swap $\tau_K=D\,P$ of \cref{d5:prop:cmrel}(iii), where
$D=(-1)^{\mathsf p\otimes\mathsf p}$ and $\omega=(-1)^{\mathsf p}$ on the basis of equation~\cref{eq:realisation}.

\begin{proposition}[Two exchanges]
\label{e1:prop:swap}
\begin{enumerate}[label=(\roman*),leftmargin=2.2em]
\item $P(e_a\otimes\id)P=u_a=\Gamma^{(1)}E_{(2,a)}$. So $P$ exchanges the ungraded copies and maps the vectors $E_{(1,a)}$ of $\Cl(6,6)$ to elements of
grade $7$. It commutes with the diagonal Lorentz generators $\Delta_{\mu\nu}$, with $N$, with $\omega_S$ (so it preserves $W$) and with $H_u$ at
$m_1=m_2$, but not with $H_g$.
\item $\tau_KE_{(1,a)}\tau_K=E_{(2,a)}$. $\tau_K$ commutes with $\Delta_{\mu\nu}$, $N$, $\omega_S$ and $H_g$ at $m_1=m_2$, but not with $H_u$.
\item $D=\half(\id+\Gamma^{(1)}+\Gamma^{(2)}-\Gamma^{(1)}\Gamma^{(2)})$. So $\tau_K=P$ except on the sector $\Gamma^{(1)}=\Gamma^{(2)}=-1$, where $\tau_K=-P$.
\end{enumerate}
\end{proposition}

\begin{proof}
(i) $P(A\otimes B)P=B\otimes A$. Diagonal bivectors, $N=N^{(1)}+N^{(2)}$, $\omega_S=J\otimes J$ and $H_u$ at equal masses are symmetric under it,
while $P(\omega\otimes e_1)P=e_1\otimes\omega$. (ii) The first identity restates \cref{d5:prop:cmrel}(iii). Hence
$\tau_K\Gamma^{(1)}\tau_K=\Gamma^{(2)}$ and $\tau_Ku_1\tau_K=e_1\otimes\omega\neq u_1$, while bivectors $E_{(1,a)}E_{(1,b)}$ go to $E_{(2,a)}E_{(2,b)}$. (iii) $D$ is
diagonal with entry $-1$ exactly where both $\omega$-eigenvalues are $-1$.
\end{proof}

Each swap is a symmetry of its own reading and not of the other: $P$ of the additive Hamiltonian, $\tau_K$ of the graded one (fixed; e1-5a,
e1-5b). For the Lorentz content let $C=-\sum_{\mu<\nu}\eta^{\mu\mu}\eta^{\nu\nu}\Delta_{\mu\nu}^2$ be the diagonal Casimir, normalised as in \cref{sec:dualpair}. It is $\tfrac32$ on one constituent, it satisfies $C(C-3)(C-4)=0$ on $S_2$, and it takes the values $0$, $3$ and $4$ on
$(0,0)$, $(\half,\half)$ and $(1,0)\oplus(0,1)$. On one constituent each boost is $J\Gamma$ times its dual rotation,
$e_1e_4=J\Gamma\,e_5e_6$ and cyclically, so $\Gamma=\pm1$ are the two Weyl halves (e1-5c).

\begin{proposition}[Statistics sectors]
\label{e1:prop:stat}
The dimensions of the eigenspaces of $C$ in the sectors of $P$ and $\tau_K$ are those of \cref{e1:tab:stat}. In particular:
\begin{enumerate}[label=(\roman*),leftmargin=2.2em]
\item On $W$, $P$ is $\C$-linear. Its $-1$ eigenspace, $\Lambda^2_\C S$, carries $2(0,0)\oplus(\half,\half)$, and its $+1$ eigenspace, $\mathrm{Sym}^2_\C S$,
carries $(1,0)\oplus(0,1)\oplus(\half,\half)$. In the chirality sectors, $\Gamma^{(1)}=\Gamma^{(2)}=\pm1$ gives spin one in $P=+1$ and spin zero in
$P=-1$, while the mixed sector ($16$ dimensions, $C=3$) splits $8+8$.
\item On $W$, $\tau_K$ is $\C$-linear, and its two eigenspaces have the same Casimir content, $C=0,3,4$ on $2,8,6$ dimensions. The spin-one part of
$\tau_K=+1$ lies in $\Gamma^{(1)}=\Gamma^{(2)}=+1$, and that of $\tau_K=-1$ in $\Gamma^{(1)}=\Gamma^{(2)}=-1$: the two sectors are mirror images.
\item On $\overline W$, $P$ is antilinear, $PJ^{(1)}=-J^{(1)}P$, and $J^{(1)}$ exchanges its two eigenspaces. They carry the same content.
\end{enumerate}
\end{proposition}

\begin{proof}
All entries are traces of products of commuting integer matrices: the projectors $\half(\id\pm X)$ for $X\in\{\omega_S,P,\tau_K,\Gamma^{(1)},\Gamma^{(2)}\}$ and
the spectral projectors of $C$ (e1-5d to e1-5f). (i) $PJ^{(1)}P=J^{(2)}$, which equals $J^{(1)}$ on $W$. (ii) Use \cref{e1:prop:swap}(iii). (iii)
$J^{(2)}=-J^{(1)}$ on $\overline W$, and $J^{(1)}$ commutes with $C$.
\end{proof}

\begin{table}[t]
\centering
\small
\begin{tabular}{@{}lccc@{}}
\toprule
sector & on $W=S\otimes_\C S$ ($32$) & on $\overline W=S\otimes_\C\overline S$ ($32$) & on $S_2$ ($64$)\\
\midrule
$P=-1$ & $0^4\ 3^8$ & $0^2\ 3^8\ 4^6$ & $0^6\ 3^{16}\ 4^6$\\
$P=+1$ & $3^8\ 4^{12}$ & $0^2\ 3^8\ 4^6$ & $0^2\ 3^{16}\ 4^{18}$\\
$\tau_K=-1$ & $0^2\ 3^8\ 4^6$ & $0^2\ 3^8\ 4^6$ & $0^4\ 3^{16}\ 4^{12}$\\
$\tau_K=+1$ & $0^2\ 3^8\ 4^6$ & $0^2\ 3^8\ 4^6$ & $0^4\ 3^{16}\ 4^{12}$\\
\bottomrule
\end{tabular}
\caption{Casimir content $c^d$ ($C=c$ on $d$ real dimensions) of the exchange sectors under the diagonal interior Lorentz algebra.
Checks e1-5d to e1-5f.}
\label{e1:tab:stat}
\end{table}

The decompositions of $\Lambda^2$ and $\mathrm{Sym}^2$ of a Dirac spinor are inherited \cite{WeinbergQFT}; their appearance on $W$, and only on $W$,
is fixed. The Koszul column restates \cref{d5:prop:spin}(iii) on $S_2$ (fixed there) and refines it on $W$. The two $\tau_K$
sectors are isomorphic as real representations, since the realification of $(1,0)$ is that of $(0,1)$, but they are mirror images once $J$ is
the imaginary unit (fixed). So the ungraded swap sorts the complex two-particle space by statistics in the textbook way. Identical fermions
are antisymmetric under exchange \cite{Pauli40}, and the antisymmetric square of a Dirac spinor has only scalar and vector content at a single
momentum mode (inherited \cite{WeinbergQFT}). The Clifford exchange $\tau_K$ does not separate symmetric from antisymmetric content. Its two
sectors have the same real Lorentz content and differ only in the chirality of their spin-one part, so its eigenvalue is a chirality label,
not a statistics label. Fermi statistics
for two copies is therefore a constructed requirement: antisymmetry under $P$ on $W$. The algebra supplies $\tau_K$, the exchange that
preserves the vectors of $\Cl(6,6)$, and $\tau_K$ is blind to it.

% --- source: emergence/e1_two_from_one#7 (Answer: what makes the second rung carry two particles)
\subsection{Answer: what makes the second rung carry two particles}
\label{e1:sec:answer}

\Cref{e1:tab:answer} collects the structures. Read together they answer the question in four sentences; the particle language is an
interpretation under the stated assumptions, as \cref{e8:sec:main} says. (1) As given by the tower, and read in the Dirac picture, the second
rung carries one particle. Its vectors $E_{(k,a)}$ give one Dirac clock along a unit vector of $\R^{6,6}$, of mass $M=(m_1^2+m_2^2)^{1/2}$, with
no simultaneous nonzero-mass Dirac constraint of the graded form on the copies. Under the selected unit $J_{\rm sel}$ copy~2 carries no linear Dirac operator (fixed;
\cref{e1:prop:rest,e1:prop:select}). Read as particle content, the rung is one particle with a larger internal space (interpretive). (2) The algebra also contains, uniquely up to
conjugation, a second commuting copy, the generators dressed by the chirality of the first. With it the copies have separate mass shells and
additive energies (fixed; \cref{e1:prop:copies,e1:prop:rest}). (3) Three further structures make that a two-particle system in the sense of
ordinary quantum mechanics: the subspace $W$ with the complex structure $J|_W$, the additive Hamiltonian $H_u$, and the ungraded exchange $P$
with antisymmetry. None of them is selected. $W$ is singled out only together with $\overline W$, and the diagonal action selects no imaginary unit at
even rungs (constructed; obstructed as a selection; \cref{e1:cor:nocx,e1:prop:stat}). (4) The tower's own counting operator $N$ counts quanta
of the split structure, not Lorentz-invariant particles. No gauge-admissible, binding-compatible field changes it (fixed;
\cref{e1:thm:Nrule,e1:prop:Nlorentz}).

\begin{table}[t]
\centering
\footnotesize
\begin{tabular}{@{}>{\raggedright\arraybackslash}p{3.6cm}>{\raggedright\arraybackslash}p{6.4cm}>{\raggedright\arraybackslash}p{4.6cm}@{}}
\toprule
structure & what it does at $n=2$ & class\\
\midrule
graded vectors $E_{(2,a)}$ & one Dirac clock of mass $M$; no simultaneous nonzero-mass graded Dirac constraints; $\tau_K$ as exchange & fixed (they define $\Cl(6,6)$)\\
commuting copy $u_a=\Gamma^{(1)}E_{(2,a)}$ & two mass shells, energies $\pm m_1\pm m_2$ & fixed as an algebra (unique up to $\id\otimes GL(8,\R)$); use as vectors constructed\\
common complex structure $J|_W$ & $W=S\otimes_\C S$, $32$ dimensions, total charge $\pm2$ & $W\oplus\overline W$ fixed by $\omega_S$; $J|_W$ constructed; selection obstructed ($k_I=2$), copy~2 internal ($k_I=1$)\\
additive Hamiltonian $H_u$ & threshold $m_1+m_2$; two-torus rest flow & constructed; not a vector (grades $1$ and $7$)\\
ungraded exchange $P$ & $\Lambda^2_\C S=2(0,0)\oplus(\half,\half)$ on $W$ & constructed; $\tau_K$ fixed but statistics-blind\\
Fock number $N$ & counts quanta; conserved by paired rotations, changed by $\pm2$ by paired boosts & fixed; as a particle number interpretive (not Lorentz- or $J$-invariant)\\
\bottomrule
\end{tabular}
\caption{Structures that bear on a two-particle reading of the second rung.}
\label{e1:tab:answer}
\end{table}

\paragraph{Which process takes one to two.} Three candidates have been tested. The step from rung $1$ to rung $2$ is a change of arena,
not a dynamical process. In the graded reading it enlarges the particle and does not add one. In the Fock reading the Lie algebra
$\so{3n,3n}$ changes $N$ only by $0$ or $\pm2$, and among binding-compatible fields only the paired boosts do it, which are not gauge fields.
A step from one quantum to two would need an odd element, which is a field operator and not a symmetry (fixed;
\cref{e1:prop:Nbiv,e1:thm:Nrule}). In the two-particle reading the step is the tensor embedding $\psi\mapsto\psi\otimes\phi_0$, which intertwines
$A\mapsto A\otimes\id$ but needs a state $\phi_0$ of the second particle that nothing selects (constructed). In ordinary physics
the process is creation from a field, and the paper's second quantisation of one constituent already contains it (\cref{thm:spectra}). Whether some mechanism of the arena selects $W$ and $J|_W$ dynamically is open.

\paragraph{Up the ladder.} What has been proved at every rung is this. The commuting copies exist, dressed by $\Gamma^{(1)}\cdots\Gamma^{(k-1)}$.
The graded masses add in quadrature and the ungraded ones additively. The rule of \cref{e1:thm:Nrule} holds. The common complex subspace
$J^{(1)}=\cdots=J^{(n)}$ has dimension $2\cdot4^n$. The diagonal action selects an imaginary unit, $\omega_S=J^{(1)}\cdots J^{(n)}$, exactly at
odd $n$. On the common subspace it acts as $(-1)^{(n-1)/2}J$, so at odd rungs the selected unit and the common complex structure agree up to
sign there, while at even rungs there is no selected unit to compare (fixed; e1-1d, e1-2f, e1-3h, e1-4c, e1-4d). The physics of the odd rungs, whose diagonal sector is the vector-like matter of \cref{sec:actions},
is not examined here.

\begin{table}[t]
\centering
\footnotesize
\begin{tabular}{@{}>{\raggedright\arraybackslash}p{8.3cm}>{\raggedright\arraybackslash}p{4.4cm}>{\raggedright\arraybackslash}p{2.2cm}@{}}
\toprule
statement & class & checks\\
\midrule
\cref{e1:prop:copies}: commuting copy, commutant $\id\otimes M_8(\R)$, grade-$7$ generators, chirality dressing, rung $n$ & restated (\cref{d5:prop:graded}, \cref{d5:prop:fock}, \cref{app:realisation}); commutant dimension and uniqueness fixed; exponential law inherited & e1-1a to e1-1d\\
\cref{e1:prop:rest}: quadrature against addition, rotated clock, joint mass shells, rest flows, rung $n$ & fixed; restated at $m_1=m_2$ (\cref{d6:prop:dirac}(iv)); two-body kinematics inherited & e1-2a to e1-2g\\
\cref{e1:prop:W}: $W=S\otimes_\C S$, $\omega_S$, what preserves $W$, $C_b|_W=M_{12}|_W$ & fixed; tensor algebra inherited & e1-3a, e1-3b, e1-3f\\
\cref{e1:prop:select}: per-particle units $\pm J^{(k)}$; $J_{\rm sel}=J^{(1)}\Gamma^{(2)}$; no unit at $k_I=2$; parity of $n$ & fixed; restated (\cref{thm:i}, \cref{prop:diagtower}; \cref{d5:prop:vacuum}); obstructed & e1-3c to e1-3e, e1-3g, e1-3h\\
\cref{e1:cor:nocx}: no selected two-particle complex structure & fixed (as stated); reading interpretive & ---\\
\cref{e1:prop:Nbiv}: $\ad N=K_n$, $\mathfrak{gl}(3n,\R)$, pairs, parity & restated (\cref{d5:prop:fock}); inherited; realisation fixed & e1-4a, e1-4b, e1-4g\\
\cref{e1:thm:Nrule}: compact $\Leftrightarrow$ conserves $N$; boosts change it by $\pm2$; no gauge field creates pairs & fixed & e1-4c, e1-4d\\
\cref{e1:prop:Nlorentz}: copy rotation hops quanta; $N$ against $\lint$, $J$, $\lG$ & fixed; particle reading interpretive & e1-4e, e1-4f\\
\cref{e1:prop:swap}: $P$ against $\tau_K$ & fixed; $\tau_K$ restated (\cref{d5:prop:cmrel}) & e1-5a, e1-5b\\
\cref{e1:prop:stat}, \cref{e1:tab:stat}: statistics sectors & fixed; Dirac$\otimes$Dirac inherited; $\tau_K$ on $S_2$ restated (\cref{d5:prop:spin}) & e1-5c to e1-5f\\
\cref{e1:sec:answer}: which process; up the ladder; dynamical selection of $W$ & constructed; fixed; open & see text\\
\bottomrule
\end{tabular}
\caption{Claims of \cref{e1:sec:main} and their classes. The checks are the \texttt{PASS} lines printed by the script
\texttt{e1\_two\_from\_one.py}.}
\label{e1:tab:claims}
\end{table}

% --- source: emergence/e2_processes#0 ((head))
% Track e2 -- How one becomes two: vertices, thresholds and pair emergence.
% New references: sections/bib_e2.tex (WeinbergQFT and Chevalley54 are defined in sections/bib_e1.tex).

% --- source: emergence/e2_processes#1 (How one becomes two: processes)
\section{How one becomes two: processes}
\label{e2:sec:main}

\Cref{e1:sec:main} asked which structures make a rung carry several particles. This section asks which processes take one to two. The word
``particle'' has four precise senses in the arena, and each has its own process. A particle can be a spin module of a Lorentz algebra, and
then the process is a vertex, an equivariant map. It can be a mass shell, and then the process is a decay above a threshold. It can be a
field, as in the field-strength reading (\cref{ch:ch03_readings}), and then the process is a duplication of fields. And it can be a quantum of the Fock
structure of the tower, and then the process is a change of the number $N$. The answers, in order, are these.
(1) Under the common interior Lorentz action no equivariant map takes a spin-$\half$ module into the space of two constituents, in either
direction. Maps from a fermion to a fermion and a boson exist, and Clifford multiplication provides them (the selection rule inherited, the
counts in the realisation fixed; \cref{e2:sec:vertex}). (2) At equal momentum labels the top level $E_+$ of the graded composite of two
masses lies strictly below the free-pair energy $E_1+E_2$ iff $m_1m_2\neq0$; at rest the gap is $m_1+m_2-M$, with $(m_1+m_2)^2-M^2=2m_1m_2$
(fixed). This compares two operators on one space, not a bound state with its decay products: at fixed total momentum the graded
levels with large relative momentum lie above the free-pair threshold, so ``binding'' is not established (\cref{e2:sec:threshold}). (3) One null field never becomes a massive pair covariantly, and two non-collinear null fields always form one
(restated from \cref{sec:energy}; \cref{e2:sec:field}). (4) Boosts create pairs of Fock quanta, with the amplitudes of a fermionic Bogoliubov
transformation, and so do the compact generators that do not commute with the binding root, among them the spatial rotations of the interior
four-plane and the charge $J$. The compact generators that commute with the binding root never do (fixed; \cref{e2:sec:quanta}). Under the
common binding action a Lorentz-invariant number of quanta exists exactly at even rungs, and its quanta are shared equally by two constituents
(fixed; \cref{e2:sec:invnumber}). (5) The copy rotation moves one quantum from one constituent to another and never doubles it; at unequal
rates it does so by Rabi mixing (fixed; \cref{e2:sec:hop}). \Cref{e2:sec:answer} collects the processes with their classes. Identifiers
e2-$n$ refer to the \texttt{PASS} lines of \texttt{e2\_processes.py} (\cref{e2:tab:claims}).

\paragraph{Notation.} As in \cref{e1:sec:main}: $S=\R^8$ with the realisation of equation~\cref{eq:realisation}, the tower generators
$E_{(k,a)}=\omega^{\otimes(k-1)}\otimes e_a\otimes\id^{\otimes(n-k)}$ on $S_n=S^{\otimes n}$ (equation~\cref{eq:tower}), $V=\R^{3,3}$,
$V_n=\R^n\otimes V$, $\eta_n=\id_n\otimes\eta$, $G_n=\id_n\otimes G$ and $K_n=\id_n\otimes K$ with $Ke_i=e_{i+3}$, $Ke_{i+3}=e_i$. The spin lift of
$A\in\so{\eta_n}$ is $\sigma(A)=\tfrac14\sum_{a,b}A_{ab}\eta_{bb}E_aE_b$, so that $[\sigma(A),E_c]=\sum_aA_{ac}E_a$. The interior Lorentz algebra
$\lint$ has the rotations $S_1=\half e_5e_6$, $S_2=\half e_6e_4$, $S_3=\half e_4e_5$ and the boosts $T_a=\half e_1e_{a+3}$; it acts on the
interior four-plane $E=\langle e_1,e_4,e_5,e_6\rangle$ with $\gamma^\mu=(e_1,e_4,e_5,e_6)$, and trivially on the surplus pair $\langle e_2,e_3\rangle$
(\cref{sec:dirac}). The field-strength algebra $\lG$, the commutant of $G$ in $\so{3,3}$ (\cref{thm:GLorentz}), has the boosts
$K_i=\half(e_ke_{j+3}-e_je_{k+3})$ and the rotations $L_k=-[K_i,K_j]$, for $(i,j,k)$ cyclic. $J=e_2e_3$. The Fock modes are
$a^\pm_{k,i}=\half(E_{(k,i)}\pm E_{(k,i+3)})$, $N^{(k)}=\sum_ia^+_{k,i}a^-_{k,i}$ and $N=\sum_kN^{(k)}$ (\cref{d5:prop:fock}); $a^-_{k,i}$ is the
transpose of $a^+_{k,i}$. On $S$ the basis vectors of the realisation are the Fock states, and $|n_1n_2n_3\rangle$ denotes the one with
$N_i=n_i$; the Fock vacuum $|000\rangle$ is the first basis vector, and $|0\rangle$ is the first basis vector of $S_n$ (e2-1). $\varepsilon=\bigl(\begin{smallmatrix}0&-1\\1&0\end{smallmatrix}\bigr)$
acts on the copy index. The copy rotation on vectors is $\varepsilon\otimes\id_6$ and the copy boost is $\varepsilon\otimes G$ (\cref{d5:sec:main}). We write
\[
\hat C:=\sigma(\varepsilon\otimes\id_6)=-\half\textstyle\sum_a\eta_{aa}E_{(1,a)}E_{(2,a)} ,
\]
the negative of $C_b$ of \cref{d5:prop:cmrel}(ii), which lifts the copy rotation with the opposite orientation. Dimensions of
spaces of equivariant maps are real dimensions.

% --- source: emergence/e2_processes#2 (Vertices)
\subsection{Vertices}
\label{e2:sec:vertex}

A vertex that takes one particle into two is an equivariant linear map from the module of the one into the module of the two. The
ungraded and graded copies of $\Cl(3,3)$ in $\Cl(6,6)$ have the same bivectors, $E_{(2,a)}E_{(2,b)}=\id\otimes e_ae_b$ (e2-1;
\cref{e1:prop:copies}), so every Lorentz action built from bivectors is the same operator on $S_2=S\otimes S$ in both readings, and the
statements below hold for both.

\begin{proposition}[Vertices]
\label{e2:prop:vertex}
Let $\lint$ act on $S_n$ diagonally, on every constituent, unless stated otherwise.
\begin{enumerate}[label=(\roman*),leftmargin=2.2em]
\item $\dim\Hom_{\lint}(S,S_2)=\dim\Hom_{\lint}(S_2,S)=0$. The rotation $e^{2\pi S_3}$ is $-\id$ on $S$ and $(-1)^n\id$ on $S_n$, so
$\Hom_{\lint}(S,S_n)=0$ for every even $n$. At $n=3$, $\dim\Hom_{\lint}(S,S_3)=160$.
\item The spaces of $\lint$-equivariant maps have the dimensions
\[
\dim\End_{\lint}(S)=8,\quad\dim\Hom_{\lint}(S,S\otimes E)=8,\quad\dim\Hom_{\lint}(S,S\otimes V)=24,
\]
\[
\dim\Hom_{\lint}(S,S\otimes\Lambda^2E)=8,\qquad\dim\Hom_{\lint}(S,S\otimes\Lambda^2V)=32.
\]
\item The Clifford multiplications $c_1(\psi)=\sum_\mu\eta_{\mu\mu}\gamma_\mu\psi\otimes e_\mu$ and
$c_2(\psi)=\sum_{\mu<\nu}\eta_{\mu\mu}\eta_{\nu\nu}\gamma_\mu\gamma_\nu\psi\otimes e_\mu\wedge e_\nu$ are $\lint$-equivariant, and
$\Hom_{\lint}(S,S\otimes E)=c_1\circ\End_{\lint}(S)$, $\Hom_{\lint}(S,S\otimes\Lambda^2E)=c_2\circ\End_{\lint}(S)$. Without the factors $\eta_{\mu\mu}$, $c_1$ is not
equivariant.
\item Under the spectator action, $\lint$ on constituent~$1$ alone, $\dim\Hom(S,S_2)=64$. Under the mixed action, $\lint$ on constituent~$1$ and
$\lG$ on constituent~$2$, $\dim\Hom(S,S_2)=16$ and $\dim\Hom(S,S\otimes V^{(2)})=8$, where $V^{(2)}$ is the vector space of constituent~$2$. On $S$
the $\lG$ Casimir $\sum_iK_i^2-\sum_kL_k^2$ equals $3\,\id$, and $e^{2\pi L_3}=\id$.
\end{enumerate}
\end{proposition}

\begin{proof}
(i) $e^{2\pi S_3}=e^{\pi e_4e_5}=\cos\pi+\sin\pi\,e_4e_5=-\id$ since $(e_4e_5)^2=-\id$. The diagonal lift is a product of $n$ commuting copies,
one per constituent. An equivariant $T$ intertwines the group elements, so $T(-\psi)=(-1)^nT\psi$, and $T=0$ for even $n$; the same holds for
$\Hom(S_n,S)$. For $n=3$, $S\otimes\C=2(\half,0)\oplus2(0,\half)$ by (ii), and the joint weights of the commuting pair $(S_3,T_3)$ on $S_3\otimes\C$
contain $(\half,0)$ and $(0,\half)$ forty times each, so $\dim_\C\Hom=2\cdot40+2\cdot40$. (ii) Over $\C$, $\lint\cong\slr{2,\C}$ and
$S\otimes\C=2(\half,0)\oplus2(0,\half)$, $E\otimes\C=(\half,\half)$, $\Lambda^2E\otimes\C=(1,0)\oplus(0,1)$, while $V=E\oplus2\cdot\mathbf1$ and
$\Lambda^2V=\Lambda^2E\oplus2E\oplus\mathbf1$. The Clebsch--Gordan rule gives $2^2+2^2=8$ for each of $\End(S)$, $S\to S\otimes E$ and $S\to S\otimes\Lambda^2E$,
hence $8+2\cdot8$ and $8+2\cdot8+8$. (iii) $\sum_\mu\eta_{\mu\mu}\gamma_\mu\otimes e_\mu$ is the invariant element of $\Cl(E)\otimes E$ defined by $\eta$; the
equivariance, its failure without $\eta$, and the ranks $8$ of the composites are computed exactly. (iv) Constituent~$2$ is a trivial
module under the spectator action, so $S_2\cong8S$ and the count is $8\cdot8$. On $S$ the $\lG$ Casimir is $3\,\id$, the value on a
four-vector, so $S\otimes\C$ is two copies of $(\half,\half)$, and $\dim\Hom(S,S\otimes2E)=2\cdot8$. $V$ is the field-strength module of $\lG$
(\cref{thm:emfield}), $V\otimes\C=(1,0)\oplus(0,1)$, so $\Hom(S,S\otimes V^{(2)})\cong\Hom(S,S\otimes\Lambda^2E)$.
\end{proof}

Part (i) is the angular-momentum selection rule: a spin-$\half$ state cannot become a state of integer spin, and two spin-$\half$
constituents have integer spin \cite{EhrenfestOppenheimer31,WeinbergQFT}. It restates \cref{prop:orthospin} and
\cref{d6:prop:dirac}(v) as a statement about vertices (inherited; verified in the realisation, e2-2, e2-6). Parts (ii) and (iii) are
inherited Clebsch--Gordan counts, fixed in the realisation (e2-3, e2-4). The vertex $c_1$ is the structure of the coupling
$\bar\psi\gamma^\mu\psi A_\mu$, and $c_2$ that of the Pauli coupling $\bar\psi\gamma^\mu\gamma^\nu\psi F_{\mu\nu}$. In both the boson is an
added tensor factor $E$ or $\Lambda^2E$, not a second constituent. Part (iv) is fixed (e2-5). Under the mixed action the second constituent
itself is the boson. A fermion of constituent~$1$ has an equivariant vertex into $S\otimes V^{(2)}$, where $V^{(2)}$ is read through $\lG$ as a
field strength, the frame field of constituent~$2$ in the field-strength reading. The mixed action is one of the five Lorentz actions at
$n=2$ (\cref{thm:algebra}; \cref{d5:sec:actions}), and nothing selects it, so this vertex is constructed. Under the diagonal action a
fermion has vertices into the odd rungs only, and the first is $n=3$.

% --- source: emergence/e2_processes#3 (Thresholds)
\subsection{Thresholds}
\label{e2:sec:threshold}

The free Dirac Hamiltonian of one constituent is $H(\bp)=me_1+\sum_ap_ae_1e_{a+3}$, with $H(\bp)^2=E^2\id$, $E=(m^2+|\bp|^2)^{1/2}$ (equation~\cref{eq:Hp}). For two constituents with masses $m_1,m_2\ge0$ and momenta $\bp,\mathbf q$ put $E_1=(m_1^2+|\bp|^2)^{1/2}$, $E_2=(m_2^2+|\mathbf q|^2)^{1/2}$,
\[
a=E_1^2+E_2^2,\qquad b=E_1^2E_2^2-m_1^2m_2^2=m_1^2|\mathbf q|^2+m_2^2|\bp|^2+|\bp|^2|\mathbf q|^2 .
\]
The ungraded two-particle operator is $H_u=H(\bp)\otimes\id+\id\otimes H(\mathbf q)$. The graded composite is
$H_g=H^{(1)}(\bp)+H^{(2)}(\mathbf q)$, with $H^{(k)}(\bp)=m_kE_{(k,1)}+\sum_ap_aE_{(k,1)}E_{(k,a+3)}$ built from the tower's vectors.

\begin{proposition}[Thresholds]
\label{e2:prop:threshold}
\begin{enumerate}[label=(\roman*),leftmargin=2.2em]
\item $\spec H_u=\{\pm E_1\pm E_2\}$, the four sign choices on $16$ dimensions each, and $E_1+E_2\ge m_1+m_2$, with equality iff $\bp=\mathbf q=0$.
\item $H_g^2=a\,\id+X$ with $X$ traceless and $X^2=4b\,\id$, and $E_{(1,1)}E_{(2,1)}$ anticommutes with $H_g$. Hence $\spec H_g=\{\pm E_+,\pm E_-\}$ with
\[
E_\pm^2=E_1^2+E_2^2\pm2\bigl(E_1^2E_2^2-m_1^2m_2^2\bigr)^{1/2},
\]
each on $16$ dimensions when the four values are distinct; when values coincide the multiplicities add, and $E_-=0$ exactly when $m_1m_2=0$
and $E_1=E_2$ (for $m_1=m_2=0$ and $\bp=\mathbf q=(1,0,0)$ the spectrum is $-2$, $0$, $2$ on $16$, $32$, $16$ dimensions). $E_+$ is the largest
eigenvalue. At rest $\spec H_g=\{\pm M\}$, $M=(m_1^2+m_2^2)^{1/2}$, each on $32$
dimensions.
\item $\max(E_1,E_2)\le E_+\le E_1+E_2$, and $(E_1+E_2)^2-E_+^2=2\bigl((b+m_1^2m_2^2)^{1/2}-b^{1/2}\bigr)$, which is positive iff $m_1m_2\neq0$. If
$m_1m_2=0$, then $E_\pm=|E_1\pm E_2|$: the graded composite has the additive spectrum.
\end{enumerate}
\end{proposition}

\begin{proof}
(i) $H(\bp)\otimes\id$ and $\id\otimes H(\mathbf q)$ commute, are diagonalisable with eigenvalues $\pm E_1$ and $\pm E_2$ on four dimensions each,
and the joint eigenspaces have dimension $4\cdot4$. $E_k\ge m_k$. (ii) Write $H^{(1)}=m_1A+P$ and $H^{(2)}=m_2B+Q$ with $A=E_{(1,1)}$, $B=E_{(2,1)}$,
and $P$, $Q$ the momentum terms. Each $H^{(k)}$ squares to $E_k^2$, as on one constituent. $A$ anticommutes with $P$ and $B$, and commutes with $Q$;
$B$ anticommutes with $Q$ and commutes with $P$; $P$ and $Q$ commute. Hence $X=\{H^{(1)},H^{(2)}\}=2(m_1AQ+m_2PB+PQ)$. The squares of the three
terms are $m_1^2|\mathbf q|^2$, $m_2^2|\bp|^2$ and $|\bp|^2|\mathbf q|^2$, and the three mutual anticommutators vanish by the same rules, so
$X^2=4b$. Each term is a product of distinct blades, so $X$ is traceless, and $H_g^2$ has the eigenvalues $a\pm2b^{1/2}$ on $32$ dimensions
each. $W=E_{(1,1)}E_{(2,1)}$ anticommutes with $A$, $B$ and every term of $P$ and $Q$ by the blade rule (equation~\cref{eq:bladecomm}), and
$W^2=-\id$, so $WH_gW^{-1}=-H_g$; $W$ commutes with $X$, so each eigenspace of $H_g^2$ carries $\pm E_\pm$ on equal halves. (iii) $E_+^2-E_1^2=E_2^2+2b^{1/2}\ge0$, and $(E_1+E_2)^2=a+2E_1E_2=a+2(b+m_1^2m_2^2)^{1/2}$.
\end{proof}

Part (i) is the threshold of two free particles, inherited: a particle of mass $M$ at rest decays into two of masses $m_1,m_2$ only if
$M\ge m_1+m_2$, since the final energy at zero total momentum is $E_1(\bp)+E_2(-\bp)\ge m_1+m_2$ (e2-7). Part (ii) extends
\cref{d6:prop:dirac}(iv) and \cref{e1:prop:rest}(i) from rest to all momenta. The closed form and (iii) are fixed and
elementary (e2-8: $(H_g^2-a)^2=4b$ exactly at four integer data, the inequalities in $200$ random samples to $10^{-9}$). What follows,
and what does not. At rest the graded composite has the single level $M=(m_1^2+m_2^2)^{1/2}$, below the free-pair threshold $m_1+m_2$; the gap
$m_1+m_2-M$ is positive iff $m_1m_2\neq0$, and it is not proportional to $m_1m_2$ (at $m_1=m_2=m$ it is $(2-\sqrt2)\,m$, linear in $m$); what is
proportional to the product is $(m_1+m_2)^2-M^2=2m_1m_2$. At equal momentum labels, $E_+(\bp,\mathbf q)<E_1(\bp)+E_2(\mathbf q)$ iff $m_1m_2\neq0$, by (iii).
It does not follow that the graded composite lies below the threshold of the free pair at the same total momentum. At $m_1=m_2=m$ and
$\bp=-\mathbf q$, the top level $E_+$ exceeds $m_1+m_2=2m$ as soon as $|\bp|>m/2$ (equality at $|\bp|=m/2$, exactly), and at $\bp=-\mathbf q=(0,0,2)$, $m=1$,
it is $E_+\approx4.45$ (checks y2-8 and z1-2a); at each total momentum only the lowest graded levels lie below $((m_1+m_2)^2+|\mathbf P|^2)^{1/2}$.
So the inequality is not a bound-state condition, and the graded composite is not shown to be stable against dissociation. Both operators act
on the same space $S_2$, nothing in the tower turns one into the other, and neither is derived from a common interacting theory with
asymptotic states; ``binding'' and ``dissociation'' therefore compare two readings and describe no dynamical process. Reading $m_1+m_2-M$
as a binding energy is not supported.

% --- source: emergence/e2_processes#4 (Fields: one null field and two)
\subsection{Fields: one null field and two}
\label{e2:sec:field}

At the level of fields three results of the energy analysis (\cref{sec:energy}) answer the question directly. We restate them as statements about one and two.

\begin{corollary}[One field into two]
\label{e2:cor:field}
\begin{enumerate}[label=(\roman*),leftmargin=2.2em]
\item One null field does not become a massive pair. Every map commuting with the diagonal $SO(3,\C)$ sends $(z,0)$, $z$ null, to
$(\alpha z,\beta z)$, whose Gram matrix $Q$ vanishes, so $M_{\rm pair}=\half(\sigma_1+\sigma_2)=0$ (\cref{d6:prop:covariant}(iii) and \cref{d6:thm:pair}(ii)).
\item Two null fields that are not collinear do: $M_{\rm pair}=|z_1\cdot z_2|>0$, and the copy rotation can move the whole of it into rest
energy of the two constituents, in kinematic analogy with $\gamma\gamma\to e^+e^-$ (\cref{d6:thm:pair}(iv),(v) and \cref{d6:ex:bw}; \cite{BreitWheeler34});
$M_{\rm pair}$ is a rest-energy density, not an invariant mass.
\item The closure moves nothing between constituents: for every Hamiltonian the per-constituent invariants of the momentum fields are
constant, so no closure dynamics fills an empty constituent or makes a null one massive (\cref{d6:cor:nopair}).
\end{enumerate}
\end{corollary}

All three are restated (fixed there; re-verified for the Gram matrices, e2-9). Read as one and two: at the field level ``one becomes two''
covariantly only as ``two become two''. A single null field is shared between the constituents as literal replicas, and a massive pair
needs a second, non-collinear field or a non-covariant third body that takes the momentum (\cref{d6:ex:R1}).

% --- source: emergence/e2_processes#5 (Quanta: pairs emerge under boosts)
\subsection{Quanta: pairs emerge under boosts}
\label{e2:sec:quanta}

The tower carries a counting operator of its own, the Fock number $N$. We now ask which Lie-algebra elements change it, and with which
amplitudes. Throughout, ``pair'' means two Fock quanta of $N$ on the finite space $S_n$, and ``creates a pair'' means that a Lie-algebra element
has matrix elements with $\Delta N=\pm2$, or that its exponential gives the Fock vacuum a component with $N=2$. These are statements about a
fixed finite-dimensional space at one momentum label. They are not statements about observable particle production: that would need a
physical vacuum, a positive inner product preserved by the evolution, a particle-number observable defined from a Poincar\'e-invariant
state space, and a dynamics that generates the transformation. None of these is specified here, and boosts are not unitary in the
realisation. The first lemma refines \cref{e1:prop:Nbiv} and \cref{e1:thm:Nrule}.

\begin{lemma}[The number and the split structure]
\label{e2:lem:N}
\begin{enumerate}[label=(\roman*),leftmargin=2.2em]
\item $N=\tfrac{3n}2\id+\sigma(K_n)$, and $\ad N=K_n$ on $V_n$. $(-1)^N=-\omega$ at $n=1$ and is the volume element of $\Cl(6,6)$ at $n=2$; it commutes with
every bivector.
\item $\so{3n,3n}=\Cent(K_n)\oplus\mathrm{Anti}(K_n)$, the parts commuting and anticommuting with $K_n$, of dimensions $9n^2$ and $9n^2-3n$. The lift of an
element of $\Cent(K_n)$ commutes with $N$. The lift of an element of $\mathrm{Anti}(K_n)$ maps each eigenspace $\{N=k\}$ into $\{N=k-2\}\oplus\{N=k+2\}$:
it changes $N$ by exactly $\pm2$.
\item An element of the maximal compact subalgebra $\so{3n}\oplus\so{3n}$ of $\so{3n,3n}$ (an antisymmetric matrix in the realisation; called
\emph{compact} below) commutes with $N$ iff it commutes with $G_n$, iff it is a paired rotation
$\bigl(\begin{smallmatrix}R&0\\0&R\end{smallmatrix}\bigr)$, $R\in\so{3n}$, in the basis of the temporal directions followed by their partners. The
restriction to $\so{3n}\oplus\so{3n}$ is needed: the element $\bigl(\begin{smallmatrix}A&B\\B&A\end{smallmatrix}\bigr)$ with $A$ the rotation of
$(e_1,e_2)$ and $B=\operatorname{diag}(\half,-\half,0)$ lies in $\so{3,3}$, generates a periodic one-parameter group and commutes with $K$ and $N$, but
not with $G$ (y2-13).
\item Let $X\in\so{3n,3n}$ be compact in this sense or Euclidean-symmetric (a pure boost); for other $X$ the conclusion fails (y2-14). If $X\notin\Cent(K_n)$, then $\sigma(X)|0\rangle$ has a nonzero
component with $N=2$: $e^{t\sigma(X)}|0\rangle$ contains a pair at first order in $t$.
\end{enumerate}
\end{lemma}

\begin{proof}
(i) and (ii) restate \cref{e1:prop:Nbiv}: $N_i=\half(\id-e_ie_{i+3})$ and $\sigma(K)=-\half\sum_ie_ie_{i+3}$, $K$ is an $\eta$-antisymmetric involution,
and $\ad K$ has eigenvalues $0,\pm2$ on $\Lambda^2V_n$. (iii) A compact element is $X=\operatorname{diag}(A,D)$ with $A,D\in\so{3n}$. In the same
basis $K_n=\bigl(\begin{smallmatrix}0&\id\\\id&0\end{smallmatrix}\bigr)$ and $G_n=\bigl(\begin{smallmatrix}0&-\id\\\id&0\end{smallmatrix}\bigr)$, and
both $K_nXK_n^{-1}$ and $G_nXG_n^{-1}$ equal $\operatorname{diag}(D,A)$. So each commutation is $A=D$. By (ii) commuting with $K_n$ is
commuting with $N$, because the spin representation is faithful on bivectors. (iv) A compact $X$ is a combination of blades $e_ae_b$ with
$\eta_{aa}=\eta_{bb}$, which are antisymmetric matrices, and a Euclidean-symmetric one of blades with $\eta_{aa}=-\eta_{bb}$, which are symmetric. So
$\sigma(X)^{\!\top}=\mp\sigma(X)$. Write $\sigma(X)=R+Z+L$ with $[N,R]=2R$, $[N,Z]=0$, $[N,L]=-2L$, by (ii). Transposition exchanges raising and
lowering parts ($a^-=(a^+)^{\!\top}$), so $L=\mp R^{\!\top}$, and $X\notin\Cent(K_n)$ gives $R\neq0$. $R=\sum_{i<j}c_{ij}a^+_ia^+_j$ in the $3n$ modes,
and the states $a^+_ia^+_j|0\rangle$, $i<j$, are linearly independent, so $R|0\rangle\neq0$; it is the $N=2$ component of $\sigma(X)|0\rangle$.
\end{proof}

(i) and (ii) are inherited, the spin representation on the exterior algebra of a maximal isotropic subspace \cite{Chevalley54}, and fixed
in the realisation at $n=1,2,3$ (e2-10, e2-11). (iii) and (iv) are fixed (e2-12, e2-26). Part (iii) sharpens \cref{e1:thm:Nrule}, which works
inside $\Cent(G_n)$: among all compact elements, the number-conserving ones are exactly the compact binding-compatible ones. By (iv) a
compact element that does not commute with $G_n$ creates a pair from the vacuum. \Cref{e2:prop:bogo} gives the examples, and
e2-13 records which generators of $\lint$, $\lG$ and the copy algebras lie where: $\lint\cap\Cent(K)=\R\,T_1$, the pole boost; every
other generator of $\lint$, and $J$, has a pair part; $\lG\cap\Cent(K)$ is spanned by the $L_k$, and the $K_i$ lie in $\mathrm{Anti}(K)$; the copy
rotation $\varepsilon\otimes\id_6$ lies in $\Cent(K_2)$ and the copy boost $\varepsilon\otimes G$ in $\mathrm{Anti}(K_2)$.

\begin{proposition}[Pair amplitudes on one constituent]
\label{e2:prop:bogo}
\begin{enumerate}[label=(\roman*),leftmargin=2.2em]
\item Interior boost. $\ad T_2$ has eigenvalues $\pm1$ and $0$, so $U_\zeta=e^{\zeta T_2}$ is the boost of rapidity $\zeta$ in the plane
$(e_1,e_5)$. $U_\zeta|000\rangle=\cosh\tfrac\zeta2\,|000\rangle+\sinh\tfrac\zeta2\,|110\rangle$ and $\|U_\zeta|000\rangle\|^2=\cosh\zeta$. $U_\zeta$ is symmetric and not
orthogonal.
\item Field-strength boost. $\ad K_1$ has eigenvalues $\pm1$ (twice) and $0$ (twice), and
$e^{\zeta K_1}|000\rangle=\cosh\zeta\,|000\rangle-\sinh\zeta\,|011\rangle$. The rotations $L_k$ commute with $N$.
\item Compact generators outside $\Cent(G)$. $e^{\theta S_1}|000\rangle=\cos\tfrac\theta2\,|000\rangle+\sin\tfrac\theta2\,|011\rangle$ and
$e^{\theta J}|000\rangle=\cos\theta\,|000\rangle+\sin\theta\,|011\rangle$. A spatial rotation through $\pi$ about the axis of $S_1$ carries the vacuum to a state
of two quanta.
\end{enumerate}
\end{proposition}

\begin{proof}
Each generator is a sum of commuting blades $B$ with $B^2=\pm\id$, so its exponential is a product of factors $\cos t+\sin t\,B$ or
$\cosh t+\sinh t\,B$. In the realisation, with $\varepsilon=\sigma_x\sigma_z$ the matrix used on the copy index,
$e_1e_5=\varepsilon\otimes\varepsilon\otimes\id$, $e_5e_6=\id\otimes\sigma_x\otimes\varepsilon$ and $e_2e_3=\id\otimes\varepsilon\otimes\sigma_x$. Each maps
$|000\rangle$ to the basis vector with two occupation numbers flipped, $|110\rangle$, $|011\rangle$ and $|011\rangle$, and $(e_1e_5)^2=\id$,
$(e_5e_6)^2=(e_2e_3)^2=-\id$. The terms $\half e_3e_5$ and $-\half e_2e_6$ of $K_1$ commute, square to $\tfrac14\id$ and both flip modes $2$ and $3$, so
$e^{\zeta K_1}|000\rangle=(\cosh^2\tfrac\zeta2+\sinh^2\tfrac\zeta2\,s)|000\rangle+\cosh\tfrac\zeta2\sinh\tfrac\zeta2\,(e_3e_5-e_2e_6)|000\rangle$ with $s=-e_3e_5e_2e_6$ on
$|000\rangle$; the signs, $s=1$ and $(e_3e_5-e_2e_6)|000\rangle=-2|011\rangle$, are computed (e2-14 to e2-16).
\end{proof}

All three parts are fixed (e2-14 to e2-16, at $\zeta=0.3,1,2.5$ and $\theta=0.4,1.3,\pi$, to $10^{-11}$). The amplitudes have the form of a
fermionic Bogoliubov transformation \cite{Bogoliubov58,Berezin66}, inherited. A boost of rapidity $\zeta$ turns the vacuum into
$\cos\vartheta\,|0\rangle+\sin\vartheta\,|\text{pair}\rangle$ after normalisation, with $\tan\vartheta=\tanh\tfrac\zeta2$ for $T_2$ and $\tan\vartheta=\tanh\zeta$ for
$K_1$. The normalised persistence $\cos^2\vartheta$ is $\cosh^2\tfrac\zeta2/\cosh\zeta$ and $\cosh^2\zeta/\cosh2\zeta$, and it decreases to $\half$ as
$\zeta\to\infty$. So for these boosts the normalised pair probability $\tanh^2/(1+\tanh^2)$ stays below $\half$, which is its large-rapidity
limit, where $\cosh$ and $\sinh$ become equal. This is not a Pauli bound: under the unitary rotation $e^{\pi S_1}$ of (iii) the pair probability
is $1$, and Pauli exclusion bounds the occupation of each mode pair by one (y2-17). The normalisation is needed because a boost does not preserve the Euclidean inner
product of the realisation, which is the Fock inner product ($a^-_i=(a^+_i)^{\!\top}$): in the split realisation boosts are symmetric and rotations
antisymmetric, and the Lorentz group, being non-compact, has no finite-dimensional unitary representation other than the trivial one
\cite{WeinbergQFT} (inherited). In ordinary field theory the free-particle number is Poincar\'e invariant, and a vacuum turns into
pairs only for an observer whose modes are not related to the inertial ones by a Poincar\'e transformation, as in the Unruh effect
\cite{Unruh76}. Here the inertial boosts of the interior four-plane already mix the modes, because $N$ is built from the split structure
$K$ and not from the Lorentz algebra. The parallel with the Unruh effect is interpretive.

Part (iii) corrects an easy expectation. Compactness does not protect the number of quanta. A spatial rotation of the interior four-plane
and the charge $J$, both compact, create pairs unitarily, and at $\theta=\pi$, resp.\ $\theta=\tfrac\pi2$, the vacuum is turned wholly into a pair.
What protects $N$ is commutation with the binding root (\cref{e2:lem:N}(iii)).

\begin{proposition}[Pairs across constituents]
\label{e2:prop:copyboost}
At $n=2$, let $U_\zeta=e^{\zeta\sigma(\varepsilon\otimes G)}$, the copy boost of rapidity $\zeta$.
\begin{enumerate}[label=(\roman*),leftmargin=2.2em]
\item $\sigma(\varepsilon\otimes G)=-\sum_{i=1}^3\bigl(a^+_{1,i}a^+_{2,i}+a^-_{2,i}a^-_{1,i}\bigr)$. Hence
$U_\zeta|0\rangle=\prod_i\bigl(\cosh\zeta-\sinh\zeta\,a^+_{1,i}a^+_{2,i}\bigr)|0\rangle$.
\item $\langle0|U_\zeta|0\rangle=\cosh^3\zeta$; the weight of $U_\zeta|0\rangle$ in $\{N=2k\}$ is $\binom3k\cosh^{2(3-k)}\zeta\,\sinh^{2k}\zeta$, of total $\cosh^32\zeta$, and
there is none in odd $N$.
\item $N^{(1)}=N^{(2)}$ on $U_\zeta|0\rangle$: every pair that emerges has one quantum in each constituent.
\end{enumerate}
\end{proposition}

\begin{proof}
(i) $\varepsilon\otimes G$ anticommutes with $K_2$, so its lift lies in $\Lambda^2L_+\oplus\Lambda^2L_-$ (\cref{e1:prop:Nbiv}(ii)), and it pairs the modes
$(1,i)$ and $(2,i)$ because $\varepsilon\otimes G$ maps the plane $(e^{(1)}_i,e^{(1)}_{i+3})$ to $(e^{(2)}_i,e^{(2)}_{i+3})$. The coefficients are
computed exactly, with $\sigma(\varepsilon\otimes G)$ symmetric. The three terms commute. On the two-dimensional span of $|0\rangle$ and
$a^+_{1,i}a^+_{2,i}|0\rangle$ the $i$-th term acts as $-\sigma_x$ and it annihilates the states with one quantum in the pair $(1,i),(2,i)$. (ii)
and (iii) follow by expanding the product.
\end{proof}

Fixed (e2-17, exact for (i), and for (ii), (iii) at $\zeta=0.3,1,2$ to relative $10^{-9}$); the lift of the copy boost is checked against
$\varepsilon\otimes G$ on all twelve generators. This is pair emergence across the constituents. The copy boost is a metric connection
between two binding frames but not a gauge field (\cref{d7:prop:boost}; \cref{cor:noboost}). Each emerging pair has
one quantum in each constituent, with normalised persistence $\cosh^6\zeta/\cosh^32\zeta=(1+\tanh^2\zeta)^{-3}$.

\begin{proposition}[Lorentz-invariant states]
\label{e2:prop:invstates}
\begin{enumerate}[label=(\roman*),leftmargin=2.2em]
\item At $n=1$ no nonzero state is $\lint$-invariant. The two projectors $\half(\id\pm\omega)$ of \cref{prop:rigid} are the projectors
onto odd and even $N$. The paired-mode projector of \cref{ex:pairedmode} is the projector onto $N\ge2$.
\item At $n=2$ the states invariant under the diagonal $\lint$ form a space of dimension $8$. Its weights in the sectors $N=0,\dots,6$ (the
traces of the invariant projector restricted to them) are $0,\half,2,3,2,\half,0$. It contains no state with $N=0$ or $N=6$, and its states of
definite $N$ form a space of dimension $2$, at $N=3$, on which $N^{(1)}_i+N^{(2)}_i=1$ for $i=1,2,3$.
\item Under the spectator action no nonzero state of $S_2$ is invariant. The projector $\Pi_0=\half(\id+H)$ of \cref{thm:towervac}
commutes with $(-1)^N$ and not with $N$, and its compression to every sector $\{N=k\}$ has trace equal to half the dimension of that sector.
\end{enumerate}
\end{proposition}

\begin{proof}
(i) $S\otimes\C$ contains only spin $\half$. $\omega$ is $-(-1)^N$ (\cref{e2:lem:N}(i)). With $\varpi=\sum_ie_ie_{i+3}=\sum_i(1-2N_i)$ and
the volume element $\omega$, the projector $\half\id-\tfrac14(\varpi+\omega)$ of \cref{ex:pairedmode} vanishes exactly where at most one mode is occupied, by the
table of signs given there. (ii) and (iii) are exact kernel computations; the weights are traces of the orthogonal projector onto the
kernel. $H=E_{(1,2)}E_{(1,3)}E_{(2,1)}E_{(2,2)}$ is a blade of even grade, so it commutes with the volume element $(-1)^N$. It shares one
index with each of the pole operators of the modes $(1,2)$, $(1,3)$, $(2,1)$, $(2,2)$ and none with the other two, so it anticommutes with
those four and is a signed permutation of the Fock basis that flips exactly those four occupation numbers. Hence $H|0\rangle$ has $N=4$, so
$H$ and $\Pi_0$ do not commute with $N$, and $H$ has zero diagonal, so the compression of $\Pi_0$ to every sector has trace half its
dimension.
\end{proof}

Fixed (e2-18). At one constituent the number of quanta is not Lorentz invariant (\cref{e1:prop:Nlorentz}), while its parity is: the
Lorentz-invariant projectors of \cref{prop:rigid} are the parity projectors, and the paired-mode projector, a function of $N$
that is not a function of the parity, is not Lorentz invariant. The Fock vacuum is not Lorentz invariant at either rung. At the second rung the Lorentz scalars of the diagonal action are spread
over all numbers from $1$ to $5$; the two scalars of definite number carry one quantum in each mode pair $i$, shared between the
constituents. The Lorentz-invariant vacuum of \cref{thm:towervac} is not a function of $N$. It selects half of every number sector,
and so it is not a Fock vacuum.

% --- source: emergence/e2_processes#6 (A Lorentz-invariant number at even rungs)
\subsection{\texorpdfstring{A Lorentz-invariant number at even rungs}{A Lorentz-invariant number at even rungs}}
\label{e2:sec:invnumber}

The Fock number is the lift of the involution $K_n$ (\cref{e2:lem:N}(i)), and $K_n$ commutes with no Lorentz algebra of the tower except
through the pole boost. Any other $\eta_n$-antisymmetric involution defines a number in the same way. The question is whether one commutes
with a Lorentz action.

\begin{theorem}[An invariant number of quanta]
\label{e2:thm:invnumber}
\begin{enumerate}[label=(\roman*),leftmargin=2.2em]
\item Let $X\in\so{3n,3n}$ with $X^2=\id$. Its eigenspaces $V_\pm$ are totally null of dimension $3n$, and $N_X:=\tfrac{3n}2\id+\sigma(X)$ has spectrum
$0,\dots,3n$ with multiplicities $\binom{3n}k$. If $X$ commutes with a Lie algebra $\mathfrak l\subset\so{3n,3n}$, then $N_X$ commutes with $\sigma(\mathfrak l)$.
\item Under the common binding action (the diagonal $\lG$, $(k_I,k_B,k_S)=(0,n,0)$ in \cref{thm:algebra}) such an $X$ exists iff $n$ is even.
At $n=2$ the only ones are $\pm\varepsilon\otimes G$, with $V_\pm=\{X^{(2)}=\pm GX^{(1)}\}$, the null subspaces of \cref{d6:prop:pairs}(ii). At
$n=4$, $(\varepsilon\oplus\varepsilon)\otimes G$ is one; it is not unique, since each of the three pairings of the four constituents gives one.
\item At $n=2$ none exists under the other four actions: the commutants, of dimensions $7$, $1$, $28$ and $15$ for $(2,0,0)$, $(1,1,0)$, $(1,0,1)$ and
$(0,1,1)$, consist of compact elements for the first two and of elements vanishing on the Lorentz-carrying subspace of constituent~$1$ for the
last two. At $n=3$ the action $(0,2,1)$ admits $(\varepsilon\otimes G)\oplus K$.
\item At $n=2$, with $N_K:=N_{\varepsilon\otimes G}$: the diagonal-$\lG$-invariant states form a space of dimension $4$, with weights $1,2,1$ at
$N_K=0,3,6$. The vacuum $|0_K\rangle$ of $N_K$ is $\lG$-invariant, has weights $\tfrac18,\tfrac38,\tfrac38,\tfrac18$ at $N=0,2,4,6$, and equals the
limit of $U_\zeta|0\rangle/\|U_\zeta|0\rangle\|$ as $\zeta\to-\infty$. Every vector of $V_\pm$ has components of equal length in the two constituents.
\end{enumerate}
\end{theorem}

\begin{proof}
(i) For $v,w\in V_+$, $\eta(v,w)=\eta(Xv,w)=-\eta(v,Xw)=-\eta(v,w)$, so $V_+$ is totally null, and likewise $V_-$; as $V_n=V_+\oplus V_-$ and
null subspaces of $\R^{3n,3n}$ have dimension at most $3n$, both have dimension $3n$. Two pairs of complementary maximal null subspaces are
conjugate under $O(\eta_n)$ (Witt's theorem), so $X=gK_ng^{-1}$, and a lift $\hat g$ of $g$ to the Pin group gives
$\sigma(X)=\hat g\sigma(K_n)\hat g^{-1}$. The spectrum of $\sigma(K_n)=N-\tfrac{3n}2$ is that of \cref{e2:lem:N}. Finally $\sigma$ is a Lie algebra
homomorphism. (ii) The commutant of the diagonal $\lG$ in $\so{3n,3n}$ is $\{A\otimes\id_6+B\otimes G:A,B\in\so n\}$, which acts on $z\in\C^n\otimes\C^3$
as $(A+iB)\otimes\id_3$ (\cref{d7:prop:exist}(i), and \cref{d5:prop:so6C}(iii) at $n=2$; commuting with the diagonal $\lG$ already forces commuting with $G_n$, since $\End_{\lG}(V)=\operatorname{span}\{\id,G\}$). $X^2=\id$ requires $(A+iB)^2=\id$, so $A+iB$ is invertible; but a complex antisymmetric
matrix of odd size has $\det(A+iB)=\det(A+iB)^{\!\top}=-\det(A+iB)=0$. For even $n$, $A+iB=i\varepsilon\oplus\cdots\oplus i\varepsilon$ squares to $\id$. At
$n=2$ the solutions of $(a\,\varepsilon\otimes\id_6+b\,\varepsilon\otimes G)^2=\id$ are $a=0$, $b=\pm1$. (iii) and (iv) are computed. For the limit, $|0_K\rangle$ is the
eigenvector of $\sigma(\varepsilon\otimes G)$ for its lowest eigenvalue $-3$, and $\langle0|0_K\rangle\neq0$. On $V_\pm$, $|X^{(2)}|=|GX^{(1)}|=|X^{(1)}|$.
\end{proof}

Part (i) is inherited, the correspondence between polarisations and Fock structures \cite{Chevalley54,Berezin66}; the spectrum and the
commutation are fixed in the realisation at $n=2$ (e2-19) and on $S_4$, of dimension $4096$, where $\sigma((\varepsilon\oplus\varepsilon)\otimes G)$ was
checked to lift its vector and to commute with the six lifted diagonal $\lG$ generators in sparse integer arithmetic (e2-22). Part (ii) is
fixed; the parity argument is general, and the exact checks cover $n=2,3,4$ (e2-19, e2-21). Part (iii) is fixed and re-verifies the
dimensions of \cref{thm:algebra} (e2-21). Part (iv) is fixed (e2-20; the limit to $10^{-9}$ at $\zeta=-12$).

So a Lorentz-invariant number of quanta exists, under the common binding action, exactly at even rungs. At the second rung it is unique up
to $N_K\mapsto6-N_K$. Its quanta are not quanta of one constituent: $V_\pm$ is the graph of $\pm G$ from constituent~$1$ to constituent~$2$, and
$N_K$ does not commute with $N^{(1)}$, since $\sigma(\varepsilon\otimes G)$ changes $N^{(1)}$ and $N^{(2)}$ by one each (\cref{e2:prop:copyboost}). Its
vacuum, read in the Fock structure of the constituents, is a condensate of pairs with one quantum in each constituent, the infinite
copy boost of the Fock vacuum (fixed). Read as a particle number, $N_K$ counts quanta that are Lorentz invariant under the binding action
and belong to two constituents at once (interpretive). The parity pattern can be set beside \cref{e1:prop:select}(iv): under the diagonal
interior action an imaginary unit is selected exactly at odd rungs, and under the diagonal binding action an invariant number exists
exactly at even rungs. The two statements concern different Lorentz actions, and the juxtaposition is a reading.

\begin{table}[!htbp]
\centering
\footnotesize
\begin{tabular}{@{}>{\raggedright\arraybackslash}p{4.0cm}>{\raggedright\arraybackslash}p{1.8cm}>{\raggedright\arraybackslash}p{5.3cm}>{\raggedright\arraybackslash}p{3.6cm}@{}}
\toprule
process & sense & result & class\\
\midrule
$S\to S_n$, common interior action & module & none for even $n$; $160$ maps at $n=3$ & inherited (selection rule); fixed\\
$S\to S\otimes E$, $S\otimes\Lambda^2E$; $S\to S\otimes V^{(2)}$ (mixed) & module & Clifford multiplication, $8$ maps each & inherited counts; fixed; mixed action constructed\\
free particle of mass $M$ $\to$ two & mass shell & only if $M\ge m_1+m_2$ & inherited\\
graded composite $\to$ free pair & mass shell & $E_+<E_1+E_2$ iff $m_1m_2\neq0$ at equal labels; above threshold at large relative momentum & fixed; binding not supported\\
one null field $\to$ massive pair & field & no covariant map; closure moves nothing & obstructed (restated)\\
two null fields $\to$ pair with $M_{\rm pair}>0$ & field & $M_{\rm pair}=|z_1\cdot z_2|$, a rest-energy density & inherited (restated)\\
boosts ($\lint$, $\lG$, copy) & quanta & pairs, fermionic Bogoliubov amplitudes & fixed; form inherited\\
compact, outside $\Cent(G_n)$ ($S_k$, $J$) & quanta & pairs, unitarily & fixed\\
compact, inside $\Cent(G_n)$ (paired rotations) & quanta & $N$ conserved & fixed\\
copy rotation & quanta & one quantum hops; Rabi at unequal rates & fixed; two-level inherited; $i$ constructed\\
$N_X$ under $(0,n,0)$ & quanta & exists iff $n$ even; vacuum a pair condensate & fixed; particle reading interpretive\\
dynamical pair-creating agent & --- & --- & open\\
\bottomrule
\end{tabular}
\caption{Processes that take one to two, by sense of ``particle''.}
\label{e2:tab:processes}
\end{table}

\begin{table}[!htbp]
\centering
\footnotesize
\begin{tabular}{@{}>{\raggedright\arraybackslash}p{8.3cm}>{\raggedright\arraybackslash}p{4.4cm}>{\raggedright\arraybackslash}p{2.2cm}@{}}
\toprule
statement & class & checks\\
\midrule
conventions: tower, ungraded copy, common bivectors, CAR, Fock vacuum & restated (\cref{d5:prop:graded}, \cref{d5:prop:fock}) & e2-1\\
\cref{e2:prop:vertex}(i): no vertex $S\leftrightarrow S_2$; $(-1)^n$; $160$ at $n=3$ & inherited; fixed in realisation & e2-2, e2-6\\
\cref{e2:prop:vertex}(ii), (iii): counts $8,8,24,8,32$; Clifford multiplication & inherited counts; fixed & e2-3, e2-4\\
\cref{e2:prop:vertex}(iv): spectator $64$; mixed $16$ and $8$; $\lG$ Casimir $3$ & fixed; mixed action constructed & e2-5\\
\cref{e2:prop:threshold}(i): free-pair threshold & inherited & e2-7\\
\cref{e2:prop:threshold}(ii), (iii): $(H_g^2-a)^2=4b$; $E_+<E_1+E_2$ iff $m_1m_2\neq0$; crossover $|\bp|=m/2$ & fixed; binding not supported & e2-8, y2-8, z1-2a\\
\cref{e2:cor:field}: one null field, two null fields, closure & restated (\cref{d6:prop:covariant}, \cref{d6:thm:pair}, \cref{d6:cor:nopair}) & e2-9\\
\cref{e2:lem:N}(i), (ii): $N=\tfrac{3n}2+\sigma(K_n)$; $\Cent(K_n)\oplus\mathrm{Anti}(K_n)$ & restated (\cref{e1:prop:Nbiv}); inherited & e2-10, e2-11\\
\cref{e2:lem:N}(iii), (iv): compact and number-conserving $=$ paired rotations; outside $\Cent(K_n)$ the vacuum acquires a pair & fixed & e2-12, e2-13, e2-26\\
\cref{e2:prop:bogo}: Bogoliubov amplitudes; pairs from $S_1$ and $J$ & fixed; form inherited; Unruh analogy interpretive & e2-14 to e2-16\\
\cref{e2:prop:copyboost}: copy boost, pairs across constituents & fixed & e2-17\\
\cref{e2:prop:invstates}: invariant states at $n=1,2$; \cref{prop:rigid}, \cref{thm:towervac} & fixed (compared with these) & e2-18\\
\cref{e2:thm:invnumber}: invariant number iff $n$ even; $N_K$ at $n=2$ & fixed; (i) inherited; reading interpretive & e2-19 to e2-22\\
\cref{e2:prop:hop}: hopping, ladder, Rabi transfer & fixed; Rabi inherited; external $i$ constructed & e2-23 to e2-25\\
\cref{e2:sec:answer}: dynamical agent of pair creation & open & ---\\
\bottomrule
\end{tabular}
\caption{Claims of \cref{e2:sec:main}, their classes and the checks printed by \texttt{e2\_processes.py}.}
\label{e2:tab:claims}
\end{table}

% --- source: emergence/e2_processes#7 (The copy rotation moves one quantum)
\subsection{The copy rotation moves one quantum}
\label{e2:sec:hop}

\begin{proposition}[Hopping and Rabi transfer]
\label{e2:prop:hop}
\begin{enumerate}[label=(\roman*),leftmargin=2.2em]
\item $\hat C=\sigma(\varepsilon\otimes\id_6)=\sum_i\bigl(a^+_{2,i}a^-_{1,i}-a^+_{1,i}a^-_{2,i}\bigr)$. It commutes with $N$ and not with $N^{(1)}$;
$\hat Ca^+_{1,i}|0\rangle=a^+_{2,i}|0\rangle$ and $\hat Ca^+_{2,i}|0\rangle=-a^+_{1,i}|0\rangle$, and
$e^{\theta\hat C}a^+_{1,1}|0\rangle=\cos\theta\,a^+_{1,1}|0\rangle+\sin\theta\,a^+_{2,1}|0\rangle$. $e^{\theta\hat C}$ has no matrix element between different values
of $N$, and $e^{2\pi\hat C}=\id$.
\item At $n=3$ the copy rotations commute with $K_3$, and their lifts with $N$; the copy boosts anticommute with $K_3$.
\item With rates $\kappa_1,\kappa_2$ and coupling $\beta$, let $\mathsf H=\kappa_1N^{(1)}+\kappa_2N^{(2)}+i\beta\hat C$, with an external $i$. $\mathsf H$ is
Hermitian and commutes with $N$. On the sector $N=1$, in the basis $a^+_{1,i}|0\rangle$, $a^+_{2,i}|0\rangle$, it is
$\bigl(\begin{smallmatrix}\kappa_1&-i\beta\\ i\beta&\kappa_2\end{smallmatrix}\bigr)\otimes\id_3$, whose eigenvalues are the lines $\omega_\pm$ of
\cref{d5:thm:bfield2}(ii). From $a^+_{1,1}|0\rangle$ the probability of finding the quantum in constituent~$2$ is
$s\sin^2(\Omega t/2)$, with $s=4\beta^2/((\kappa_2-\kappa_1)^2+4\beta^2)$ and $\Omega=((\kappa_2-\kappa_1)^2+4\beta^2)^{1/2}$.
\end{enumerate}
\end{proposition}

\begin{proof}
(i) $\varepsilon\otimes\id_6$ commutes with $K_2$ and maps the vector of each mode of constituent~$1$ to the same mode of constituent~$2$, so its
lift lies in $L_+\otimes L_-$ and is a hopping term; the coefficients are computed. The rest follows from $[\hat C,a^\pm_{1,i}]=a^\pm_{2,i}$ and
$[\hat C,a^\pm_{2,i}]=-a^\pm_{1,i}$, and from $\hat C$ being a sum of six commuting terms of square $-\tfrac14$ (\cref{d5:prop:cmrel}(ii)). (ii)
$GK=-KG$. (iii) $\hat C$ is real antisymmetric, so $i\beta\hat C$ is Hermitian. On the sector $N=1$, $\hat C$ acts as $\varepsilon\otimes\id_3$, and the
two-level formula follows.
\end{proof}

Fixed (e2-23 to e2-25; the transfer probability at $(\kappa_1,\kappa_2,\beta)=(1,3,0.4)$ to $2\times10^{-16}$, with weight below $10^{-24}$
outside $N=1$, and complete transfer at equal rates at $t=\pi/(2\beta)$). The two-level formula is inherited \cite{Rabi37}. The external $i$ is
constructed, as in \cref{d6:prop:dirac}(i); on fields the same matrix acts with $G_2$ as the imaginary unit and gives the per-copy
rest energies of \cref{d6:sec:gauge}. Read as one and two: one quantum can be shared between constituents, in any proportion at equal
rates and up to the fraction $s$ at unequal rates, but it is never doubled, because $\hat C$ commutes with $N$. At equal rates the
coupling is removable (\cref{d5:thm:bfield2}(iii)), so which constituent carries the quantum is gauge dependent if the copy rotation is
gauged, as the per-copy rest energy is (\cref{d6:prop:gauge}). At unequal rates the partial transfer cannot be removed while the rates are kept fixed. Up the ladder the
copy rotations $\so n$ hop quanta among $n$ constituents and the copy boosts create pairs across them (e2-24).

% --- source: emergence/e2_processes#8 (Answer: which processes take one to two)
\subsection{Answer: which processes take one to two}
\label{e2:sec:answer}

\Cref{e2:tab:processes} lists the processes. Read together they say the following.
In the sense of Lorentz modules, one fermion never becomes two under the common interior action. It becomes a fermion and a boson, by Clifford
multiplication, and under the mixed action the boson can be the second constituent read as a field strength (fixed; the mixed action
constructed). In the sense of mass shells, the graded composite at rest lies below the free-pair threshold iff $m_1m_2\neq0$, and at equal
momentum labels its top level lies below the free-pair energy; at fixed total momentum its levels with large relative momentum lie above
the threshold, so no bound state and no binding energy are established (fixed; binding not supported). In the
sense of fields, one null field becomes a massive pair only together with a second field (restated). In the sense of quanta, no element of
$\so{3n,3n}$ changes $N$ by an odd number, so one quantum never becomes two. The vacuum becomes partly a pair under every compact or
Euclidean-symmetric element outside $\Cent(K_n)$: the boosts of $\lint$ other than the pole boost, the boosts of $\lG$, the copy boost, and the
compact generators outside $\Cent(G_n)$, including the spatial rotations of the interior four-plane and the charge $J$ (fixed). The compact binding-compatible generators, and the copy rotation among them, never create quanta; the copy rotation
moves one (fixed). The number of Fock quanta is therefore not Lorentz invariant under any action of the tower, since every action
contains, on some constituent, the boosts of $\lint$ or of $\lG$. Under the common binding action an invariant number operator of the
specified polarisation type exists exactly at even rungs, where its quanta belong to two constituents at once (fixed); other actions with
spectator constituents can also admit one, as \cref{e2:thm:invnumber}(iii) shows. Whether a
dynamical mechanism of the arena supplies the non-binding-compatible or non-compact generator that creates a pair is open.

\paragraph{Up the ladder.} Rung by rung the following changes, each fixed. Under the common interior action $e^{2\pi S_3}=(-1)^n$, so a fermion
has vertices into odd rungs only. $\so{3n,3n}$ splits into $9n^2$ number-conserving and $9n^2-3n$ pair-creating directions, and the compact
number-conserving ones are the $\tfrac32n(3n-1)$ paired rotations. The copy rotations $\so n$ hop quanta among the constituents and the copy
boosts create pairs across them. An invariant number under the common binding action exists exactly at even $n$, uniquely up to
$N_K\mapsto6-N_K$ at $n=2$ and in several ways at $n=4$.

% --- source: rest_energy/d6_pair#0 ((head))
% Track d6 -- Duplication in motion: pair creation and the balance surface.
% Verification: verification/ch05_clifford_tower/d6_pair_balance.py (claims d6-0 ... d6-14c).

% --- source: rest_energy/d6_pair#1 (Duplication in motion: the balanced pair)
\section{Duplication in motion: the balanced pair}
\label{d6:sec:main}

Two further remarks are tested here: that the $b$-particle of \cref{sec:bfield} becomes decisive once two or more manifolds are
present, and that one should find the process by which a unitary (undivided) energy results in a duplication, from $\Cl(3,3)$ to
$\Cl(6,6)$. Duplication is taken literally: two copies of the arena whose tangent spaces add, $V_2=V\oplus V\cong\R^{6,6}$, which is grade
one of $\Cl(6,6)\cong\Cl(3,3)\hat\otimes\Cl(3,3)$ (equation~\cref{eq:tower}; the graded tensor product is inherited). It is followed in
motion, inside the closure of \cref{sec:closure}, and at a single field point. In brief:
(1) with one constituent and Lorentz-class inertia a constant nonzero rate forces a null momentum field, while with two or more a
codimension-two balance surface carries massive constituents at an exactly constant rate with bounded frames (\cref{d6:thm:balance},
restating \cref{prop:several} with the dimension count added). This is the precise sense in which a duplicate does what one copy
cannot, and it needs no $b$-field. (2) The closure never creates a pair (\cref{d6:cor:nopair}). (3) The admissible $b$-field between the
manifolds, the copy rotation, duplicates a field literally, and no Lorentz-covariant operation makes a massive pair out of a single null
field (\cref{d6:prop:covariant}); a natural candidate that rotates a single direction between the copies is a different, non-covariant
element (\cref{d6:ex:R1}).
(4) The rest-energy density of a duplicated field point is $M_{\rm pair}=\half(\sigma_1+\sigma_2)$, and the copy rotation can move all of it into
the two copies (\cref{d6:thm:pair,d6:ex:bw}); the kinematics of $\gamma\gamma\to e^+e^-$ is an analogy only. (5) The rest energy of each copy separately is
not invariant under the copy rotation, hence gauge dependent if that rotation is gauged, and with equal inertias the closure's own symmetry maps every balanced pair to one null constituent
(\cref{d6:prop:gauge}): pair creation is physical only where the copies are distinguished. (6) On the Dirac side the graded duplication adds
Dirac masses in quadrature (\cref{d6:prop:dirac}). (7) The zero-point energy stays zero, but an inter-manifold $b$ at unequal rates can
bring one line of the motion to zero exactly, and there the motion is unbounded (\cref{d6:prop:zpe}).

\paragraph{Conventions.} As in \cref{sec:restclosure} for one constituent. For two, $V_2=V\oplus V$ with coordinates $(X^{(1)};X^{(2)})$,
$\eta_2=\id_2\otimes\eta$, $G_2=\id_2\otimes G$; in the tower $e^{(1)}_a=e_a\otimes\id$, $e^{(2)}_a=\omega\otimes e_a$, $\omega=e_1\cdots e_6$,
and $V_2$ is spanned by the twelve $e^{(k)}_a$. $Z=(z_1,z_2)\in\C^2\otimes\C^3$, $z_k=z(X^{(k)})$, so $G_2$ acts as $i$; per copy
$E_6^{(k)},E_0^{(k)},\mathbf S^{(k)}$ and $q_k=z_k\cdot z_k$ are those of \eqref{en:defs}, and $q_{\rm tot}=q_1+q_2$. The Gram matrix
$Q_{kl}=z_k\cdot z_l$ (complex bilinear, symmetric) has singular values $\sigma_1\ge\sigma_2\ge0$; $H_{\rm sum}=\sum_kz_kz_k^\dagger$. The
\emph{copy rotation} is generated by $E:=\epsilon\otimes\id_6$, $\epsilon=\bigl(\begin{smallmatrix}0&-1\\1&0\end{smallmatrix}\bigr)$:
$R_\varphi=e^{\varphi E}$ maps $(z_1,z_2)\mapsto(\cos\varphi\,z_1-\sin\varphi\,z_2,\ \sin\varphi\,z_1+\cos\varphi\,z_2)$, identically in all six
directions; $E^2=-\id$, $[E,G_2]=0$, $E$ is antisymmetric for $\eta_2$ and for the Euclidean form. $T(z)$ is the stress tensor of
\cref{en:rainich}, $P(t)=T(z)t$ the density seen by a unit future timelike observer $t$, $P_{\rm sum}(t)=\sum_kT(z_k)t$, and
\[
M_{\rm pair}:=\inf_t|P_{\rm sum}(t)|,
\]
the rest-energy density of the duplicated field point: a local invariant of the summed stress tensor, not the invariant mass of two photons. In the closure the per-copy quantities are those of the momentum fields
$\pi^{(k)}=\eta^{-1}p^{(k)}$, $w_k=q(z(\pi^{(k)}))$; inertias are of Lorentz class, $N_k=(\nu_{k,1}\eta^{-1}+\nu_{k,2}G\eta^{-1})/m_k$, with
$c_k:=\nu_k/m_k\in\C$ and $R_{\rm eff}e^{i\varphi_{\rm eff}}:=\sum_kc_kw_k(0)$ (\cref{prop:several}). Two Lorentz actions occur: the
per-copy algebras $\lG^{(1)}\oplus\lG^{(2)}$, and the diagonal $\mathfrak l_{G,\Delta}$ ($k_B=2$ in \cref{thm:algebra}). Each $E_0^{(k)}$ is invariant under both; $M_{\rm pair}$, defined with one observer for both copies, is invariant under the diagonal action but not under the per-copy one, since a boost of one copy alone changes it. A constituent is called \emph{massive} when its momentum field is not null, $E_0(\pi^{(k)})\neq0$:
massive in the sense of the dictionary of \cref{s0:sec:question}, an interpretive reading that does not involve $mc^2=\hbar\kappa$. Rates are rates:
$mc^2=\hbar\kappa$ is an added identification (\cref{sec:restclosure}). Identifiers d6-$n$ refer to \texttt{d6\_pair\_balance.py}
(integrations DOP853, relative tolerance $10^{-12}$; d6-0 checks the conventions).

% --- source: rest_energy/d6_pair#2 (Mass with a constant rate needs a second constituent)
\subsection{Mass with a constant rate needs a second constituent}
\label{d6:sec:balance}

\begin{theorem}[Constant rate with massive constituents]
\label{d6:thm:balance}
Let one rotor act on $n$ constituents with Lorentz-class inertias and all $c_k\neq0$.
\begin{enumerate}[label=(\roman*),leftmargin=2.2em]
\item Along every solution $w_k(\tau)=e^{2i\theta(\tau)}w_k(0)$, so $R_{\rm eff}$ is constant. The rate $h/I_0$ is constant and nonzero iff
$R_{\rm eff}=0$ and $h(0)\neq0$; then $\theta=\kappa\tau$ and every frame moves on the circle
$x_k(\tau)=x_k(0)+\bigl(e^{2\kappa\tau G}-\id\bigr)(2\kappa G)^{-1}N_kp^{(k)}(0)$, of diameter $|N_kp^{(k)}(0)|/|\kappa|$.
\item For $n=1$, $R_{\rm eff}=0$ iff $E_0(\pi)=0$: no datum with a constant nonzero rate has a massive momentum field.
\item For $n\ge2$ the balance set $\mathcal B=\{\sum_kc_kw_k=0\}\subset\R^{6n}$ is, away from $0$, a smooth submanifold of codimension two; its
part on which every $E_0(\pi^{(k)})$ vanishes has dimension $4n<6n-2$, so the data with a constant nonzero rate and some $E_0(\pi^{(k)})>0$ are
open and dense in $\mathcal B\times\{h\neq0\}$. For $n=2$ both constituents are then massive, with $|c_1|E_0(\pi^{(1)})=|c_2|E_0(\pi^{(2)})$.
\end{enumerate}
\end{theorem}

\begin{proof}
(i) With a common rotor the momenta evolve as $p^{(k)}=e^{-\theta G}p^{(k)}(0)$ (\cref{prop:LP} and the closing remark of \cref{thm:closureunique}, $\mathfrak s_n=\R\ltimes_{G_n}V_n$), so $w_k=e^{2i\theta}w_k(0)$ (\cref{cl:frozen}) and, by \cref{thm:inertiareq}(ii) per
constituent, $U(\theta)=\sum_k\half p^{(k)\top}N_kp^{(k)}=\half\Re\bigl(e^{2i\theta}\sum_kc_kw_k(0)\bigr)$, whence
$I_0\ddot\theta=R_{\rm eff}\sin(2\theta+\varphi_{\rm eff})$. If $R_{\rm eff}=0$, $\ddot\theta=0$; if $R_{\rm eff}>0$ and $h$ is constant, then
$\sin(2\theta+\varphi_{\rm eff})\equiv0$, so $\theta$ is constant and $h=0$. Lorentz-class $N_k$ anticommutes with $G$, so
$\dot x_k=e^{\theta G}N_ke^{-\theta G}p^{(k)}(0)=e^{2\theta G}N_kp^{(k)}(0)$ (\cref{thm:drift}), which integrates to the circle.
(ii) $R_{\rm eff}=|c|\,|w(p)|=2|c|E_0(\pi)$. (iii) $F(p)=\sum_kc_kw(p^{(k)})$ has $dF_p(\delta)=2\sum_kc_k\,z(\pi^{(k)})\cdot z(\delta\pi^{(k)})$. At
$p\neq0$ some $z(\pi^{(k)})\neq0$, and $u\mapsto2c_kz(\pi^{(k)})\cdot u$ maps $\C^3$ onto $\C$, so $0$ is a regular value of $F$ on
$\R^{6n}\setminus0$. The null set $\{w=0\}\subset\R^6$ is the cone $|\bB|=|\bE|$, $\bB\cdot\bE=0$, of dimension four, and the all-null part of
$\mathcal B$ is its $n$-fold product. Points with $w(p^{(1)})=1$, $w(p^{(2)})=-c_1/c_2$ show that the massive part is not empty.
\end{proof}

Items (i) and (ii) restate \cref{prop:several} and \cref{cor:constrate} (fixed there); (iii) is fixed, the dimension count being added.
Checks: $\half p^{\!\top}Np=\Re(\nu w(p))/(2m)$ exactly, and the rank of $dF$ at random balanced points for $n=1,2,3$ (d6-1a,b); integrations of
the full two-constituent closure, rotor, momenta and frames (d6-2a--d; $I_0=1$, momenta drawn at random from the script's seeded generator, so
the values quoted are those of that sample). With equal inertias and $p^{(2)}=-Gp^{(1)}$ ($E_0(\pi^{(k)})=4.95$) the
rate stays constant to $1.6\times10^{-14}$ over $\tau\in[0,60]$ and the frames lie on the circles of (i) to $4\times10^{-12}$; a generic balanced
pair, and unequal inertias ($m=1,2.5$; $\nu=1+0.3i,\,0.4-0.7i$), give the same to $10^{-13}$. Off the surface (the unequal-inertia pair with $p^{(2)}$ perturbed at random, $R_{\rm eff}=0.0752$, rotating) the
rotor follows $I_0\ddot\theta=R_{\rm eff}\sin(2\theta+\varphi_{\rm eff})$ to $2\times10^{-10}$ and frame~1 drifts by $0.1795$ per half turn, exactly
additively (\cref{thm:freedrift} for several constituents).

\begin{remark}[The Lorentz requirement is what empties (ii)]
\label{d6:rem:lorentzreq}
$N=\operatorname{diag}(1,-1,0,-1,1,0)/m$ anticommutes with $G$ but is not $\lG$-invariant. For $p_0=e^1+e^2$, $z(\pi_0)=(1,1,0)$ and
$E_0(\pi_0)=1$, while $U(\theta)=\half p_0^{\!\top}e^{2\theta G}Np_0\equiv0$: the rate is exactly constant and the frame moves on a circle (d6-1c). One
constituent can thus combine a constant rate, a massive momentum field and a bounded, moving frame only through a cancellation between directions
that breaks Lorentz invariance: a bounded clamped frame needs $N_cp_0=0$, and a nonzero Lorentz-class $N_a$ then keeps $U$ constant only on null
data. The balance surface is the Lorentz-invariant form of the same cancellation, between constituents. Without the bounded frame no cancellation
is needed: a duality-class inertia, $[N,G]=0$, gives a constant rate for every datum, massive ones included, with $U=\half p_0^{\!\top}Np_0$ constant,
and the frame drifts at the velocity $Np_0$ (\cref{cl:flat}(iii) and \cref{thm:drift}; restated). Fixed, as a counterexample to
(ii) without the hypothesis.
\end{remark}

\begin{corollary}[The closure creates no pair]
\label{d6:cor:nopair}
For every Hamiltonian on $\mathfrak s_n^*$, including any coupling of the constituents through $H$, $p^{(k)}(\tau)=e^{-\theta(\tau)G}p^{(k)}(0)$ with
one angle $\theta$. Hence $E_6$, $E_0$ and $\mathbf S$ of each $\pi^{(k)}$, the Hermitian Gram matrix $\bar z_k\cdot z_l$ and the moduli $|Q_{kl}|$
of the momentum fields are constants of motion, and $Q\mapsto e^{2i\theta}Q$. No closure Hamiltonian turns a null constituent into a massive
one, fills an empty copy, or moves rest energy between constituents.
\end{corollary}

\begin{proof}
The Lie--Poisson equations of $\mathfrak s_n$ are \eqref{eq:restLPgen} with $G_n=\bigoplus_kG$; the first integrates to $p=e^{-\theta G_n}p_0$,
$\dot\theta=H_h$.
\end{proof}

Obstructed, for the momentum field and every $H$; it extends \cref{cl:frozen} to $n$ constituents. With a random Hamiltonian that couples
the two copies (quartic in $h$, trigonometric in $p$) these quantities stay constant to $2\times10^{-13}$ and $Q$ follows $e^{2i\theta}Q$ to the
same accuracy (d6-3). Inside the closure a pair is therefore a label of the initial data, carried and not generated, and the balance surface is
where such labels are compatible with a constant rate.

% --- source: rest_energy/d6_pair#3 (The balanced pairs)
\subsection{The balanced pairs}
\label{d6:sec:pairs}

\begin{proposition}[Balanced pairs]
\label{d6:prop:pairs}
Let $z_k=z(\pi^{(k)})$.
\begin{enumerate}[label=(\roman*),leftmargin=2.2em]
\item Equal inertias. With $a=\half(z_1-iz_2)$ and $b=\half(z_1+iz_2)$, $q_1+q_2=4\,a\cdot b$ and $q_1=a\cdot a+2a\cdot b+b\cdot b$. The balance
surface is the quadric $a\cdot b=0$, and on it $q_2=-q_1$ and $E_0^{(1)}=E_0^{(2)}=\half|a\cdot a+b\cdot b|$.
\item $EG_2=\epsilon\otimes G$ is an $\eta_2$-antisymmetric involution, the noncompact generator of $\mathfrak g(0,2,0)\cong\so{2,\C}$. Its
eigenspaces $V_\pm=\{X^{(2)}=\pm GX^{(1)}\}$ are six-dimensional and totally $\eta_2$-null, and $E=\mp G_2$ on $V_\pm$. A momentum field lies in
$V_+$ iff $p^{(2)}=-Gp^{(1)}$ iff $b=0$, and in $V_-$ iff $p^{(2)}=Gp^{(1)}$ iff $a=0$; so $V_\pm$ lie in the balance surface, and on them the
copy rotation $R_\varphi$ is the duality rotation $e^{\mp\varphi G_2}$ of both copies.
\item On $V_+$, $z_2=iz_1$: $w_2=-w_1$, and the two momentum fields have equal $E_6$, $\mathbf S$ and $E_0$; when $q_1>0$ the first is of
magnetic and the second of electric type. $T(z_2)=T(z_1)$, so the summed stress tensor is $2T(z_1)$ and $M_{\rm pair}=2E_0^{(1)}$, although
$q_{\rm tot}=0$.
\item Unequal inertias. If $q_1\neq0$ the balanced pairs are $z_2=\rho\Lambda z_1$ with $\Lambda\in SO(3,\C)$ and $\rho^2=-c_1/c_2$, and
$|c_1|E_0^{(1)}=|c_2|E_0^{(2)}$.
\end{enumerate}
\end{proposition}

\begin{proof}
(i) Expand. (ii) $(\epsilon\otimes G)^2=\epsilon^2\otimes G^2=\id$ and $(\epsilon\otimes G)^\dagger=\epsilon^{\!\top}\otimes G^\dagger=-\epsilon\otimes G$.
The eigenvalue $+1$ means $-GX^{(2)}=X^{(1)}$, i.e.\ $X^{(2)}=GX^{(1)}$, and then
$\eta_2(X,Y)=\eta(X^{(1)},Y^{(1)})+\eta(GX^{(1)},GY^{(1)})=0$ since $G$ is an anti-isometry; on $V_+$,
$E(X^{(1)},X^{(2)})=(-GX^{(1)},X^{(1)})=-G_2(X^{(1)},X^{(2)})$. Finally $\eta^{-1}p^{(2)}=G\eta^{-1}p^{(1)}$ iff $p^{(2)}=-Gp^{(1)}$, as $\eta G=-G\eta$, and
$z_2=iz_1$ iff $b=0$. (iii) $T$ is duality invariant (\cref{en:rainich}), and \cref{en:rainich} gives $|2T(z_1)t|=2E_0^{(1)}$ for
every $t$. (iv) $c_1q_1+c_2\rho^2q(\Lambda z_1)=c_1(q_1-q_1)=0$; conversely $q(z_2/\rho)=q_1\neq0$, and $SO(3,\C)$ is transitive on each complex
sphere $\{q=q_1\}$ (\cref{en:orbits}).
\end{proof}

The first sentence of (iii) restates the last sentence of \cref{prop:several} (fixed there); (i), (ii) and (iv) are fixed (that
$V_\pm$ are null, and are exactly where the copy rotation is a duality rotation, could have failed); ``magnetic and electric pair'' and
``duality-balanced'' are readings. Checked exactly (d6-4a--c) and numerically (d6-4d). So the balanced pair of \cref{prop:several} is
a magnetic and an electric constituent a quarter duality turn apart, with equal rest energies, whose duality phases the rotor turns together.
Its total $q$ vanishes, but $q_{\rm tot}$ is not the rest energy of the pair: its summed density has rest energy $2E_0^{(1)}$ for every
observer. Thus $\half|q_{\rm tot}|$ is what the rotor sees ($R_{\rm eff}=|\nu|\,|q_{\rm tot}|/m$ for equal inertias) and $M_{\rm pair}\ge\half|q_{\rm tot}|$ is
what an observer of the two copies sees (\cref{d6:thm:pair}). The block-diagonal inertia is the per-copy Lorentz requirement: the invariant
symmetric forms on $V_2^*$ number four under $\lG^{(1)}\oplus\lG^{(2)}$, six under $\mathfrak l_{G,\Delta}$ alone (cross terms,
$U=\half\Re\tr(CQ(\pi))$, $C$ complex symmetric, balance $\tr(CQ(\pi))=0$, again of codimension two) and two under $\Cent(G_2)\cap\so{6,6}$, all of
Lorentz class (d6-4e; the $n=2$ analogue of \cref{thm:inertia3}).

% --- source: rest_energy/d6_pair#4 (What a copy rotation can and cannot duplicate)
\subsection{What a copy rotation can and cannot duplicate}
\label{d6:sec:covariant}

\begin{proposition}[Lorentz-covariant duplication]
\label{d6:prop:covariant}
\begin{enumerate}[label=(\roman*),leftmargin=2.2em]
\item The real linear maps of $V_2$ commuting with $\mathfrak l_{G,\Delta}$ are the $A\otimes\id_3$, $A\in M_2(\C)$, acting on $Z$ with $G_2=i$. Inside
$\so{6,6}$ they form $\operatorname{span}\{E,EG_2\}\cong\so{2,\C}$, which is $\mathfrak g(0,2,0)$ of \cref{thm:algebra}; its compact part
is $\{R_\varphi\}$.
\item $A\otimes\id_3$ maps $Q\mapsto AQA^{\!\top}$ and $Z\mapsto ZA^{\!\top}$ ($Z$ the $3\times2$ matrix of columns $z_k$). Hence $R_\varphi$ preserves
$H_{\rm sum}$, so every summed density $P_{\rm sum}(t)$, and $q_{\rm tot}$, $\det Q$, $\sigma_1$, $\sigma_2$. The unitary $A$, the group $U(2)$ of
energy-preserving duality rotations of two free Maxwell fields (end of \cref{sec:spintower}; \cite{GaillardZumino81}), preserve
$H_{\rm sum}$, $\sigma_1$ and $\sigma_2$ but not $q_{\rm tot}$, and $\uu2\cap\so{6,6}=\R E$.
\item Every map $F:\C^6\to\C^6$, linear or not, commuting with the diagonal $SO(3,\C)$ sends $(z,0)$ with $q(z)=0\neq z$ to $(\alpha z,\beta z)$:
two null constituents. In particular $R_\varphi(z,0)=(\cos\varphi\,z,\sin\varphi\,z)$.
\end{enumerate}
\end{proposition}

\begin{proof}
(i) $V$ is an irreducible $\lG$-module with commutant $\operatorname{span}\{\id,G\}\cong\C$ (\cref{prop:lGhoms}), so the commutant
of $\mathfrak l_{G,\Delta}$ on $\R^2\otimes V$ is $M_2(\R)\otimes\operatorname{span}\{\id,G\}=M_2(\C)$; antisymmetry for $\eta_2$ leaves $A\in\C\epsilon$, which
is \cref{prop:commutantsn} (row ``diagonal $\lG$'', dimension $2$). (ii) $Q=Z^{\!\top}Z$ and $H_{\rm sum}=ZZ^\dagger$, and $T$ is linear in
$H$ (\cref{en:split}). (iii) For $\Lambda\in SO(3,\C)$ fixing $z$, $F(z,0)=F(\Lambda z,0)=\Lambda F(z,0)$. In the action $u\mapsto w\times u$
of $\so{3,\C}$ (\eqref{bf:eq:w}) the stabiliser algebra of a null $z$ is $\{w:w\times z=0\}=\C z$, whose common fixed vectors satisfy
$z\times v=0$, i.e.\ $v\in\C z$.
\end{proof}

Item (i) restates \cref{prop:commutantsn} and \cref{thm:algebra} (fixed there), (ii) is inherited linear algebra and restates the remark
on $U(n)$ at the end of \cref{sec:spintower}, and (iii) is an obstruction; all checked exactly (d6-5a--c). This settles in the negative the reading of the copy rotation
as a pair-creating operation. The copy rotation acting identically in all six directions is the admissible $b$-field between the two manifolds,
and it duplicates a field literally: the undivided energy and the rest energy of $(z,0)$ are shared in the same ratio
$\cos^2\varphi:\sin^2\varphi$, a null field becomes two null fields and a massive one two massive replicas. It cannot produce a balanced massive
pair from a single null field, and no other Lorentz-covariant operation can.

\begin{theorem}[The rest-energy density of a duplicated field point]
\label{d6:thm:pair}
Let $z_1,z_2\in\C^3$, with Gram matrix $Q$ of singular values $\sigma_1\ge\sigma_2$.
\begin{enumerate}[label=(\roman*),leftmargin=2.2em]
\item $T(z_1)\cdot T(z_2)=|z_1\cdot z_2|^2$, and as mixed tensors $\bigl(T(z_1)+T(z_2)\bigr)^2=\tfrac14\|Q\|_F^2\,\id+T(z_1\times z_2)$, with
$(z_1\times z_2)\cdot(z_1\times z_2)=\det Q$.
\item $|P_{\rm sum}(t)|^2=\tfrac14\|Q\|_F^2+u_t(z_1\times z_2)$ for every observer $t$, $u_t(w)$ being the energy density of $w$ seen by $t$. Hence
$M_{\rm pair}=\half(\sigma_1+\sigma_2)$, attained exactly in the rest frames of $z_1\times z_2$ when it is not null.
\item In every frame, $E_0^{(1)}+E_0^{(2)}\le M_{\rm pair}\le|P_{\rm sum}(t)|\le E_6^{(1)}+E_6^{(2)}$, and $M_{\rm pair}\ge\half|q_{\rm tot}|$.
\item $\max_\varphi\bigl(E_0^{(1)}+E_0^{(2)}\bigr)(R_\varphi Z)=M_{\rm pair}$.
\item For two null fields $M_{\rm pair}=|z_1\cdot z_2|=\bigl(E_6^{(1)}E_6^{(2)}\bigr)^{1/2}(1-\cos\vartheta_{12})$, $\vartheta_{12}$ the angle between
their Poynting vectors; for $(z,0)$ with $z$ null, $M_{\rm pair}=0$.
\end{enumerate}
\end{theorem}

\begin{proof}
(i) Polynomial identities, verified symbolically in twelve real variables (d6-6a); in spinor form \cite{PenroseRindler84}, $T=\sum_k\phi_k\bar\phi_k$ and
$\phi^k_{AB}\phi^{l\,B}{}_C=\half(\phi^k\cdot\phi^l)\epsilon_{AC}+\chi^{kl}_{AC}$ with $\chi^{kl}$ antisymmetric in $(k,l)$, so the cross terms of $T^2$
pair the symmetric $Q$ with the antisymmetric $\chi$ and drop out. The last identity is the complex Lagrange identity. (ii) $T$ is
$g$-self-adjoint, so $|P(t)|^2=t\cdot T^2t$; insert (i). By \cref{en:min}, $u_t(w)\ge E_0(w)=\half|\det Q|$, with equality iff $t$ is a
rest frame of $w$, and $\tfrac14\|Q\|_F^2+\half|\det Q|=\tfrac14(\sigma_1+\sigma_2)^2$. For null $w\neq0$ the infimum is approached by Doppler boosts;
for $w=0$, $|P(t)|=\half\sigma_1$ for all $t$. (iii) Each $P_k(t)$ is future causal of length $E_0^{(k)}$ (\cref{en:rainich}), future causal
vectors obey $|P_1+P_2|\ge|P_1|+|P_2|$ and $|P|\le P^0$, and $\sigma_1+\sigma_2\ge|\tr Q|$. (iv) Write $q_{1,2}(\varphi)=t\pm d(\varphi)$ and
$Q_{12}(\varphi)=r(\varphi)$ with $t=\half q_{\rm tot}$. The pair $(d,r)$ turns by $2\varphi$, $d^2+r^2=t^2-\det Q$ is fixed, and
$(|q_1|+|q_2|)^2=2|t|^2+2|d|^2+2|t^2-d^2|\le2|t|^2+2|d|^2+2|r|^2+2|\det Q|=(\sigma_1+\sigma_2)^2$, by $|t^2-d^2|=|\det Q+r^2|\le|\det Q|+|r|^2$.
Equality holds exactly when $\overline{\det Q}\,r(\varphi)^2\ge0$, hence for every $\varphi$ when $\det Q=0$; as $\varphi$ varies, $r(\varphi)$ runs over an ellipse centred at $0$ or a
segment through $0$, so such $\varphi$ exists. (v) $Q=\bigl(\begin{smallmatrix}0&r\\r&0\end{smallmatrix}\bigr)$, $r=z_1\cdot z_2$; and
$T(z_k)=k_k\otimes k_k$ with $k_k$ future null, $k_k^0=(E_6^{(k)})^{1/2}$, so $|r|=k_1\cdot k_2$ by (i), and $|\mathbf S^{(k)}|=E_6^{(k)}$.
\end{proof}

Items (i), (iii) and (v) are inherited: ordinary electrodynamics of two field strengths, with the reverse triangle inequality for future
causal momentum densities. The closed form of $M_{\rm pair}$ in (ii) and its saturation by the copy rotation in (iv) are fixed; their priority is
not established (compare Kaiser's local invariants \cite{Kaiser11} and Thompson's results on complex symmetric matrices \cite{Thompson79}). Checked symbolically (i), at the rest frames of $z_1\times z_2$ to $10^{-10}$ against $300$ random
observers (ii), by maximising over $\varphi$ (relative deviation $5\times10^{-16}$) (iv), and on random and null pairs (iii), (v) (d6-6a--d).
Since $T(z_1\times z_2)$ is traceless, the trace of (i) gives the pair invariant of \cref{d5:rem:twoE0},
$E_{\rm pair}=(\tfrac14\tr T_{\rm tot}^2)^{1/2}=\half\|Q\|_F$, and (ii) gives $M_{\rm pair}^2=E_{\rm pair}^2+\half|\det Q|$. The two agree exactly when
$\det Q=0$, which holds on every duplicate $(\alpha z,\beta z)$ of a single field (\cref{d5:thm:dup}); for the counter-propagating null pair of
\cref{d6:ex:bw}, $E_{\rm pair}=\sqrt2$ while $M_{\rm pair}=2$. $M_{\rm pair}$ is the local rest-energy density of the summed stress tensor at a point,
not the invariant mass of two photons. For two null fields it equals $(u_1u_2)^{1/2}(1-\cos\vartheta_{12})$ by (v), whereas two photons of
energies $E_1,E_2$ have invariant mass $\bigl(2E_1E_2(1-\cos\vartheta_{12})\bigr)^{1/2}$; for $z_1=(1,i,0)$ and $z_2=(0,1,i)$ the values are $1$ and
$\sqrt2$, and the counter-propagating example of \cref{d6:ex:bw} agrees only by coincidence.

\begin{example}[Two counter-propagating null fields]
\label{d6:ex:bw}
$z_1=(1,i,0)$ and $z_2=(1,-i,0)$ are null with $\mathbf S^{(1)}=-e_3=-\mathbf S^{(2)}$, so $\mathbf S_{\rm sum}=0$, $E_6^{\rm tot}=2$,
$Q=\bigl(\begin{smallmatrix}0&2\\2&0\end{smallmatrix}\bigr)$ and $M_{\rm pair}=2$. The rotation $R_{\pi/4}$ gives $z_1'=(0,\sqrt2\,i,0)$ and
$z_2'=(\sqrt2,0,0)$: a pure electric and a pure magnetic field, both at rest, with $q=-2,+2$ and $E_0=1$ each. In this frame, the centre-of-momentum frame, the chain of
\cref{d6:thm:pair}(iii) holds with equality for the rotated pair, $E_0^{(1)}+E_0^{(2)}=M_{\rm pair}=|P_{\rm sum}|=E_6^{\rm tot}=2$, while
$M_{\rm pair}=2>\half|q_{\rm tot}|=0$: the rotated pair is at rest in the centre-of-momentum frame and takes the whole of $M_{\rm pair}$ (d6-7).
The kinematics of $\gamma\gamma\to e^+e^-$ \cite{BreitWheeler34} is an analogy, not a production process. \Cref{d6:thm:pair} adds the two stress
tensors, $T(z_1)+T(z_2)$, which differs in general from the tensor $T(z_1+z_2)$ of the coherent superposition (here $\operatorname{diag}(2,0,0,2)$ against
$\operatorname{diag}(2,-2,2,2)$); the coherent superposition is, by analogy, the standing
wave of
\cref{en:whole}(3).
\end{example}

\begin{example}[The rotation of one direction]
\label{d6:ex:R1}
Let $R^{(1)}_\varphi$ rotate the planes $(e^{(1)}_1,e^{(2)}_1)$ and $(e^{(1)}_4,e^{(2)}_4)$ through $\varphi$. It is compact, commutes with $G_2$ and lies in
$\Cent(G_2)\cap\so{6,6}\cong\so{6,\C}$, the $n=2$ analogue of \cref{thm:so9C}. On $((1,i,0),0)$ it gives $z_1'=(\cos\varphi,i,0)$,
$z_2'=(\sin\varphi,0,0)$: $q_1'=-\sin^2\varphi$, $q_2'=\sin^2\varphi$, a balanced pair with $E_0=\half\sin^2\varphi$ each. $E_6^{\rm tot}=1$ and $q_{\rm tot}=0$
are conserved, $\mathbf S_{\rm sum}$ goes from $-e_3$ to $-\cos\varphi\,e_3$ and $M_{\rm pair}$ from $0$ to $|\sin\varphi|$, so
$E_6^{\rm tot}=1\ge M_{\rm pair}=|\sin\varphi|\ge\sum_kE_0^{(k)}=\sin^2\varphi$, with equality throughout at $\varphi=\pi/2$ (a magnetic field along $e_1$ and
an electric one along $e_2$, at rest). $R^{(1)}$ does not commute with $\mathfrak l_{G,\Delta}$, and lies neither in $\mathfrak g(0,2,0)$ nor in
$\mathfrak g(0,1,1)=\so{3,3}$ (Lorentz algebra on copy~1, copy~2 a spectator): it is a candidate gauge transformation under no Lorentz action of
\cref{thm:algebra} with $k_I=0$ (d6-8), and under none with $k_I>0$ either: it lies outside $\mathfrak g(2,0,0)$ and
$\mathfrak g(1,0,1)=\so{5,3}$ (x6-8a), and it has components from the interior four-plane of one copy into the binding frame of the other,
which the proof of \cref{thm:algebra} excludes from $\mathfrak g(1,1,0)$. It is therefore non-covariant under all five Lorentz actions at $n=2$. Read kinematically it is an infinitely heavy third body, absorbing momentum and no energy, and
$E_6\ge E_0^{(1)}+E_0^{(2)}$ is the field-point form of the threshold $E_\gamma\ge2mc^2$ for pair production in the field of a nucleus
\cite{BetheHeitler34}. The numbers and the non-covariance are fixed (d6-8; x6-8a,b), the third-body reading is interpretive, the threshold
kinematics is inherited, and the agent is constructed: it exists, and nothing in the arena supplies it.
\end{example}

A tempting reading of pair creation applies to $((1,i,0),0)$ a rotation said to mix the copies ``identically in all six directions''. The rotation
that produces the balanced pair above is $R^{(1)}_\varphi$, which rotates two directions only, not the copy rotation $R_\varphi$. For $R^{(1)}_\varphi$ the
conservation of $E_6^{\rm tot}$ and $q_{\rm tot}$ and the non-conservation of $\mathbf S_{\rm sum}$ hold, while the copy rotation, which does act identically
in all six directions, conserves $\mathbf S_{\rm sum}$ and the whole summed stress tensor.

% --- source: rest_energy/d6_pair#5 (Is the rest energy of one copy observable?)
\subsection{Is the rest energy of one copy observable?}
\label{d6:sec:gauge}

\begin{proposition}[Gauge status of the per-copy rest energy]
\label{d6:prop:gauge}
\begin{enumerate}[label=(\roman*),leftmargin=2.2em]
\item $R_\varphi$ is admissible for the diagonal action: compact, commuting with $G_2$ and $\mathfrak l_{G,\Delta}$, and the compact part of
$\mathfrak g(0,2,0)$. On $S_2$ it is generated by $M=\half\sum_a\eta_{aa}e^{(1)}_ae^{(2)}_a$, with $\ad M=-E$ on $V_2$ and eigenvalues
$i\{-3,\dots,3\}$; $U_\varphi=\prod_a\bigl(\cos\tfrac\varphi2+\sin\tfrac\varphi2\,\eta_{aa}e^{(1)}_ae^{(2)}_a\bigr)$ has $U_{2\pi}=\id$, so the copy rotation
acts through $SO(2)$; and $U_{\pi/2}J^{(1)}U_{\pi/2}^{-1}=J^{(2)}$, $J^{(k)}=e^{(k)}_2e^{(k)}_3$.
\item $R_\varphi$ preserves $H_{\rm sum}$, $q_{\rm tot}$, $\det Q$, $\sigma_{1,2}$, $M_{\rm pair}$ and the total charge $J^{(1)}+J^{(2)}$; it does not
preserve $E_0^{(k)}$, $q_k$ or the per-copy charges (\cref{d6:ex:bw}, (i)). On $V_\pm$ it is a duality rotation of both copies, and there the
$E_0^{(k)}$ are invariant.
\item Equal inertias. The two-constituent closure is invariant under $O(\eta_2)\cap\Cent(G_2)\cong O(6,\C)$, whose compact part, the real
orthogonal group of $\C^6$, acts transitively on $\{q_{\rm tot}=0,\ |Z|=r\}$: every balanced pair is carried to one null constituent and an empty
copy at the same $E_6^{\rm tot}$; for instance $(e_1,ie_1)\mapsto((1,i,0),0)$, with $M_{\rm pair}$ going from $1$ to $0$.
\item If $c_1\neq c_2$, no nonzero element of $\mathfrak g(0,2,0)$ preserves the inertia.
\end{enumerate}
\end{proposition}

\begin{proof}
(i) Each term pairs two generators of the same square, hence is compact (\cref{thm:sig}), squares to $-\tfrac14$, and the six terms
commute; exponentiate term by term. (ii) \Cref{d6:prop:covariant}(ii) and \cref{d6:prop:pairs}(ii). (iii) Every $\Lambda$ commuting with $G_2$
preserves the Lie--Poisson structure, and for $\Lambda\in O(\eta_2)$, $\Lambda\eta_2^{-1}\Lambda^{\!\top}=\eta_2^{-1}$ and
$\Lambda G_2\eta_2^{-1}\Lambda^{\!\top}=G_2\eta_2^{-1}$, so the equal Lorentz-class inertia is invariant. A nonzero $Z$ with $q_{\rm tot}(Z)=0$ is $u+iv$
with $u,v\in\R^6$, $|u|=|v|$, $u\perp v$, and $SO(6)$ is transitive on such pairs. (iv) $sE+tEG_2$ acts on the copy index as $(s+it)\epsilon$, and
$\epsilon^{\!\top}\operatorname{diag}(c_1,c_2)+\operatorname{diag}(c_1,c_2)\,\epsilon=(c_2-c_1)\bigl(\begin{smallmatrix}0&1\\1&0\end{smallmatrix}\bigr)$, which must vanish for the form $\sum_kc_k\,z_k\cdot z_k$ to be invariant.
\end{proof}

Item (i) restates \cref{thm:algebra} and \cref{prop:orthospin} (fixed there), with the lift computed; (ii)--(iv) are fixed; checked in integer
arithmetic (d6-9a, d6-9b, d6-9c, d6-10). This decides whether ``mass creation in pairs'' is physical. \emph{If the copy rotation is gauged}, the rest energy of
each copy is not an observable: \cref{d6:ex:bw} is a change of gauge, not a process, and two counter-propagating null fields and an
electric--magnetic pair at rest are one physical state, whose gauge-invariant content $M_{\rm pair}=2$ was there before. \emph{In the
equal-inertia closure alone}, its own compact symmetry maps every balanced pair to one null constituent, changing $M_{\rm pair}$ as it does so (by (iii)), and the rotor sees only
$\half|q_{\rm tot}|=0$; the rest-energy density $M_{\rm pair}$ of the pair is invisible to the rate. $M_{\rm pair}$ becomes an invariant only through the diagonal Lorentz action, and the per-copy rest energies only through a per-copy one,
neither of which the closure selects (the Lorentz action at several constituents is left open, \cref{tab:choices}; at $n=3$ the closure symmetry meets the interior Lorentz algebra of the first constituent in $0$, \cref{thm:nocommon}, and \cref{d5:prop:so6C}(iv) is the $n=2$ analogue). \emph{If the copies are
distinguished} --- by unequal inertias, which by (iv) break $\mathfrak g(0,2,0)$, by unequal binding rates, by a projector, or by a Lorentz
action on one copy only, whose algebra $\mathfrak g(0,1,1)$ contains no copy mixing --- the per-copy rest energies are observables, but then the
closure creates no pair (\cref{d6:cor:nopair}). In the interior with $S=\operatorname{diag}(\kappa_1,\kappa_2)\otimes G+\beta E$ the momentum field
obeys $\dot\pi=S^\dagger\pi$, i.e.\ $z(\pi)(\tau)=e^{ih\tau}z(\pi)(0)$ with $h=\bigl(\begin{smallmatrix}\kappa_1&-i\beta\\i\beta&\kappa_2\end{smallmatrix}\bigr)$ on the
copy index. For $\kappa_1=\kappa_2$ this is $e^{i\kappa\tau}R_{-\beta\tau}$, and the per-copy rest energies change only by the gauge rotation
(Larmor \cite{Larmor1897}; \cref{sec:bfield}). For $\kappa_1\neq\kappa_2$ they change physically: for the pair of \cref{d6:ex:bw},
$E_0(\pi^{(1)})(\tau)=|z_1\cdot z_2|\,|U_{11}U_{12}|$, $U=e^{ih\tau}$, which for $\kappa=(1,3)$, $\beta=0.4$ peaks at $0.6897$, the Rabi value
$2\sqrt{s(1-s)}$ with $s=4\beta^2/((\kappa_2-\kappa_1)^2+4\beta^2)$, at fixed $E_6^{\rm tot}$; from $(z,0)$ with $z$ null both copies stay null for all
$\tau$ (d6-11). This is two-level mixing \cite{Rabi37}, inherited. It changes the per-copy rest energies but not $M_{\rm pair}$, so their sum stays at most
$M_{\rm pair}$ (\cref{d6:thm:pair}(iii)) without being conserved in general: for the pair of \cref{d6:ex:bw} both copies gain the same rest
energy, so none passes from one copy to the other, and from two null fields rest energy appears only where the second is non-collinear and
supplies a nonzero $M_{\rm pair}$.

\begin{remark}[A copy rotor]
\label{d6:rem:copyrotor}
The uniqueness argument of \cref{thm:closureunique} uses that equivariant maps between distinct constituents vanish, which holds for
$\lG^{(1)}\oplus\lG^{(2)}$ but not for $\mathfrak l_{G,\Delta}$: $E$ is $\mathfrak l_{G,\Delta}$-equivariant, commutes with $G_2$ and rotates $p$ without dilating it, so a
second rotor $\hat c$ with $[\hat c,X]=EX$, $[\hat c,\hat g]=0$ gives an equivariant closure $\R^2\ltimes V_2$, with Lie--Poisson equations
$\dot p=(H_hG_2^{\!\top}+H_{h_c}E^{\!\top})p$, $\dot h=-p^{\!\top}G_2\nabla H$, $\dot h_c=-p^{\!\top}E\nabla H$. With equal Lorentz-class inertias
$\dot h_c=0$, the copy angle turns uniformly, and the per-copy $E_0(\pi^{(k)})$ oscillate while the summed momentum tensor, $M_{\rm pair}(\pi)$ and
$|q_{\rm tot}(\pi)|$ stay constant; $((1,i,0),0)$ stays null--null (d6-12). The inter-manifold $b$-field made dynamical changes
the rest energies of the two manifolds at fixed pair invariants. Their sum is not conserved in general (for the data of \cref{d6:ex:bw} the copy
rotation takes it from $0$ to $2$) and never exceeds $M_{\rm pair}$ (\cref{d6:thm:pair}(iii)), so the rotor creates per-copy rest energy but no
pair invariant. Constructed.
\end{remark}

% --- source: rest_energy/d6_pair#6 (The Dirac side)
\subsection{The Dirac side}
\label{d6:sec:dirac}

\begin{proposition}[Pairs of Dirac states in the tower]
\label{d6:prop:dirac}
On $S_2=\R^8\otimes\R^8$ with the tower generators:
\begin{enumerate}[label=(\roman*),leftmargin=2.2em]
\item $J^{(1)}=J\otimes\id$ and $J^{(2)}=\id\otimes J$ (the Klein factors cancel on bivectors); they commute with each other and with the eight
interior gammas $e^{(k)}_a$, $a\in\{1,4,5,6\}$. With an external $i$, $Q_{\rm tot}=-i(J^{(1)}+J^{(2)})$ has spectrum $\{-2,0,2\}$ with
multiplicities $16,32,16$.
\item Let $\psi,\phi\in S\otimes\C$ be particles at rest, $e_1\psi=\psi$, $J\psi=i\psi$ ($H(0)=me_1$, equation~\cref{eq:Hp}), and $C=e_2$
(\cref{thm:C}). Then $\Psi=\psi\otimes C\phi$ has charges $(1,-1)$, $Q_{\rm tot}\Psi=0$, and
$-i\bigl(JH(0)\otimes\id+\id\otimes JH(0)\bigr)\Psi=2m\Psi$.
\item $C^{(1)}C^{(2)}=e^{(1)}_2e^{(2)}_2$ anticommutes with $J^{(1)}$ and $J^{(2)}$, but $(C^{(1)}C^{(2)})^2=-\id$: no state is mapped to $\pm$ itself.
The Klein-corrected $C\otimes C=-\omega^{(1)}C^{(1)}C^{(2)}$, $\omega^{(1)}=\omega\otimes\id$, squares to $\id$, and $\Psi_\pm=\psi\otimes C\psi\pm C\psi\otimes\psi$
satisfy $(C\otimes C)\Psi_\pm=\pm\Psi_\pm$ and $Q_{\rm tot}\Psi_\pm=0$.
\item The tower's mass terms anticommute: $(e^{(1)}_1+e^{(2)}_1)^2=2\,\id$, so $m(e^{(1)}_1+e^{(2)}_1)$ has spectrum $\{\pm\sqrt2\,m\}$, whereas the
two-particle operator $m(e_1\otimes\id+\id\otimes e_1)=m(e^{(1)}_1+\omega^{(1)}e^{(2)}_1)$ has spectrum $\{-2m,0,2m\}$, multiplicities $16,32,16$.
\item Under the diagonal interior action ($k_I=2$) a rotation through $2\pi$ is $+\id$ on $S_2$.
\end{enumerate}
\end{proposition}

\begin{proof}
(i) $(\omega\otimes e_2)(\omega\otimes e_3)=\id\otimes J$ as $\omega^2=\id$; the commutations are the blade rule (equation~\cref{eq:bladecomm});
$(J^{(1)}+J^{(2)})^2=-2(\id-J^{(1)}J^{(2)})$, where $\half(\id-J^{(1)}J^{(2)})$ is a projector of trace $32$. (ii) $C$ anticommutes with $e_1$ and $J$, and
$J$ commutes with $H(0)$. (iii) Generators of distinct constituents anticommute and $e_2^2=\id$; $\id\otimes e_2=\omega^{(1)}e^{(2)}_2$. (iv) The
same anticommutation; $e_1\otimes\id$ and $\id\otimes e_1$ are commuting traceless involutions with traceless product. (v) Each copy's rotation
through $2\pi$ is $-\id$, and the two commute.
\end{proof}

Inherited: Klein--Jordan--Wigner signs \cite{JordanWigner28}, the charge-conjugation parity of a particle--antiparticle pair
\cite{WeinbergQFT}, and integer spin for two spin-$\half$ constituents, which is \cref{prop:orthospin} at $k_I=2$ (restated, fixed
there); checked exactly (d6-13a--e). Two points are specific to the tower. The total charge is invariant under the copy rotation while
$U_{\pi/2}$ exchanges $J^{(1)}$ and $J^{(2)}$ (\cref{d6:prop:gauge}(i)), so which copy carries the particle is not invariant under the copy rotation, hence gauge dependent if that rotation is gauged, exactly as the
per-copy rest energy is. And by (iv) the graded duplication $\Cl(3,3)\hat\otimes\Cl(3,3)=\Cl(6,6)$ adds Dirac masses in quadrature, like two
anticommuting Clifford directions --- one particle of mass $\sqrt2\,m$ --- while the pair threshold $2m$ appears only for Klein-corrected,
ungraded operators: duplication in the tower's sense is not pair creation.

% --- source: rest_energy/d6_pair#7 (Zero-point energy of the duplicated frame)
\subsection{Zero-point energy of the duplicated frame}
\label{d6:sec:zpe}

\begin{proposition}[Zero-point energy]
\label{d6:prop:zpe}
\begin{enumerate}[label=(\roman*),leftmargin=2.2em]
\item For $\mathcal L=\sum_k\frac{m_k}2\eta\bigl(\dot X^{(k)}+\kappa_kGX^{(k)},\dot X^{(k)}+\kappa_kGX^{(k)}\bigr)$ each rate $\kappa_k$ carries Krein
signature $(3,3)$; with positive-norm quantisation the zero-point energy is $\sum_k\frac{\hbar\kappa_k}2(3-3)=0$ for all $\kappa_1,\kappa_2>0$ and $m_1,m_2\neq0$. At $\kappa_k=0$ the copy is free, its phase-space generator is not semisimple, and no zero-point energy is defined.
\item With an inter-manifold $b=\beta E$ ($m_1=m_2$, $\kappa_k\ge0$ not both zero), $S=\operatorname{diag}(\kappa_1,\kappa_2)\otimes G+\beta E$ has spectrum
$\pm i\omega_\pm$, $\omega_\pm=\half(\kappa_1+\kappa_2)\pm\bigl(\tfrac14(\kappa_1-\kappa_2)^2+\beta^2\bigr)^{1/2}$, each three times; these are the lines of \cref{d5:thm:bfield2}(ii). Every solution is
bounded unless $\beta^2=\kappa_1\kappa_2$ and $\kappa_1\neq\kappa_2$; when bounded, each nonzero-frequency line (coincident lines combined) has
Krein signature $(3,3)$ and their zero-point energy vanishes. At $\kappa_1=\kappa_2=|\beta|>0$ the zero line is semisimple and the motion
bounded, but it is not an oscillator line. At $\kappa_1=0$, $|\omega_-|\approx\beta^2/\kappa_2$.
\item At $\beta^2=\kappa_1\kappa_2$, $\kappa_1\neq\kappa_2$, $\omega_-=0$ and the phase-space generator is not semisimple: some solutions grow linearly
in $\tau$, and on that line there is no oscillator vacuum.
\end{enumerate}
The cancellation in (i) and (ii) concerns the nonzero-frequency lines of a Hamiltonian that is unbounded below; no vacuum-energy conclusion
follows from it (compare \cref{thm:EBsplit}).
\end{proposition}

\begin{proof}
(i) \Cref{mb:noground} per copy, or \cref{bf:krein}. (ii) On the $G_2=i$ sector $S$ acts on the copy index as
$i\operatorname{diag}(\kappa_1,\kappa_2)+\beta\epsilon=i(K-i\beta\epsilon)$ with $K-i\beta\epsilon$ Hermitian. $S^\dagger=K\otimes G-\beta E$ commutes with
$G_2$, so for $\mu\neq0$ its eigenspaces lie in the $G_2=\pm i$ sectors, which are totally isotropic for $h$ (a vector with $GY=iY$ has
$\bE=-i\bB$); \cref{bf:criterion}(ii) then gives boundedness except possibly at $\omega_-=0$, that is at $\beta^2=\kappa_1\kappa_2$; there the kernel of $S^\dagger$ computed in (iii) is $\eta_2$-null exactly when $\kappa_1=\kappa_2$, so the case $\kappa_1=\kappa_2=|\beta|$ is bounded as well. \Cref{bf:krein} gives the signatures. (iii) At
$\beta^2=\kappa_1\kappa_2$ the kernel of $S^\dagger$ consists of the $(v_1c,v_2c)$, $c\in\C^3$, with $(v_1,v_2)$ a nonzero multiple of $(i\kappa_2,\beta)$ if $\kappa_2>0$ and of $(0,1)$ if $\kappa_2=0$ (then $\beta=0$ and copy~$2$ is free), on which
$\eta_2=\Re\bigl((v_1^2+v_2^2)\,c\cdot c\bigr)$ and $v_1^2+v_2^2$ is a nonzero multiple of $\kappa_2(\kappa_1-\kappa_2)$, respectively of $1$, hence $\neq0$; so the kernel is not $\eta_2$-null and
\cref{bf:criterion}(i) fails.
\end{proof}

Item (i) restates \cref{thm:subfs} (fixed there) for each copy: duplication preserves the neutral signature, and the cancellation
survives unequal rates and masses. Since $E=C_v$, (ii) and (iii) specialise \cref{d5:thm:bfield2} to $m_1=m_2$. The lines in (ii) are inherited two-level mixing and restate
\cref{d5:thm:bfield2}(ii); the boundedness criterion restates \cref{d5:thm:bfield2}(v); the seesaw value $0.0495097568$ at
$(\kappa_1,\kappa_2,\beta)=(0,5,\half)$ against $\beta^2/\kappa_2=0.05$ restates \cref{d5:thm:bfield2}(vi). What is new here is the Krein signature
$(3,3)$ of each line and the vanishing zero-point energy at $n=2$, fixed by \cref{bf:krein}. (iii) restates the Jordan block in the proof
of \cref{d5:thm:bfield2}(v) (fixed there); it is the inter-manifold analogue of the exceptional ratio of \cref{bf:screen}. Checked numerically for the signatures (d6-14a,b) and with rational
data for semisimplicity: not semisimple at $(1,4,2)$, semisimple at $(1,4,1)$, $(2,2,2)$ and $(1,3,\tfrac25)$ (d6-14c). Read with the added
identification $mc^2=\hbar\kappa$ applied to the lines: an inter-manifold $b$ can give an unbound manifold the line $\beta^2/\kappa_2$ borrowed from the
binding of the other, or bring a line to zero, at unequal rates exactly at the price of unbounded motion; it never changes a rate, the spectrum of the
self-adjoint part of the shift (\cref{d7:thm:schr}(iii)), and never produces a line without an existing binding, since at
$\kappa_1=\kappa_2=0$ the lines $\pm\beta$ are removed by $X\mapsto e^{-\beta E\tau}X$.

% --- source: rest_energy/d6_pair#8 (Assessment)
\subsection{Assessment}
\label{d6:sec:assess}

\begin{table}[t]
\centering
\small
\begin{tabular}{@{}>{\raggedright\arraybackslash}p{2.7cm}>{\raggedright\arraybackslash}p{7.2cm}>{\raggedright\arraybackslash}p{3.2cm}>{\raggedright\arraybackslash}p{1.4cm}@{}}
\toprule
claim & statement & class & check\\
\midrule
Thm.~\ref{d6:thm:balance}(i) & rate constant and nonzero iff $R_{\rm eff}=0$ and $h(0)\neq0$; frames then on circles & fixed (restated, \cref{prop:several}) & d6-2a--d\\
Thm.~\ref{d6:thm:balance}(ii) & $n=1$: a constant nonzero rate with nonzero Lorentz-class inertia requires $E_0(\pi)=0$ & fixed (restated, \cref{cor:constrate}) & d6-1a\\
Thm.~\ref{d6:thm:balance}(iii) & $n\ge2$: codimension-two balance set, massive data dense & fixed & d6-1b\\
Rem.~\ref{d6:rem:lorentzreq} & without the Lorentz requirement $n=1$ suffices & fixed (counterexample) & d6-1c\\
Cor.~\ref{d6:cor:nopair} & no closure Hamiltonian creates or moves a pair & obstructed & d6-3\\
Prop.~\ref{d6:prop:pairs} & $a\cdot b=0$; $V_\pm$ null, $E=\mp G_2$; magnetic--electric pair, $M_{\rm pair}=2E_0$, $q_{\rm tot}=0$ & fixed (part restated); readings interpretive & d6-4a--e\\
Prop.~\ref{d6:prop:covariant}(i,ii) & commutant $M_2(\C)$, $\mathfrak g(0,2,0)$, invariants of $R_\varphi$ and $U(2)$ & restated / inherited & d6-5a,b\\
Prop.~\ref{d6:prop:covariant}(iii) & no covariant map makes a pair from one null field & obstructed & d6-5c\\
Thm.~\ref{d6:thm:pair} & $M_{\rm pair}=\half(\sigma_1+\sigma_2)$, the chain, $\max_\varphi$ & (i), (iii), (v) inherited; (ii), (iv) derived here, priority unestablished & d6-6a--d\\
Ex.~\ref{d6:ex:bw} & two null fields $\to$ electric and magnetic at rest & inherited & d6-7\\
Ex.~\ref{d6:ex:R1} & rotation of one direction: balanced pair from a null field, covariant under none of the five Lorentz actions at $n=2$; third body & fixed (numbers, non-covariance); agent constructed; reading interpretive & d6-8; x6-8a,b\\
Prop.~\ref{d6:prop:gauge} & $R_\varphi$ admissible; per-copy $E_0$ not invariant; $O(6,\C)$ maps pairs to null fields & fixed (part restated) & d6-9a--c, 10\\
interior & distinguished copies: per-copy rest energy changed by two-level mixing, sum $\le M_{\rm pair}$ & inherited & d6-11\\
Rem.~\ref{d6:rem:copyrotor} & a dynamical copy rotor changes $E_0^{(k)}$ and their sum at fixed pair invariants & constructed & d6-12\\
Prop.~\ref{d6:prop:dirac} & Klein signs, $C$-parity, masses in quadrature & inherited & d6-13a--e\\
Prop.~\ref{d6:prop:zpe} & formal cancellation of the zero-point energy of the nonzero lines of the indefinite Hamiltonian; exceptional ratio $\beta^2=\kappa_1\kappa_2$ & fixed ((i) restated; lines, boundedness and (iii) restate \cref{d5:thm:bfield2}) & d6-14a--c\\
answers & duplication as copy rotation, pair creation and balance & interpretive & ---\\
\bottomrule
\end{tabular}
\caption{The claims of \cref{d6:sec:main}, their classes and the checks of \texttt{d6\_pair\_balance.py} (d6-0 verifies the conventions; x6-8a,b are printed by \texttt{x6\_pair\_balance\_referee.py}).}
\label{d6:tab:claims}
\end{table}

\emph{On the $b$-particle with two manifolds.} The remark is partly borne out. With two manifolds there is an admissible $b$-field between
them, the copy rotation, the compact part of $\mathfrak g(0,2,0)$; it commutes with the binding root $G_2$, and with the binding $K\otimes G$ only at equal rates (\cref{d5:thm:bfield2}(iii)), keeps every orbit bounded except when
$\beta^2=\kappa_1\kappa_2$ with $\kappa_1\neq\kappa_2$, or when $\kappa_1=\kappa_2=0$ and nothing binds (\cref{d5:thm:bfield2}(v)), and changes the rest energy of each manifold, their sum bounded by $M_{\rm pair}$ but not conserved in general, which on one manifold only the anti-diagonal part of $b$
can do, and only by tearing or by screening (\cref{en:transport,bf:screen}); with unequal rates it mixes them, lending a line to an unbound
manifold or bringing one to zero, while the rates themselves never change (\cref{d7:thm:schr}(iii)). It does not generate an invariant mass: the copy rotation preserves $M_{\rm pair}$, the summed stress tensor and $q_{\rm tot}$, and the
interior flow $e^{ih\tau}$, unitary on the copy index, preserves $M_{\rm pair}$, the summed stress tensor and $E_6^{\rm tot}$ at any rates but
$|q_{\rm tot}|$ only at equal rates (for the pair of \cref{d6:ex:bw} at $\kappa=(1,3)$, $\beta=0.4$ it leaves $0$); with equal rates the copy rotation is removable along a worldline (\cref{d5:thm:bfield2}(iii)); the rest energy of each copy is not invariant under the copy rotation, hence gauge dependent if that rotation is gauged; and the one thing two constituents do in
the closure that one cannot, carrying rest energy at a constant rate, needs no $b$-field at all: it is the balance surface.

\emph{On how a unitary energy results in a duplication.} The mathematical process is the graded tensor product (inherited), under which the
undivided energy becomes $E_6^{\rm tot}=\half|Z|^2$; the Lorentz-covariant, energy--momentum-preserving duplications form $U(2)\otimes\id_3$, of
which the metric ones form $O(2)\otimes\id_3$ (the real orthogonal $A$); its identity component, the copy rotation $SO(2)$, is the part generated by an admissible $b$-field, and the other component contains the swap of the copies and $\operatorname{diag}(1,-1)$, which are metric but not copy rotations (\cref{d6:prop:covariant}). The physical process it
describes shares the energy of one copy between two as literal replicas, preserving the type. From one null field, a balanced magnetic--electric pair requires either
a second, non-collinear null field, with the inherited kinematics of $\gamma\gamma\to e^+e^-$ ($\sum_kE_0^{(k)}\le M_{\rm pair}$, saturated), or a
non-covariant agent that takes the momentum, with the kinematics of pair production on a nucleus ($\sum_kE_0^{(k)}\le E_6^{\rm tot}$). From one massive field $z$ the copy rotation never balances $(z,0)$, since it preserves $q_{\rm tot}=q(z)\neq0$; the non-metric equal split with a relative quarter turn of duality, $z\mapsto(z,iz)/\sqrt2$ (\cref{d5:thm:dup}(ii) at $\theta=\pi/2$), is covariant and energy-preserving, lands on $V_+$ and gives the balanced pair of \cref{d6:prop:pairs}(iii), with $E_0^{(k)}=\half E_0(z)$ each and $M_{\rm pair}=E_0(z)$ unchanged (\cref{d5:prop:balance}). In motion
the closure conserves the content of each copy and carries the pair, on the balance surface, at a constant rate. The decisive caveat is
\cref{d6:prop:gauge}: every statement that a mass has been created in one copy needs a structure that distinguishes the copies, and the
symmetric closure and the gauged copy rotation provide none. Whether such a structure is physical is open.

% --- source: emergence/e3_ladder#0 ((head))
% Track e3 -- Up the ladder: the rung Cl(3n,3n).
% New references: sections/bib_e3.tex (Howe89). Also cited: Chevalley54, LawsonMichelsohn89, Pauli40, WeinbergQFT (bib_e1.tex),
% EhrenfestOppenheimer31 (bib_e2.tex).

% --- source: emergence/e3_ladder#1 (Up the ladder)
\section{Up the ladder}
\label{e3:sec:main}

\Cref{e1:sec:main,e2:sec:main} treat the second rung. This section treats the rung $\Cl(3n,3n)$ for every $n$ and the step from $n$ to
$n+1$. The question is the same as at $n=2$: in what precise sense, and under which added structures, does rung $n$ carry $n$ particles
rather than one particle with a larger internal space, and what changes from one rung to the next. The answers, with their classes, are
these. (1) $\Cl(3n,3n)=M_{8^n}(\R)$ at every rung. The step embeds rung $n$ in rung $n+1$ with multiplicity eight and adds three Fock
modes, one frame and one copy direction; the commuting copy of the new constituent is its generators dressed by the volume element of
rung $n$ (fixed; \cref{e3:sec:rung}). (2) For $n$ equal masses the graded rest Hamiltonian is one clock of mass $\sqrt n\,m$, and the
ungraded one is $n$ clocks of mass $m$, with spectrum $m(n-2j)$ of multiplicity $\binom nj4^n$ (fixed, partly restated from
\cref{e1:prop:rest}; \cref{e3:sec:mass}). (3) A rotation through $2\pi$ of the diagonal interior Lorentz action is $(-1)^n$ on $S_n$, and one
of the diagonal field-strength action is $+\id$. The interior composite has half-integer spin exactly at odd $n$, a fermion under the
spin--statistics reading, and the field-strength composite has integer spin at every rung (restated from \cref{prop:orthospin}, with the second half fixed; \cref{e3:sec:spin}). (4) On vectors
the copy rotations of all six directions span the commutant of the diagonal $\so{3,3}$, and the copy group $O(n)$ has the diagonal $\so{3,3}$
as its commutant. On $S_n$ the Casimirs satisfy $C_{\rm cop}+C_{\rm diag}=\tfrac34n(n+4)$. This is Howe's skew duality $(O(n),\so6)$ on the spinor module of $\so{6n}$, and at
$n=2,3$ the joint decomposition has the predicted multiplicity-free form (fixed; the duality inherited; \cref{e3:sec:dual}). (5) The Koszul
swaps realise the permutation matrices of that $O(n)$, twisted by the sign character. Their symmetric and antisymmetric sectors are mirror images, of equal dimension and
never empty. The plain permutations of the common complex space give $\Lambda^n_\C$ of the Dirac module, which vanishes for $n\ge5$; this
concerns one momentum label of the finite construction, with antisymmetry imposed, and is not a bound on fermions in physical space (fixed;
the exterior powers inherited; \cref{e3:sec:sym}). (6) The algebra singles out no rung, and its invariant forms single out only the parity of $n$, through a period of four. Odd rungs are
singled out by the interior Lorentz action, even rungs by the binding action, and $n=3$ by the projector $P_*$ (restated and fixed; \cref{e3:sec:table}). (7) Going up the ladder
enlarges one particle; it becomes the addition of a particle only through the constructed structures of \cref{e1:sec:main}, now needed at every
rung (interpretive, drawn from the fixed statements; \cref{e3:sec:answer}). Identifiers e3-$n$ refer to the \texttt{PASS} lines of
\texttt{e3\_ladder.py} (\cref{e3:tab:claims}).

\paragraph{Notation.} As in \cref{e1:sec:main}: $S=\R^8$, $S_n=S^{\otimes n}$, $E_{(k,a)}=\omega^{\otimes(k-1)}\otimes e_a\otimes\id^{\otimes(n-k)}$, and
$V_n=\R^n\otimes V=\R^{3n,3n}$ with $\eta_n=\id_n\otimes\eta$ and $G_n=\id_n\otimes G$. $\Gamma^{(k)}=E_{(k,1)}\cdots E_{(k,6)}=\id\otimes\cdots\otimes\omega\otimes\cdots\otimes\id$ is the
volume element of constituent $k$, and $\Gamma_n=\Gamma^{(1)}\cdots\Gamma^{(n)}$. The ungraded generators are
$u_{(k,a)}=\id^{\otimes(k-1)}\otimes e_a\otimes\id^{\otimes(n-k)}$, $J^{(k)}=E_{(k,2)}E_{(k,3)}$, and $a^\pm_{k,i}=\half(E_{(k,i)}\pm E_{(k,i+3)})$, $N=\sum_{k,i}a^+_{k,i}a^-_{k,i}$
(\cref{d5:prop:fock}); $|0\rangle$ is the first basis vector of $S_n$, the Fock vacuum. The diagonal generators are
$\Delta_{ab}=\half\sum_kE_{(k,a)}E_{(k,b)}$. All fifteen span the \emph{diagonal} $\so{3,3}$; those with $a,b\in\{1,4,5,6\}$ span the diagonal $\lint$; and the
paired rotations $\Delta_{ij}-\Delta_{i+3,j+3}$ and paired boosts $\Delta_{i,j+3}-\Delta_{j,i+3}$ span the diagonal $\lG$ (\cref{thm:GLorentz}). The
copy rotations are $C^{kl}=\half\sum_a\eta_{aa}E_{(k,a)}E_{(l,a)}$ for $k<l$, with $C^{12}=C_b$ of \cref{d5:prop:cmrel}(ii), and $M_{kl}$ are their
interior parts (\cref{sec:diagonal}). The Casimirs are $C_{\rm diag}=-\sum_{a<b}\eta^{aa}\eta^{bb}\Delta_{ab}^2$ over all six directions,
$C_\Delta$ the same over the interior four (\cref{sec:dualpair}), $C_G=-\sum_iL_i^2+\sum_iK_i^2$ for the paired rotations $L_i$ and boosts
$K_i$ of $\lG$, $C_{\rm cop}=-\sum_{k<l}(C^{kl})^2$ and $C^{\rm int}_{\rm cop}=-\sum_{k<l}M_{kl}^2$. A Lorentz Casimir takes the value
$2[j_L(j_L+1)+j_R(j_R+1)]$ on $(j_L,j_R)$. We write $\mathfrak S_n$ for the symmetric group, to keep $S_n$ for the module.

% --- source: emergence/e3_ladder#2 (The rung and the step)
\subsection{The rung and the step}
\label{e3:sec:rung}

\begin{proposition}[The rung and the step]
\label{e3:prop:rung}
\begin{enumerate}[label=(\roman*),leftmargin=2.2em]
\item Every blade of $\Cl(3n,3n)$ other than $\pm\id$ is traceless on $S_n$. Hence the $64^n$ blades are linearly independent and
$\Cl(3n,3n)\cong M_{8^n}(\R)$, with $S_n$ its unique irreducible module. The volume element is $E_{(1,1)}\cdots E_{(n,6)}=\Gamma_n=\omega^{\otimes n}$, and
$(-1)^N=(-1)^n\Gamma_n$.
\item $E^{(n+1)}_{(k,a)}=E^{(n)}_{(k,a)}\otimes\id$ for $k\le n$ and $E_{(n+1,a)}=\Gamma_n\otimes e_a$. So $\Cl(3n,3n)\otimes\id\subset\Cl(3n+3,3n+3)$, and
$S_{n+1}$ restricted to $\Cl(3n,3n)$ is eight copies of $S_n$. The commutant of $\Cl(3n,3n)\otimes\id$ is $\id\otimes M_8(\R)$, generated by the
ungraded generators $\id\otimes e_a=(\Gamma_n\otimes\id)E_{(n+1,a)}$ of the new constituent.
\item $N_{n+1}=N_n\otimes\id+\id\otimes N_1$, $\spec N_n=\{0,\dots,3n\}$ with multiplicities $\binom{3n}j$, and $[N,E_{(k,i)}]=E_{(k,i+3)}$,
$[N,E_{(k,i+3)}]=E_{(k,i)}$. Each step adds three modes.
\item Every diagonal bivector satisfies $\sum_{k\le n+1}E_{(k,a)}E_{(k,b)}=\bigl(\sum_{k\le n}E_{(k,a)}E_{(k,b)}\bigr)\otimes\id+\id\otimes e_ae_b$. Hence
$S_{n+1}\cong S_n\otimes S$ and $S_n\cong S^{\otimes n}$ as modules of the diagonal $\so{3,3}$, and so of the diagonal $\lint$ and $\lG$.
\end{enumerate}
\end{proposition}

\begin{proof}
(i) Moving the generators of each constituent to their own factor gives $e_I=\pm\bigotimes_je_{I_j}\omega^{p_j}$ with $p_j=|I_{j+1}|+\cdots+|I_n|$ mod $2$, and
the trace is multiplicative. On $S$, $\tr(e_J\omega^p)\ne0$ only for $(J,p)=(\emptyset,0)$ and for $J$ the full index set with $p=1$. The last
factor has $p_n=0$, so a nonzero trace needs $I_n=\emptyset$; then $p_{n-1}=0$, and so on, down to $I=\emptyset$. Trace orthogonality gives
independence, and $64^n=(8^n)^2$. The volume element is $\prod_k\Gamma^{(k)}$ because $\prod_aE_{(k,a)}=(\omega^{\otimes(k-1)})^6\otimes\omega\otimes\id$. On one
constituent $(-1)^N=-\omega$ (\cref{d5:prop:fock}), and $N=\sum_kN^{(k)}$. (ii) The formulas are equation~\cref{eq:tower} at rungs $n$
and $n+1$. The commutant of $M_{8^n}(\R)\otimes\id$ in $M_{8^{n+1}}(\R)$ is $\id\otimes M_8(\R)$, and $\Gamma_n^2=\id$. (iii) $N=\sum_{k,i}\half(\id-E_{(k,i)}E_{(k,i+3)})$,
the terms are commuting diagonal involutions with traceless products, and the commutators are those of \cref{d5:prop:fock}. (iv)
$E_{(n+1,a)}E_{(n+1,b)}=\Gamma_n^2\otimes e_ae_b=\id\otimes e_ae_b$.
\end{proof}

Part (i) is fixed in the realisation (e3-1a, exhaustive over all $64^n$ blades at $n\le4$). That $\Cl(p,q)\cong M_{2^{(p+q)/2}}(\R)$ for $p-q\equiv0\pmod8$
is inherited \cite{LawsonMichelsohn89}: the real structure of the algebra is the same at every rung, so it singles out none. Parts (ii) and
(iii) restate \cref{d5:prop:graded} and~\cref{d5:prop:fock} and \cref{e1:prop:copies}(ii),(v) at the last factor (e3-1b, with the commutant
exhaustive over blades at $n=1,2$; e3-1c, spectrum also at $n=5$). Part (iv) is fixed (e3-6c, $n\le3$). Under (iv) every statement about the
diagonal action is a statement about tensor powers of the one-constituent module, and the step is tensoring with $S$.

% --- source: emergence/e3_ladder#3 (Mass up the ladder: one clock or nn clocks)
\subsection{Mass up the ladder: one clock or \texorpdfstring{$n$}{n} clocks}
\label{e3:sec:mass}

Put $H_g=m\sum_kE_{(k,1)}$, built from the vectors of the rung, and $H_u=m\sum_ku_{(k,1)}$, built from the commuting copies.

\begin{proposition}[Mass up the ladder]
\label{e3:prop:mass}
\begin{enumerate}[label=(\roman*),leftmargin=2.2em]
\item The $E_{(k,1)}$ are pairwise anticommuting traceless involutions. So $H_g^2=nm^2\id$ and $\spec H_g=\{\pm\sqrt n\,m\}$, each of multiplicity
$8^n/2$. Every $J^{(k)}$ commutes with $H_g$, and $Z_g=J^{(1)}H_g$ satisfies $Z_g^2=-nm^2\id$: the rest flow $e^{-\lambda Z_g}$ is one clock of
angular frequency $\sqrt n\,m$.
\item The ungraded generators are $u_{(k,a)}=\Gamma^{(1)}\cdots\Gamma^{(k-1)}E_{(k,a)}$, with no further sign. The $u_{(k,1)}$ commute, square to $\id$, and
every product over a nonempty subset is traceless. So $\spec H_u=\{m(n-2j):j=0,\dots,n\}$, the value $m(n-2j)$ with multiplicity $\binom nj4^n$.
\item $T^{(k)}_u=J^{(k)}u_{(k,1)}=\id\otimes\cdots\otimes T\otimes\cdots\otimes\id$ are $n$ commuting complex structures: $n$ clocks of angular frequency $m$.
\item $H_g^{(n+1)}=H_g^{(n)}\otimes\id+\Gamma_n\otimes me_1$ with anticommuting summands, and $H_u^{(n+1)}=H_u^{(n)}\otimes\id+\id\otimes me_1$ with commuting
summands. The graded mass grows in quadrature, $\sqrt n\,m\to\sqrt{n+1}\,m$, and the ungraded energies add.
\end{enumerate}
\end{proposition}

\begin{proof}
(i) Graded generators of distinct constituents anticommute, and each $E_{(k,1)}$ is traceless because $\tr e_1=0$. $J^{(k)}$ is a bivector sharing no
index with any $E_{(l,1)}$, so it commutes with each (equation~\cref{eq:bladecomm}), and $(J^{(k)})^2=-\id$. (ii) $\Gamma^{(1)}\cdots\Gamma^{(k-1)}=\omega^{\otimes(k-1)}\otimes\id$ and $\omega^2=\id$, so the
product is $\id^{\otimes(k-1)}\otimes e_a\otimes\id$. Products of the $u_{(k,1)}$ are tensor products of factors $e_1$ and $\id$, traceless unless all
factors are $\id$. Commuting involutions whose nonempty products are traceless have all $2^n$ joint signs with equal multiplicity $8^n/2^n=4^n$, and $j$ minus signs give
$m(n-2j)$. (iii) $J\,e_1=e_2e_3e_1=T$ on the $k$th factor. (iv) \Cref{e3:prop:rung}(ii); $\Gamma_n$ anticommutes with every generator of $\Cl(3n,3n)$.
\end{proof}

At general masses (i) and (ii) are \cref{e1:prop:rest}(vi), restated; the multiplicities, the clocks (iii) and the step (iv) are fixed
(e3-2a and e3-2b, exact over signed permutations at $n\le5$, eigenvalue counts in floating point at $n\le3$; e3-6c). At $n=2$, (i) is the
Dirac duplication of \cref{d6:prop:dirac}(iv). So rung $n$ read through its vectors carries one Dirac mass, $\sqrt n\,m$, along the unit
vector $n^{-1/2}\sum_kE_{(k,1)}$ of $\R^{3n,3n}$, with one clock; read through its commuting copies it carries $n$ masses $m$ with $n$ clocks and
threshold $nm$. The quadrature is the statement that the copies of the graded tower are directions of one momentum space (\cref{e1:sec:rest}).

% --- source: emergence/e3_ladder#4 (Spin parity)
\subsection{Spin parity}
\label{e3:sec:spin}

Let $R_n=\half\sum_kE_{(k,4)}E_{(k,5)}$, the rotation about the third axis of the diagonal $\lint$, and
$L_n=\half\sum_k\bigl(E_{(k,1)}E_{(k,2)}-E_{(k,4)}E_{(k,5)}\bigr)$, the paired rotation of the $(1,2)$ and $(4,5)$ planes in the diagonal $\lG$. Both
generate rotations of period $2\pi$ on vectors.

\begin{proposition}[Spin parity up the ladder]
\label{e3:prop:spin}
\begin{enumerate}[label=(\roman*),leftmargin=2.2em]
\item $R_n$ is a sum of $n$ commuting terms $\half b_k$ with $b_k^2=-\id$. Hence $e^{2\pi R_n}=(-1)^n\id$, and the eigenvalues of $2R_n$ are $i$ times
integers congruent to $n$ modulo $2$. Under the diagonal $\lint$ the whole of $S_n$ has half-integer spin for odd $n$ and integer spin for even
$n$. The Casimir spectra are
\[
\begin{array}{@{}ll@{}}
n=1: & C_\Delta=\tfrac32\,(8),\\
n=2: & 0\,(8),\ 3\,(32),\ 4\,(24),\\
n=3: & \tfrac32\,(160),\ \tfrac{11}2\,(288),\ \tfrac{15}2\,(64),\\
n=4: & 0\,(160),\ 3\,(1024),\ 4\,(864),\ 8\,(864),\ 9\,(1024),\ 12\,(160),
\end{array}
\]
where the values $0,\tfrac32,3,4,\tfrac{11}2,8,\tfrac{15}2,9,12$ are those of $(0,0)$, $(\half,0)\oplus(0,\half)$, $(\half,\half)$, $(1,0)\oplus(0,1)$, $(1,\half)\oplus(\half,1)$, $(1,1)$, $(\tfrac32,0)\oplus(0,\tfrac32)$,
$(\tfrac32,\half)\oplus(\half,\tfrac32)$ and $(2,0)\oplus(0,2)$.
\item $L_n$ commutes with $G_n$ on vectors and is a sum of $2n$ commuting terms $\pm\half b$ with $b^2=-\id$. Hence $e^{2\pi L_n}=+\id$ at every rung. The
diagonal $\lG$ has Casimir $C_G=3$ on all of $S$; $0,4,8$ on $4,24,36$ dimensions at $n=2$; and $3,9,15$ on $128,256,128$ at $n=3$: the content of
$(\half,\half)^{\otimes n}$, integer spin at every rung.
\item The central element $-\id$ of $\mathrm{Spin}(3,3)$, which is $e^{\pi e_4e_5}=e^{\pi e_1e_2}$ on $S$, acts on $S_n$ through the diagonal
$\mathrm{Spin}(3,3)$ as $(-1)^n$. The rotation through $2\pi$ of $\lint$ is this element, and that of $\lG$ is not.
\item For the actions of \cref{thm:algebra}, with $\lint$ on $k_I$ constituents, $\lG$ on $k_B$ and trivially on the rest, the rotation
through $2\pi$ is $(-1)^{k_I}$.
\end{enumerate}
\end{proposition}

\begin{proof}
(i) $(E_{(k,4)}E_{(k,5)})^2=-E_{(k,4)}^2E_{(k,5)}^2=-\id$, and bivectors of distinct constituents share no index and commute. So
$e^{2\pi R_n}=\prod_k(\cos\pi+b_k\sin\pi)$. Each listed value of $2[j_L(j_L+1)+j_R(j_R+1)]$ is attained by exactly one unordered pair $(j_L,j_R)$, which identifies the
modules. (ii) $e_ie_j-e_{i+3}e_{j+3}$ commutes with $G$ (\cref{d5:prop:so6C}(i)). The $2n$ terms share no index, and
$e^{2\pi L_n}=\prod_k(-\id)(-\id)$. By \cref{e3:prop:rung}(iv), $S_n\otimes\C\cong(S\otimes\C)^{\otimes n}=2^n(\half,\half)^{\otimes n}$ under the diagonal $\lG$,
since $S\otimes\C=2(\half,\half)$ (\cref{prop:halfspinors}); the Clebsch--Gordan series gives the listed values. (iii) The diagonal
$\mathrm{Spin}(3,3)$ acts on $S^{\otimes n}$ by $g\mapsto g^{\otimes n}$ (\cref{e3:prop:rung}(iv)). $e^{2\pi\cdot\frac12e_4e_5}$ is $-\id$, while the double
rotation $e^{2\pi\cdot\frac12(e_1e_2-e_4e_5)}=e^{\pi e_1e_2}e^{-\pi e_4e_5}$ is $+\id$. (iv) Products of the factors of (i) and (ii) over the respective
constituents.
\end{proof}

Part (i) restates \cref{prop:orthospin} and \cref{prop:diagtower}(i) (fixed there), with the element and the spectra re-verified (e3-3a,
exact at $n\le5$; e3-3c, floating-point spectra at $n\le4$). Part (ii) extends \cref{prop:halfspinors} from $n=1$ to every rung and is
fixed (e3-3b, e3-3c). Part (iii) is fixed (e3-3b), and (iv) restates \cref{prop:orthospin}. The parity of (i) is the
Ehrenfest--Oppenheimer rule, a composite of $n$ spin-$\half$ particles is a fermion iff $n$ is odd (inherited \cite{EhrenfestOppenheimer31}).
The rule holds equally for one particle whose internal space is the $n$-fold tensor product of spinors. So the spin parity does not decide
between one particle and $n$. Under the field-strength reading it does not arise: every rung carries integer spin, $\mathrm{Spin}(3,3)\cong SL(4,\R)$
restricted to $\lG$ integrates to $SO^+(1,3)$ on every $S_n$, and the diagonal $\lG$ composite is a boson at every rung.

\begin{remark}[Comparison with \cref{thm:copyparity}]
\label{e3:rem:copyparity}
\cref{thm:copyparity} says that a copy-invariant four-vector in $\Lambda^g(\R^n\otimes\R^{1,3})$ needs $n$ odd and $g$ odd, $n\le g\le3n$. The
statement concerns the algebra and (i) concerns the module, and they are logically independent. Both are products over the $n$ constituents:
in (i) each constituent contributes a factor $-1$ to the $2\pi$ rotation, and in the theorem the copy invariant of odd grade needs the
Levi-Civita symbol of $SO(n)$ with its $n$ indices. They meet on the Dirac block. At odd $n$ the module contains spin $\half$, and at $n=3$ the
copy-singlet four-vector $Q_\mu$ of grade three builds a gauge-invariant $\Cl(1,3)$ on that block (\cref{thm:blockgamma}). The check e3-6b
re-verifies that $Q_\mu$ commutes with the interior copy rotations $M_{kl}$ and satisfies $[\Delta_{ab},Q_c]=\eta_{bc}Q_a-\eta_{ac}Q_b$, and that the $n=2$
analogue $\sum_\nu\eta_{\nu\nu}(\gamma^{(1)}_\mu\gamma^{(2)}_\nu+\gamma^{(1)}_\nu\gamma^{(2)}_\mu)$ commutes with $M_{12}$ but is not a vector. At even $n$ there is
neither a fermion to carry gamma matrices nor a copy-invariant four-vector to build them. Under the diagonal $\lG$ neither statement has
an analogue: the spin is integral at every rung, and grade one is $n[(1,0)\oplus(0,1)]$, $n$ field strengths and no four-vector (\cref{sec:spintower}). The juxtaposition is interpretive; the facts are restated and fixed.
\end{remark}

% --- source: emergence/e3_ladder#5 (A dual pair up the ladder)
\subsection{A dual pair up the ladder}
\label{e3:sec:dual}

\begin{theorem}[The copy algebra and the diagonal $\so{3,3}$]
\label{e3:thm:dual}
\begin{enumerate}[label=(\roman*),leftmargin=2.2em]
\item On vectors, the commutant of the diagonal $\so{3,3}$ in $\so{3n,3n}$ is $\so n\otimes\id_6$, of dimension $\binom n2$, spanned by the vector
actions of the $C^{kl}$. The commutant of $\so n\otimes\id_6$ together with the copy permutations is $\id_n\otimes\so{3,3}$ for $n\ge2$; that of
$\so n\otimes\id_6$ alone has dimension $36$ at $n=2$ and $15$ for $n\ge3$. The weights $\eta_{aa}$ are forced: $\sum_aE_{(k,a)}E_{(l,a)}$ does not commute
with the diagonal $\so{3,3}$.
\item On $S_n$, $C^{kl}$ is the spin lift of $(e_ke_l^{\!\top}-e_le_k^{\!\top})\otimes\id_6$: $[C^{kl},E_{(l,a)}]=E_{(k,a)}$ and $[C^{kl},E_{(k,a)}]=-E_{(l,a)}$. The
$C^{kl}$ commute with the fifteen $\Delta_{ab}$ and close on $\so n$, with $[C^{kl},C^{lm}]=C^{km}$. Each is a sum of six commuting terms of square
$-\tfrac14$, so $e^{2\pi C^{kl}}=\id$. In Fock form
\[
C^{kl}=\sum_{i=1}^3\bigl(a^+_{k,i}a^-_{l,i}-a^+_{l,i}a^-_{k,i}\bigr),\qquad C^{kl}a^+_{l,i}|0\rangle=a^+_{k,i}|0\rangle,
\]
so the copy algebra conserves $N$ and moves quanta between constituents.
\item As operators on $S_n$,
\[
C_{\rm cop}+C_{\rm diag}=\tfrac34\,n(n+4)\,\id,
\]
that is $\tfrac{15}4,9,\tfrac{63}4,24,\tfrac{135}4,\tfrac{81}2$ at $n=1,\dots,6$. The interior identity of \cref{thm:casid},
$C^{\rm int}_{\rm cop}+C_\Delta=\half n(n+2)\,\id$, holds at every $n$. Both constants are $nd(d+n-2)/8$, at $d=6$ and $d=4$.
\item At $n=2$ the commutant of the diagonal $\so{3,3}$ in $M_{64}(\R)$ has dimension $14$ and is spanned by the words in $C^{12}$ and the Koszul
swap $\tau_K$, with $C^{12}\tau_K=-\tau_KC^{12}$: it is the image of the group algebra of $O(2)$. The words in $C^{12}$ alone span $7$ dimensions.
\end{enumerate}
\end{theorem}

\begin{proof}
(i) $V=\R^{3,3}$ is absolutely irreducible under $\so{3,3}$, so by Schur's lemma the commutant of $\id_n\otimes\so{3,3}$ in $\End(\R^n\otimes V)$ is
$\End(\R^n)\otimes\id_6$. An element $A\otimes\id_6$ is $\eta_n$-antisymmetric iff $A^{\!\top}=-A$. Conversely, $\R^n$ is absolutely irreducible under
the group generated by $SO(n)$ and the permutation matrices for $n\ge2$, and under $\so n$ alone for $n\ge3$; at $n=2$ the commutant of $\so2$
is $\operatorname{span}\{\id,\varepsilon\}$, which adds $\varepsilon\otimes\mathrm{Sym}_\eta(V)$, of dimension $21$. The unweighted sum lifts
$(e_ke_l^{\!\top}-e_le_k^{\!\top})\otimes\eta$, and $\eta$ does not commute with the boosts. (ii) The spin lift $\sigma(X)=\tfrac14\sum_{A,B}X_{AB}\eta_{BB}E_AE_B$
satisfies $[\sigma(X),E_C]=\sum_AX_{AC}E_A$ and is a Lie algebra homomorphism; the vector matrices commute by (i). The six terms
$\half\eta_{aa}E_{(k,a)}E_{(l,a)}$ share no index, and $(E_{(k,a)}E_{(l,a)})^2=-\eta_{aa}^2\id$. For the Fock form substitute
$E_{(k,i)}=a^+_{k,i}+a^-_{k,i}$, $E_{(k,i+3)}=a^+_{k,i}-a^-_{k,i}$: $C^{kl}=\sum_i(a^+_{k,i}a^-_{l,i}+a^-_{k,i}a^+_{l,i})$, and $a^-_{k,i}a^+_{l,i}=-a^+_{l,i}a^-_{k,i}$ for
$k\ne l$ by the CAR of \cref{d5:prop:fock}. (iii) The proof of \cref{thm:casid} uses neither $d=4$ nor the signature of the
interior: with $B^{kl}_{ab}=E_{(k,a)}E_{(k,b)}E_{(l,a)}E_{(l,b)}$ it gives $C_{\rm diag}=\tfrac{nd(d-1)}8-\half\mathsf S$ and $C_{\rm cop}=\binom n2\tfrac d4+\half\mathsf S$
with the same cross term $\mathsf S=\sum_{k<l}\sum_{a<b}\eta^{aa}\eta^{bb}B^{kl}_{ab}$. At $d=6$ the constant is $\tfrac{15n}4+\tfrac{3n(n-1)}2=\tfrac34n(n+4)$.
(iv) Computed: the nullity of $\sum_X\ad_X^{\!\top}\ad_X$ over the fifteen generators, and the rank of the words.
\end{proof}

Part (i) is fixed (e3-4a, exact over $\Q$ at $n\le5$); at $n=3$ its interior analogue is \cref{thm:diagcent}. Part (ii) restates
\cref{d5:prop:cmrel}(ii) at $n=2$ and \cref{d7:prop:isospin}(i), and is fixed at $n\le4$ (e3-4b, e3-4f); the Fock form at $n=2$ is
\cref{e2:prop:hop}(i) with the opposite orientation. Part (iii) is inherited: the constant $\tfrac34n(n+4)$ is the case $d_{\rm N}=n$, $k_{\rm N}=3$ of Neerg{\aa}rd's
relation $C_{\mathfrak o(d_{\rm N})}+C_{\mathfrak o(2k_{\rm N})}=\tfrac14k_{\rm N}d_{\rm N}(d_{\rm N}+2k_{\rm N}-2)$ \cite[eq.~(73)]{Neergard19}, and the Clifford-algebra proof is an independent
derivation in the real realisation. The identity was checked exactly
as matrices on $S_n$ at $n\le4$ and blade by blade in the abstract Clifford algebra at $n\le6$; the interior identity is restated and checked in
the same two ways (e3-4c). Part (iv) is fixed (e3-4e; nullity in floating point with gap $>1$, rank exact). The $\eta$-weighting is where the
signature enters: a copy rotation pairs two generators of the same square, so it is compact (\cref{thm:diagcent}), and $C_{\rm cop}\ge0$.

\paragraph{Howe's skew duality.} By \cref{d5:prop:fock}, $S_n\otimes\C$ is the exterior algebra on the $3n$ creation operators
$a^+_{k,i}$, that is $\Lambda(\C^n\otimes\C^3)$ with $k$ the copy index and $i$ the mode index. The complexified diagonal $\so{3,3}$ is spanned by
the quadratics $\sum_kx_ky_k$ with $x,y\in\{a^\pm_{\cdot,i}\}$, the $O(n)$-invariant contractions over the copy index, which form, up to multiples
of $\id$ and by the substitution used in (ii),
$\Lambda^2\C^3\oplus\mathfrak{gl}(3)\oplus\Lambda^2\C^3\cong\so{6,\C}$. By (ii) the copy rotations are the $\so n\subset\mathfrak{gl}(n)$ acting on the copy index.
This is the setting of Howe's skew $(O(n),\so{2m})$ duality on $\Lambda(\C^n\otimes\C^m)$, here with $m=3$ \cite{Howe89}: the two actions generate mutual
commutants, and the module decomposes without multiplicity as $\bigoplus_\lambda\rho_\lambda\otimes\pi_{\lambda'}$, with $\rho_\lambda$ an irreducible
representation of $O(n)$ and $\pi_{\lambda'}$ one of $\so6$ determined by it. The theorem, and the linear relation between the two Casimirs, are
inherited; (iii) is its constant and (iv) its commutant at $n=2$ in the tower. The $O(n)$ contains more than $SO(n)$: at $n=2$ the reflection is
$\tau_K$, and \cref{e3:sec:sym} shows that the Koszul swaps realise the permutation matrices at every rung.

\begin{table}[t]
\centering
\footnotesize
\begin{tabular}{@{}cccclcc@{}}
\toprule
$n$ & copy label & $C_{\rm cop}$ & $C_{\rm diag}$ & $\so6$ modules $(\lambda_1,\lambda_2,\lambda_3)$ & dimensions & rank in $S_n$\\
\midrule
$2$ & $|j|=0$ & $0$ & $9$ & $(1,1,1)$, $(1,1,-1)$ & $10+10$ & $1\cdot20=20$\\
 & $|j|=1$ & $1$ & $8$ & $(1,1,0)$ & $15$ & $2\cdot15=30$\\
 & $|j|=2$ & $4$ & $5$ & $(1,0,0)$ & $6$ & $2\cdot6=12$\\
 & $|j|=3$ & $9$ & $0$ & $(0,0,0)$ & $1$ & $2\cdot1=2$\\
\midrule
$3$ & $j=0$ & $0$ & $\tfrac{63}4$ & $(\tfrac32,\tfrac32,\pm\tfrac32)$ & $20+20$ & $1\cdot40=40$\\
 & $j=1$ & $2$ & $\tfrac{55}4$ & $(\tfrac32,\tfrac32,\pm\tfrac12)$ & $36+36$ & $3\cdot72=216$\\
 & $j=2$ & $6$ & $\tfrac{39}4$ & $(\tfrac32,\tfrac12,\pm\tfrac12)$ & $20+20$ & $5\cdot40=200$\\
 & $j=3$ & $12$ & $\tfrac{15}4$ & $(\tfrac12,\tfrac12,\pm\tfrac12)$ & $4+4$ & $7\cdot8=56$\\
\bottomrule
\end{tabular}
\caption{Joint decomposition of $S_n$ under the copy algebra and the diagonal $\so{3,3}$ (\cref{e3:thm:dual}). At $n=2$, $j$ is the
$\so2$ weight; at $n=3$, $j$ is the copy isospin. $C_{\rm diag}=\langle\lambda,\lambda+2\rho\rangle$ with $\rho=(2,1,0)$, so that the spinor has
$\tfrac{15}4$ and the vector $5$. Ranks are real dimensions of $S_n$, equal to complex dimensions of $S_n\otimes\C$. Check e3-4d.}
\label{e3:tab:joint}
\end{table}

\begin{proposition}[Joint decomposition]
\label{e3:prop:joint}
The joint spectrum of $(C_{\rm cop},C_{\rm diag})$ on $S_2$ and $S_3$ is that of \cref{e3:tab:joint}. Every rank is the dimension of the copy module
times the sum of the dimensions of the listed $\so6$ modules, and every $\so6$ Casimir equals $\tfrac34n(n+4)-C_{\rm cop}$. At $n=4$,
$C_{\rm cop}=0,3,4,8,9,12,15,16,19,24$ on $70,560,504,810,1024,200,320,450,144,14$ dimensions, the values $2[j_+(j_++1)+j_-(j_-+1)]$ of the
$\so4=\su2\oplus\su2$ modules $(0,0)$, $(\half,\half)$, $(1,0)\oplus(0,1)$, $(1,1)$, $(\tfrac32,\half)\oplus(\half,\tfrac32)$, $(2,0)\oplus(0,2)$, $(\tfrac32,\tfrac32)$,
$(2,1)\oplus(1,2)$, $(\tfrac52,\half)\oplus(\half,\tfrac52)$, $(3,0)\oplus(0,3)$, each paired with $\so6$ modules of Casimir $24-C_{\rm cop}$ whose
dimensions reproduce the ranks.
\end{proposition}

\begin{proof}
At $n=2,3$ the minimal polynomial of $C_{\rm cop}$ is checked exactly and the ranks are traces of its Lagrange projectors, in integer
arithmetic. The $\so6$ dimensions are the Weyl dimension formula of $D_3$. At $n=4$, $C_{\rm cop}$ commutes with $N$ (by (ii)) and is
diagonalised on each of the thirteen number sectors.
\end{proof}

This is fixed (e3-4d), and the $n=3$ copy spectrum restates \cref{d7:prop:isospin}(ii). The joint spectra lie on the line of
\cref{e3:thm:dual}(iii). Copy type and $\so{3,3}$ type determine each other, as $j_{\rm cop}$ and Lorentz type do in \cref{cor:howe}. Two features recur up the ladder. The sector of largest copy Casimir carries the smallest $\so{3,3}$ module: the trivial module at
$n=2,4$ and the half-spinors $\mathbf4,\bar{\mathbf4}$ of one constituent at $n=3$. And the copy isospin is integral at $n=2,3$, while at $n=4$ the
modules of $\so4=\su2\oplus\su2$ include doublets of each factor, always with $j_++j_-$ integral, as single-valuedness requires (\cref{e3:thm:dual}(ii)).

% --- source: emergence/e3_ladder#6 (The symmetric group)
\subsection{The symmetric group}
\label{e3:sec:sym}

Let $\tau_k=\id^{\otimes(k-1)}\otimes\tau_K\otimes\id^{\otimes(n-k-1)}$, with $\tau_K=D\,P$ the Koszul swap of \cref{d5:prop:cmrel}(iii). For a
surplus direction $s\in\{2,3\}$ let $\Theta_s=E_{(1,s)}\allowbreak E_{(2,s)}\cdots E_{(n,s)}$, and let $C_3=\sum_{a<b,\,c<d,\,e<f}\allowbreak\varepsilon_{abcdef}L^{ab}L^{cd}L^{ef}$ be the cubic invariant of
the diagonal $\so{3,3}$, with $L^{ab}=\eta^{aa}\eta^{bb}\Delta_{ab}$.

\begin{proposition}[The Koszul symmetric group]
\label{e3:prop:koszul}
\begin{enumerate}[label=(\roman*),leftmargin=2.2em]
\item $\tau_k^2=\id$, $(\tau_k\tau_{k+1})^3=\id$ and $\tau_k\tau_l=\tau_l\tau_k$ for $|k-l|>1$, so the $\tau_k$ define an action of $\mathfrak S_n$ on $S_n$. Moreover
$\tau_kE_{(j,a)}\tau_k=E_{(s_kj,a)}$ and $\tau_k=\tfrac18\prod_a(E_{(k,a)}-E_{(k+1,a)})$: $\tau_k$ is the lift to $\mathrm{Pin}(3n,3n)$ of the transposition matrix of
copies $k,k+1$, a product of six reflections. It commutes with the diagonal $\so{3,3}$, with $N$, with $H_g$ and with $C_{\rm cop}$, and conjugates
$C^{kl}$ as the permutation does. On the vacuum $\tau_k|0\rangle=-|0\rangle$.
\item At $n=3$ the $\mathfrak S_3$-isotypic parts have the content
\begin{center}
\begin{tabular}{@{}lccc|cccc@{}}
\toprule
 & $C_\Delta=\tfrac32$ & $\tfrac{11}2$ & $\tfrac{15}2$ & $C_{\rm cop}=0$ & $2$ & $6$ & $12$\\
\midrule
trivial & $24$ & $48$ & $16$ & $20$ & $36$ & $20$ & $12$\\
sign & $24$ & $48$ & $16$ & $20$ & $36$ & $20$ & $12$\\
standard & $112$ & $192$ & $32$ & $0$ & $144$ & $160$ & $32$\\
\bottomrule
\end{tabular}
\end{center}
(real dimensions; the standard rows are twice the multiplicity of the two-dimensional representation). The Dirac block, $C_\Delta=\tfrac32$,
contains $24$ totally symmetric and $24$ totally antisymmetric dimensions.
\item $\Theta_s$ is invertible, anticommutes with every $\tau_k$, commutes with the diagonal $\lint$ and with every $C^{kl}$, and anticommutes with $C_3$. It
maps the Koszul-symmetric part isomorphically onto the Koszul-antisymmetric part, as modules of $\lint\oplus\so n$, and exchanges the two
chiralities of $\so6$. At $n=3$ the eigenvalues of $8C_3$ are $-630,-210,-126$ on the symmetric part of the copy sectors $j=0,1,2$ and
$+630,+210,+126$ on the antisymmetric part.
\item The Koszul-symmetric part and the Koszul-antisymmetric part of $S_n$ have the same dimension, $\sum_k\binom{n-k+3}3\binom4k$, which is
$8,32,88,192,360$ at $n=1,\dots,5$. Neither is ever zero.
\end{enumerate}
\end{proposition}

\begin{proof}
(i) The relations are checked. $(E_{(k,a)}-E_{(k+1,a)})^2=2\eta_{aa}$, and the reflection in this vector exchanges $E_{(k,a)}$ and $E_{(k+1,a)}$; the product
over the six directions acts on vectors as the transposition matrix $P$, because the sign $(-1)^6$ of the twisted adjoint action is $+1$. Hence
$\tau_k\sigma(X)\tau_k^{-1}=\sigma(PXP^{-1})$ for every $X\in\so{\eta_n}$ and $\tau_kv\tau_k^{-1}=Pv$ for every vector. $P$ commutes with $\id_n\otimes\so{3,3}$ and with
$K_n$, and fixes $\sum_kE_{(k,1)}$, so $\tau_k$ commutes with the $\Delta_{ab}$, with $N=\tfrac{3n}2+\sigma(K_n)$ (\cref{e2:lem:N}(i)) and with $H_g$; it permutes
the $C^{kl}$ up to sign and so fixes $C_{\rm cop}$. (ii) Traces of the isotypic projectors
$\frac{\dim\rho}{6}\sum_\sigma\chi_\rho(\sigma)\tau_\sigma$ times the spectral projectors of $C_\Delta$ and $C_{\rm cop}$, in exact integer arithmetic. (iii) $\Theta_s$ is a
blade. Against $E_{(k,a)}-E_{(k+1,a)}$ it gives the sign $(-1)^n$ for $a\ne s$ and $(-1)^{n-1}$ for $a=s$, so against $\tau_k$ the sign is $(-1)^{6n-1}=-1$.
The interior bivectors share no index with $\Theta_s$; the terms of $C^{kl}$ share none or two; a diagonal bivector $\Delta_{sb}$ shares one, so
$\Ad\Theta_s$ acts on the diagonal $\so{3,3}$ as the reflection of the direction $s$, an outer automorphism, which reverses $C_3$. (iv) $\tau_\sigma$
permutes the basis tensors of $S^{\otimes n}$ with the Koszul sign of $\omega$-parity, and $\omega$ has four even and four odd basis vectors. The image of
the symmetriser is spanned by orbit sums, and an orbit contributes iff no odd basis vector repeats, so its dimension is the number of
monomials of degree $n$ in four even and four odd variables, $\sum_k\binom{n-k+3}3\binom4k$; for the antisymmetriser exchange the roles of even
and odd, which gives the same number.
\end{proof}

Part (i) restates \cref{d5:prop:cmrel}(iii) at $n=2$ and is fixed at every rung (e3-5a at $n=3,4$; e3-5e, relations and vacuum at
$n\le5$). Parts (ii)--(iv) are fixed (e3-5b; e3-5c, exact at $n\le4$, the $C_3$ eigenvalues in floating point to $10^{-6}$; e3-5e). Since $\tau_k$
and the natural action $\Lambda(s_k)$ of the transposition on $\Lambda(\C^n\otimes\C^3)$ implement the same automorphism of $\Cl(3n,3n)$, they differ by
a scalar, and the vacuum fixes it: $\tau_\sigma=\operatorname{sgn}(\sigma)\Lambda(\sigma)$. The Koszul swaps are therefore the permutation matrices of
Howe's $O(n)$, twisted by the sign character. Part (iii) generalises the mirror statement of \cref{e1:prop:stat}(ii) to every rung: the
Koszul sign does not sort the module by Lorentz type or by copy type, and its two sectors carry the same content up to the reflection of
one direction.

\begin{proposition}[The ungraded $n$-particle space]
\label{e3:prop:ungraded}
Let $P_\sigma$ permute the tensor factors of $S_n$ without signs, and let $W_n=\{J^{(1)}=\cdots=J^{(n)}\}$, of dimension $2\cdot4^n$, with $J:=J^{(1)}|_{W_n}$ as
imaginary unit, so that $W_n\cong S\otimes_\C\cdots\otimes_\C S=(\C^4)^{\otimes_\C n}$ (\cref{e1:prop:select}(iv)).
\begin{enumerate}[label=(\roman*),leftmargin=2.2em]
\item $P_\sigma J^{(k)}P_\sigma^{-1}=J^{(\sigma k)}$, so $P_\sigma$ preserves $W_n$ and is $\C$-linear there.
\item The totally antisymmetric part of $W_n$ under $P_\sigma$ is $\Lambda^n_\C\C^4$, of real dimension $2\binom4n=8,12,8,2,0$ at $n=1,\dots,5$, and the
totally symmetric part is $\mathrm{Sym}^n_\C\C^4$, of real dimension $2\binom{n+3}3=8,20,40,70,112$.
\item Under the diagonal $\lint$, $\Lambda^2_\C\C^4$ has $C_\Delta=0$ on $4$ and $3$ on $8$ dimensions, $2(0,0)\oplus(\half,\half)$; $\Lambda^3_\C\C^4$ has
$C_\Delta=\tfrac32$ on all $8$, a Dirac spinor; and $\Lambda^4_\C\C^4$ has $C_\Delta=0$, a scalar.
\end{enumerate}
\end{proposition}

\begin{proof}
(i) $P_\sigma(A_1\otimes\cdots\otimes A_n)P_\sigma^{-1}$ permutes the factors. (ii) Exact traces of $\frac1{n!}\sum_\sigma\operatorname{sgn}(\sigma)P_\sigma$ and of $\frac1{n!}\sum_\sigma P_\sigma$, composed with the projector
$2^{1-n}\prod_k(\id-J^{(k)}J^{(k+1)})$ onto $W_n$, as signed permutations. (iii) The projector onto the antisymmetric part, as an exact integer matrix,
satisfies the stated polynomial identities with $C_\Delta$. Since $\C^4=(\half,0)\oplus(0,\half)$,
$\Lambda^2=\Lambda^2(\half,0)\oplus(\half,0)\otimes(0,\half)\oplus\Lambda^2(0,\half)$, $\Lambda^3\C^4\cong(\C^4)^*\otimes\Lambda^4\C^4$ and $\Lambda^4\C^4=\det$.
\end{proof}

The representation theory is inherited \cite{WeinbergQFT}; its realisation on $W_n$ is fixed (e3-5d). With the plain exchange and
antisymmetry imposed, rung $n$ on $W_n$ has the content of $n$ identical Dirac fermions at one momentum label, and exclusion holds: at most
four, since a Dirac spinor has four states, so there is no antisymmetric state at $n\ge5$ (exclusion inherited \cite{Pauli40}). The scope is
narrow. The vanishing of $\Lambda^n_\C\C^4$ is a statement about one mode of a finite construction; it is not a limit on the number of fermions in
physical space, which would need a continuum of momenta. The antisymmetry is imposed, not derived: neither exchange of the tower proves a
spin--statistics connection, which needs a local relativistic field theory with positive energy. The two exchanges therefore differ at
every rung, not only at $n=2$ (\cref{e1:prop:swap}). The Koszul exchange is the one the algebra supplies, since it permutes the vectors of
$\Cl(3n,3n)$. Its antisymmetric sector has dimension $88$ at $n=3$ and contains all three Lorentz types, and at $n=5$ it has dimension $360$
where the plain one is empty. It is not a statistics label. The plain exchange is a statistics label, but only on the constructed space $W_n$,
with the constructed generators $u_{(k,a)}$ as gamma matrices.

% --- source: emergence/e3_ladder#8 (The step from nn to n+1n+1, and what going up the ladder mea)
\subsection{The step from \texorpdfstring{$n$}{n} to \texorpdfstring{$n+1$}{n+1}, and what going up the ladder means}
\label{e3:sec:answer}

\paragraph{What the step preserves.} By \cref{e3:prop:rung}(ii) the step keeps rung $n$ as the subalgebra $\Cl(3n,3n)\otimes\id$, with all of its
operators. Under the spectator action the interior four-plane stays in the first factor, so the tower is one Lorentz frame with nested
internal centralisers (\cref{prop:frame}). The diagonal module is tensored with $S$ (\cref{e3:prop:rung}(iv)), and the dual pair $(O(n),\so{3,3})$ sits inside
$(O(n+1),\so{3,3})$.

\paragraph{What the step adds.} Three Fock modes, $N_{n+1}=N_n\otimes\id+\id\otimes N_1$ (\cref{e3:prop:rung}(iii)). One frame: $V_{n+1}=V_n\oplus V$, which
under the diagonal $\lG$ is one more field strength $(1,0)\oplus(0,1)$ and under the diagonal $\lint$ one more potential $(\half,\half)$ with two
scalars (\cref{sec:spintower}). One copy direction: the copy space $\R^n$ becomes $\R^{n+1}$, the copy algebra gains the $n$ rotations
$C^{k,n+1}$, and the Casimir constant grows by $\tfrac34(2n+5)$ (\cref{e3:thm:dual}(iii); e3-6c). The new constituent has its own commuting
copy, $(\Gamma_n\otimes\id)E_{(n+1,a)}$.

\paragraph{What the step changes.} The graded mass grows in quadrature and the ungraded energies add (\cref{e3:prop:mass}(iv)). The spin parity
under the diagonal $\lint$ flips, and that under the diagonal $\lG$ does not (\cref{e3:prop:spin}). Every parity-dependent selection of
\cref{e3:sec:table} switches on or off.

\paragraph{One particle or several.} Read together, the statements of this section answer the question in four sentences. (1) As given by
its vectors, every rung carries one Dirac mass $\sqrt n\,m$ along a unit vector of $\R^{3n,3n}$ with one clock, and copies that are the
components of a vector of the copy group $O(n)$ (fixed); read as particle content, this is one particle (interpretive). None of the tests
used here, mass shells, complex structures, exchanges or counting operators, is a particle count by itself: a particle interpretation needs a
state space with a positive inner product, a positive-energy representation of the Poincar\'e group and asymptotic states, and none is
constructed here. That group is the Howe dual of the diagonal $\so{3,3}$, its type is locked to
the $\so{3,3}$ type by $C_{\rm cop}+C_{\rm diag}=\tfrac34n(n+4)$, and its Lie algebra moves quanta between constituents without creating any. So
the copy index behaves as a flavour index of one particle, and going up the ladder adds a flavour direction and three modes to that particle
(fixed statements, read interpretively; \cref{e3:prop:mass,e3:thm:dual}). (2) The Lorentz content does not decide. The parity $(-1)^n$ under
the interior action is the parity of $n$ spin-$\half$ particles and equally that of one particle in an $n$-fold tensor product, and under
the field-strength action every rung is bosonic (fixed; inherited rule; \cref{e3:prop:spin}). (3) A rung carries $n$ particles in the sense of
ordinary quantum mechanics only with the structures of \cref{e1:sec:main}, now needed at every rung: the commuting copies
$u_{(k,a)}=\Gamma^{(1)}\cdots\Gamma^{(k-1)}E_{(k,a)}$ (fixed as an algebra), the additive Hamiltonian $H_u$ with its $n$ clocks, the common complex
space $W_n$ with $J|_{W_n}$, and the plain exchange with antisymmetry (constructed; the unit selected up to sign at odd $n$ only). With them
the rung has the content of $n$ identical Dirac fermions at one momentum label, with exclusion $n\le4$ at that label (fixed; inherited;
antisymmetry imposed; \cref{e3:prop:ungraded}). (4) The
exchange the algebra supplies is the Koszul $\mathfrak S_n$, the permutation matrices of $O(n)$ up to sign, and it is not a statistics label: its
two sectors are mirror images and never empty (fixed; \cref{e3:prop:koszul}).

\paragraph{Which process goes up the ladder.} Within a rung, the copy algebra hops quanta among the $n$ constituents and conserves $N$
(\cref{e3:thm:dual}(ii); \cref{e2:prop:hop}). Between rungs there is no process in the algebra: rung $n+1$ is not an element or a state of
rung $n$, and the embedding $\Cl(3n,3n)\otimes\id$ makes rung $n$ the part of rung $n+1$ that does not see the last constituent. In the
two-particle reading the step is the tensor embedding $\psi\mapsto\psi\otimes\phi_0$, which needs a state $\phi_0$ that nothing selects
(constructed, as at $n=2$). Whether a dynamical mechanism selects $W_n$, the additive Hamiltonian and a rung is open.

\begin{table}[!htbp]
\centering
\footnotesize
\begin{tabular}{@{}>{\raggedright\arraybackslash}p{8.2cm}>{\raggedright\arraybackslash}p{4.5cm}>{\raggedright\arraybackslash}p{2.2cm}@{}}
\toprule
statement & class & checks\\
\midrule
\cref{e3:prop:rung}(i): $\Cl(3n,3n)=M_{8^n}(\R)$, volume $\omega^{\otimes n}$, $(-1)^N=(-1)^n\Gamma_n$ & fixed in the realisation; classification inherited & e3-1a\\
\cref{e3:prop:rung}(ii), (iii): step formulas, eight copies, commutant $\id\otimes M_8(\R)$; Fock modes & restated (\cref{d5:prop:graded}, \cref{d5:prop:fock}; \cref{e1:prop:copies}) & e3-1b, e3-1c\\
\cref{e3:prop:rung}(iv): diagonal bivectors additive, $S_{n+1}=S_n\otimes S$ & fixed & e3-6c\\
\cref{e3:prop:mass}(i), (ii): $\pm\sqrt n\,m$; $m(n-2j)$ with $\binom nj4^n$; exact dressing & restated (\cref{e1:prop:rest}(vi), \cref{d6:prop:dirac}(iv)); multiplicities fixed & e3-2a, e3-2b\\
\cref{e3:prop:mass}(iii), (iv): one clock against $n$ clocks; the step & fixed & e3-2a, e3-2b, e3-6c\\
\cref{e3:prop:spin}(i), (iv): $e^{2\pi R_n}=(-1)^n$; $(-1)^{k_I}$; $C_\Delta$ spectra & restated (\cref{prop:orthospin}, \cref{prop:diagtower}(i)); spectra fixed & e3-3a, e3-3c\\
\cref{e3:prop:spin}(ii), (iii): $e^{2\pi L_n}=+\id$; $(\half,\half)^{\otimes n}$; central element & fixed (extends \cref{prop:halfspinors}) & e3-3b, e3-3c\\
\cref{e3:rem:copyparity}: spin parity against copy-invariant four-vectors & restated (\cref{thm:blockgamma}, \cref{thm:copyparity}); comparison interpretive & e3-6b\\
\cref{e3:thm:dual}(i): mutual commutants on vectors; $\eta$-weighting forced & fixed & e3-4a\\
\cref{e3:thm:dual}(ii): spin lifts, $\so n$, single-valued, Fock form & fixed; restated at $n=2,3$ (\cref{d5:prop:cmrel}, \cref{d7:prop:isospin}) & e3-4b, e3-4f\\
\cref{e3:thm:dual}(iii): $C_{\rm cop}+C_{\rm diag}=\tfrac34n(n+4)$; interior identity & inherited (Neerg{\aa}rd, eq.~(73)); independent proof; interior restated (\cref{thm:casid}) & e3-4c, z1-3a\\
\cref{e3:thm:dual}(iv): commutant at $n=2$ is the image of $\R[O(2)]$ & fixed & e3-4e\\
Howe's skew duality $(O(n),\so6)$ on $\Lambda(\C^n\otimes\C^3)$ & inherited \cite{Howe89} & e3-4f\\
\cref{e3:prop:joint}, \cref{e3:tab:joint}: joint decomposition at $n=2,3,4$ & fixed; $n=3$ copy spectrum restated (\cref{d7:prop:isospin}) & e3-4d\\
\cref{e3:prop:koszul}(i): Koszul $\mathfrak S_n$, Pin lift, commutation, vacuum sign & fixed; restated at $n=2$ (\cref{d5:prop:cmrel}(iii)) & e3-5a, e3-5e\\
\cref{e3:prop:koszul}(ii)--(iv): $\mathfrak S_3$ content; mirror $\Theta_s$; sector dimensions & fixed & e3-5b, e3-5c, e3-5e\\
\cref{e3:prop:ungraded}: $W_n$, $\Lambda^n_\C\C^4$, $\mathrm{Sym}^n_\C\C^4$, Lorentz content & fixed in the realisation; inherited (exterior powers, exclusion); one momentum label; antisymmetry imposed & e3-5d, z1-3c\\
\cref{e3:tab:ladder}: the ladder $n=1,\dots,5$ & fixed; restated (\cref{thm:centraliser}, \cref{thm:PS}, \cref{thm:algebra}, \cref{thm:copyparity}, \cref{prop:commutantsn}; \cref{d7:prop:isospin}) & e3-6a, e3-6b, as listed\\
\cref{e3:sec:table}: no rung selected by the algebra; odd, even, third rung & interpretive (drawn from fixed statements); period four inherited & ---\\
\cref{e3:sec:answer}: step; one particle or several; process & fixed parts as cited; reading interpretive; tensor embedding constructed; mechanism open & e3-1b, e3-1c, e3-6c\\
\bottomrule
\end{tabular}
\caption{Claims of \cref{e3:sec:main}, their classes and their checks, which are \texttt{PASS} lines of \texttt{e3\_ladder.py}.
``Restated'': established in an earlier chapter or in \cref{e1:sec:main}, re-verified here.}
\label{e3:tab:claims}
\end{table}

\section{Spin one up the tower}
\label{sec:w:spinoneupthetower}

% --- source: tex/s13_spin#6 (The tower)
\subsubsection{The tower}

\label{sec:spintower}
On the tower the two readings multiply in different ways.

\begin{proposition}[Commutants along the tower]
\label{prop:commutantsn}
On $V_n=\R^{3n,3n}$ the commutants have the following dimensions:
\begin{center}
\begin{tabular}{@{}lccccc@{}}
\toprule
commutant of & $n=1$ & $2$ & $3$ & $4$ & $5$\\
\midrule
spectator $\lint$ in $\so{3n,3n}$ & $1$ & $28$ & $91$ & $190$ & $325$\\
diagonal $\lint$ in $\so{3n,3n}$ & $1$ & $7$ & $18$ & $34$ & $55$\\
diagonal $\lG$ in $\so{3n,3n}$ & $0$ & $2$ & $6$ & $12$ & $20$\\
diagonal rotations and $G_n$ in $\so{6n}$ & $1$ & $4$ & $9$ & $16$ & $25$\\
the same inside $\so{3n,3n}$ & $0$ & $1$ & $3$ & $6$ & $10$\\
\bottomrule
\end{tabular}
\end{center}
These are the dimensions of $\so{3n-1,3n-3}$, $\so n\oplus\so{2n}$, $\so{n,\C}$, $\uu n$ and $\so n$.
\end{proposition}

Under the diagonal interior action $V_n=n\,(\tfrac12,\tfrac12)\oplus2n\,(0,0)$, so there are $n$ potentials forming the vector of the copy
algebra $\so n$, and none of them is copy-invariant at grade one for $n\ge2$. A single univectorial potential in the diagonal sector
therefore needs a choice of direction in copy space, and any such choice breaks $\so n$ to $\so{n-1}$. The obstruction belongs to grade
one only.

\begin{theorem}[Copy-invariant four-vectors]
\label{thm:copyparity}
A copy-invariant four-vector in $\Lambda^g(\R^n\otimes\R^{1,3})$ requires $n$ odd and $g$ odd with $n\le g\le3n$, and for every odd $n$
there is one at $g=n$,
\[
Q^{(n)}_b=\epsilon_{k_1\cdots k_n}\,\eta^{a_1a_2}\cdots\eta^{a_{n-2}a_{n-1}}\,\gamma^{(k_1)}_{a_1}\wedge\cdots\wedge\gamma^{(k_{n-1})}_{a_{n-1}}\wedge\gamma^{(k_n)}_b .
\]
At $n=3$ there is exactly one at each of the grades $3,5,7,9$, and at $n=2$ there is none at any grade.
\end{theorem}

\begin{proof}
An invariant element of $\Lambda^g(\R^n\otimes\R^{1,3})\otimes\R^{1,3}$ is an invariant tensor of $SO(n)\times SO^+(1,3)$ in
$(\R^n)^{\otimes g}\otimes(\R^{1,3})^{\otimes(g+1)}$. By the first fundamental theorem for the special orthogonal groups \cite{Weyl39},
such tensors are spanned by products of metrics and at most one Levi-Civita symbol. On the Lorentz side every factor has an even number
of indices, so $g+1$ is even. On the copy side $g=2d+ne$ with $e\in\{0,1\}$; $g$ odd forces $e=1$ and $n$ odd, so $g\ge n$. The Hodge
star of $\R^n\otimes\R^{1,3}$ is an equivariant isomorphism $\Lambda^g\cong\Lambda^{4n-g}$, which gives $g\le3n$. The displayed
$Q^{(n)}_b$ is invariant because $\epsilon$ is, and nonzero because the terms of $Q^{(n)}_0$ with all Lorentz indices equal to $0$
add up to $n!\,\gamma^{(1)}_0\wedge\cdots\wedge\gamma^{(n)}_0$.
\end{proof}

At $n=3$ the invariant four-vector of grade three is, up to a factor, the copy-singlet cubic $Q_\mu$ that builds the gamma matrices of
the Dirac block in \cref{thm:blockgamma}; the parity theorem explains why such an object exists at three constituents and could not
exist at two. Under the diagonal $\lG$, finally, $V_n=n\bigl[(1,0)\oplus(0,1)\bigr]$: $n$ field strengths mixed by the complex orthogonal
flavour algebra $\so{n,\C}$. The energy-preserving duality group of $n$ free Maxwell fields is $U(n)$, the maximal compact subgroup of
the Gaillard--Zumino group $Sp(2n,\R)$ \cite{GaillardZumino81}; its Lie algebra meets $\so{3n,3n}$ only in $\so n$, and its centre,
generated by $G_n$ and carrying the total helicity, lies outside $\so{3n,3n}$, as $G$ does at $n=1$. A single photon would be one
direction in this flavour space, and nothing selects one.

% --- source: tex/s13_spin#7 (The mixed rung)
\subsubsection{The mixed rung}
\label{sec:mixed}

\Cref{prop:emexclusive} forbids reading one six-dimensional space in two ways. At $n\ge2$ there is more than one constituent, and the
obvious repair is to let one constituent carry the interior reading and another the frame reading, under one Lorentz action.

\begin{proposition}[The mixed embedding]
\label{prop:mixed}
In $\Cl(6,6)$ let $\mathfrak l_{\rm mix}$ be spanned by $S_k\otimes1+1\otimes L_k$ and $T_k\otimes1+1\otimes K_k$: the interior algebra on
constituent $1$ and $\lG$ on constituent $2$, paired generator by generator. It satisfies the $\so{1,3}$ brackets, and
\[
\R^{6,6}=\underbrace{V}_{C=3}\oplus\underbrace{\Lambda^2V}_{C=4}\oplus2\cdot\mathbf1 ,
\]
so one Lorentz action carries both a four-vector and a field strength at grade one. On $S_2=\R^{64}$ the Casimir satisfies
$(C_{\rm mix}-\tfrac32)(C_{\rm mix}-\tfrac{11}2)=0$, with the Dirac value $\tfrac32$ on sixteen dimensions and $\tfrac{11}2$ on $48$.
\end{proposition}

\begin{proposition}[The cost]
\label{prop:mixedcost}
The centraliser of $\mathfrak l_{\rm mix}$ in $\so{6,6}$ is one-dimensional, spanned by the rotation of the surplus pair of
constituent $1$. At the same rung:
\begin{center}
\begin{tabular}{@{}lcl@{}}
\toprule
Lorentz action at $n=2$ & centraliser & content of $\R^{6,6}$\\
\midrule
interior on one constituent & $28$ & $V\oplus8\cdot\mathbf1$\\
$\lG$ on one constituent & $15$ & $\Lambda^2V\oplus6\cdot\mathbf1$\\
mixed, across both & $1$ & $V\oplus\Lambda^2V\oplus2\cdot\mathbf1$\\
\bottomrule
\end{tabular}
\end{center}
The mixed embedding buys the field-strength module at the price of the whole internal algebra: $\so{5,3}$ collapses to one $\uu1$.
\end{proposition}

The count is for this pairing of the two Lorentz bases at $n=2$; another pairing or a higher rung could do better and has not been
examined. Within its scope, however, the mixed rung is an honest place to test whether the two readings can be coupled, and it turns out
that they can.

\begin{theorem}[The Dirac block of the mixed rung]
\label{thm:mixeddirac}
Let $P=\tfrac18(11\,\id-2C_{\rm mix})$ and $\hat\gamma_\mu=2P\gamma^{(1)}_\mu P$, with $J_1=e_2e_3\otimes1$.
\begin{enumerate}[label=(\roman*),leftmargin=2.2em]
\item $P$ is the orthogonal projector of rank $16$ onto the block with the Dirac value, and commutes with the mixed generators.
\item $\{\hat\gamma_\mu,\hat\gamma_\nu\}=2\eta_{\mu\nu}P$, $[\hat\gamma_\mu,J_1]=0$, $J_1^2=-\id$, and $P\mathcal S_kP=\tfrac12\hat\gamma_i\hat\gamma_j$ (cyclic),
$P\mathcal T_kP=\tfrac12\hat\gamma_0\hat\gamma_k$: the compressed gammas generate the mixed Lorentz action on the block.
\item With $J_1$ as the imaginary unit the block is two complex Dirac modules, labelled by the half-spinors of constituent $2$, with
flavour algebra $\uu2$; $\mathcal H(p)=\hat\gamma^0\hat\gamma^ap_a+m\hat\gamma^0$ is symmetric and $\mathcal H(p)^2=(|p|^2+m^2)P$.
\item The block is the $\gamma$-trace part $\sum_a\gamma_a\chi\otimes\tau e_a$ of the vector-spinors (Dirac spinor of constituent $1$)
$\otimes$(four-vector of constituent $2$), and the rank-$48$ block is their $\gamma$-traceless $(1,\tfrac12)\oplus(\tfrac12,1)$ part.
\end{enumerate}
\end{theorem}

Item (iv) explains the factor $2$ in $\hat\gamma_\mu$, since $P\gamma^{(1)}_\mu$ acts on the $\gamma$-trace part as $-\tfrac12\gamma_\mu$, and it
identifies the decomposition with the classical splitting of a vector-spinor into its trace and its traceless part
\cite{RaritaSchwinger41}.

A covariant identification of the field strength of constituent $1$ with the frame field of constituent $2$ is an invariant element
of $\Lambda^2(\text{four-plane}_1)\otimes V_2$, that is, a grade-three element of the algebra commuting with the mixed generators. The
invariant space is two-dimensional, and one invariant is
\[
\Theta_F=\gamma^{(1)}_2\gamma^{(1)}_3e'_1+\gamma^{(1)}_3\gamma^{(1)}_1e'_2+\gamma^{(1)}_1\gamma^{(1)}_2e'_3
-\gamma^{(1)}_0\gamma^{(1)}_1e'_4-\gamma^{(1)}_0\gamma^{(1)}_2e'_5-\gamma^{(1)}_0\gamma^{(1)}_3e'_6 ,
\]
with $e'_a$ the generators of constituent $2$; it pairs the spatial bivectors of constituent $1$ with the magnetic directions of
constituent $2$ and the time--space bivectors with the electric ones. The other invariant is $\Theta_F$ followed by $G$, so up to scale
the identifications form the family $X=e^{\theta G}X_F$.

\begin{theorem}[A coupled model]
\label{thm:coupled}
On the block, in flavour coordinates $\R^8\otimes\R^2$ with $\hat\gamma_\mu=-\gamma_\mu\otimes1$, $J_1=J\otimes1$ and $U^\mu=\hat\gamma^0\hat\gamma^\mu$, take
for commuting real fields
\begin{equation}
\label{eq:coupledL}
\mathcal L=-\tfrac14F_{\mu\nu}F^{\mu\nu}+\psi^{\!\top}J_1U^\mu D_\mu\psi-m\,\psi^{\!\top}\hat\gamma^0\psi,\qquad D_\mu=\partial_\mu+qA_\mu J_1 ,
\end{equation}
the real form of $-\tfrac14F^2+\bar\psi(i\gamma^\mu D_\mu-m)\psi$ with $J_1$ in the place of $i$. Then the field equations are Maxwell's
equations $\partial_\mu F^{\mu\nu}=q\,j^\nu$ with $j^\nu=\psi^{\!\top}U^\nu\psi$ and $j^0=|\psi|^2$, and the Dirac equation
$J_1\hat\gamma^\mu D_\mu\psi=m\psi$; the identity $\partial_\mu j^\mu+\psi^{\!\top}J_1\mathcal E_\psi=0$ holds with $\mathcal E_\psi$ the Euler--Lagrange
derivative of the action, $\mathcal E_\psi=2\hat\gamma^0(J_1\hat\gamma^\mu D_\mu\psi-m\psi)$, and not with the Dirac expression itself; the Lagrangian is exactly invariant under $\psi\mapsto e^{-q\lambda J_1}\psi$, $A\mapsto A+\partial\lambda$; and its flavour
symmetry is the $\uu2$ spanned by $J\otimes1$, $1\otimes\epsilon$, $J\otimes\sigma_x$, $J\otimes\sigma_z$. The potential carries two
polarisations of helicity $\pm1$ and no massive mode, and each momentum carries two flavours, two helicities and two signs of frequency.
\end{theorem}

In the frame reading the source is electric. For the frame field of an arbitrary potential one has, identically,
$\partial_tz+i\curl z+i\boldsymbol{\mathcal J}=0$ and $\nabla\cdot z=i\rho$ with $(\rho,\boldsymbol{\mathcal J})=\partial_\mu F^{\mu\nu}$, and the
rotated field $e^{\theta G}X_F$ satisfies the same system with electric source $\cos\theta$ times the current and magnetic source $\sin\theta$
times it. Since the frame field takes no part in the dynamics of \eqref{eq:coupledL}, which member of the family is called the frame
field is a choice; requiring it to carry electric sources only fixes $\theta\equiv0$ modulo $\pi$, as in the textbook treatment of duality
with sources, where the angle is fixed by the convention that all charges are electric \cite{Jackson}. The duality rotation, a symmetry
of the free system and of \eqref{eq:SSL}, is not a symmetry once electric sources are present.

The coupled model is classical two-flavour spinor electrodynamics with equal charges, written in the algebra of the second rung. It
is a construction, not a selection, and it has a definite limitation that we record now and use later: the second-rung vacuum of
\cref{sec:photon} is defined for the spectator action, and under the mixed action it is preserved only by the rotations about one axis,
so the coupled model lies outside that vacuum sector.

At even rungs, finally, the diagonal interior action places integer-spin modules inside the spinor module itself: at $n=2$ the Casimir
takes the values $0,3,4$ on $8,32,24$ dimensions, so $S_2$ contains scalars of copy charge $\pm2$, four-vectors of copy charge $\pm1$ and
copy-neutral field strengths $(1,0)\oplus(0,1)$, and the copy and Lorentz Casimirs add up to $4$, the $n=2$ case of \eqref{eq:casid}. These
are modules of one spinor space, not bound states, and we do not propose them as the photon.

% --- source: emergence/e8_synthesis#0 ((head))
% Source for the tower chapter of the global manuscript; this note is not compiled separately.
% Evidence identifiers: e*-n, y*-n and z1-n are PASS lines of verify/e*.py, verify/y*_check.py and z1_phase1.py.

% --- source: emergence/e8_synthesis#1 (Synthesis: what the tower establishes about particle multipl)
\section{Synthesis: what the tower establishes about particle multiplicity}
\label{e8:sec:main}

The question of this section was whether, and in what sense, a rung $\Cl(3n,3n)$ of the tower carries several particles rather than one
particle with a larger internal space. Particle multiplicity was tested in five senses, each with its own operational meaning: separate
mass shells, a common complex structure for the copies, a counting operator, an exchange that sorts by statistics, and a process that takes
one to two. This section states the outcome claim by claim, in four classes: \emph{proved algebra} (theorems, exact computations and
restatements of fixed results), \emph{conditional or constructed structures}, \emph{obstructions within a stated class}, and
\emph{missing dynamics}. \cref{e8:tab:claims} gives the evidence for each claim.

% --- source: emergence/e8_synthesis#2 (The answer)
\subsection{The answer}
\label{e8:sec:answer}

\emph{The tower supplies the algebraic room for $n$-fold multiplicity, but it does not establish particle multiplicity.}
\begin{enumerate}[label=(\arabic*),leftmargin=2.2em]
\item \emph{What is proved.} Every rung is $M_{8^n}(\R)$, and the step embeds rung $n$ in rung $n+1$ with multiplicity eight. A second, commuting copy
of each constituent exists and is unique up to conjugation; its generators are the graded ones dressed by the volume elements of the
earlier constituents. The graded masses add in quadrature to one Dirac mass $(\sum_km_k^2)^{1/2}$, with no simultaneous nonzero-mass Dirac constraints of the graded form, while the
commuting copies have additive spectra $\{\sum_k\pm m_k\}$. A common complex structure for all copies exists only on the subspace
$W_n=\{J^{(1)}=\cdots=J^{(n)}\}$, of real dimension $2\cdot4^n$. The diagonal interior action selects an imaginary unit exactly at odd $n$, and none
at even $n$. The Fock number $N$ is conserved exactly by $\mathfrak{gl}(3n,\R)$. Within the maximal compact subalgebra, it is conserved exactly
by the paired rotations, the compact binding-compatible fields. No element of $\so{3n,3n}$ changes it by an odd number. $N$ is invariant
under no Lorentz action of the tower, and a Lorentz-invariant number of quanta exists, under the common binding action, exactly at even
$n$. On vectors, the copy algebra with the copy permutations adjoined and the diagonal $\so{3,3}$ are mutual centralisers
(\cref{e3:thm:dual}(i)); on $S_n$ the Casimirs satisfy $C_{\rm cop}+C_{\rm diag}=\tfrac34n(n+4)$, and the dual-pair statement there concerns the full
$O(n)$ action, reflections included \cite{Neergard20}. The compact field
between constituents is a one-body operator: it moves quanta and never creates one.
\item \emph{What is constructed, not selected.} A reading with $n$ particles in the sense of ordinary quantum mechanics needs three
structures together: the subspace $W_n$ with the complex structure $J|_{W_n}$, the additive Hamiltonian $H_u$ built from the commuting copies,
and the plain exchange $P$ with antisymmetry imposed. Each exists. None is selected by the algebra or by a Lorentz action of the tower.
With them, rung $n$ has the content of $n$ identical Dirac fermions at one momentum label, with exclusion $n\le4$ at that label.
\item \emph{What is obstructed.} At $n=2$ no selection principle of the paper gives a complex structure under which both copies are complex
Dirac spinors. Under the diagonal action none exists, and under the spectator action copy~2 carries no linear Dirac operator. No
gauge-admissible, binding-compatible field changes $N$. Within the binding-compatible algebra the number-changing directions are the
paired boosts, which are not gauge fields; outside it, compact generators and non-compact ones, the interior boosts of \cref{e2:prop:bogo}
among them, also mix sectors differing by two quanta. The Clifford exchange $\tau_K$ is not a statistics label. At $n=2$ none of the compact parts of
the five candidate gauge algebras contains the Pati--Salam algebra.
\item \emph{What is missing.} No state space with a positive inner product carrying a positive-energy representation of the Poincar\'e
group is constructed. There are no asymptotic states, no derived quantum theory of tower-changing processes (\cref{thm:coupled} is a classical
interacting model), no dynamics that selects $W_n$, $H_u$, an exchange or a rung,
and no proof of a spin--statistics connection. Without these, ``one particle with a larger internal space'' and ``$n$ particles'' are two
readings of the same algebra, and the tower does not decide between them.
\end{enumerate}

Three readings that earlier drafts suggested are not supported. (a) The gap $m_1+m_2-(m_1^2+m_2^2)^{1/2}$ is not a binding energy. It compares two
operators on one space at equal momentum labels, and at fixed total momentum the graded levels with large relative momentum lie above the
free-pair threshold (\cref{e2:prop:threshold}; z1-2a). (b) Fock mixing by a boost, a non-unitary operator on a finite space at one momentum
label, is not observable pair production. (c) The time-odd triangle signs of the copy field are holonomies in copy space, not physical CP
violation.

